\input{preamble}

% OK, start here.
%
\begin{document}

\title{Les axiomes d'Artin}

\maketitle

\phantomsection
\label{section-phantom}

\tableofcontents




\section{Introduction}
\label{section-introduction}

\noindent
Dans ce chapitre, nous étudions les axiomes d'Artin assurant la
représentabilité des foncteurs par des espaces algébriques. Nous renvoyons aux
articles
\cite{ArtinI}, \cite{ArtinII}, \cite{ArtinVersal}.

\medskip\noindent
Certaines notations, conventions et expressions employées dans ce chapitre
sont peu naturelles et pourront sembler prises à rebours au lecteur averti.
C'est intentionnel. On en trouvera l'explication dans
Quotients, section \ref{quot-section-conventions}.

\medskip\noindent
Soit $S$ un schéma de base localement noethérien. Soit
$$
p : \mathcal{X} \longrightarrow (\Sch/S)_{fppf}
$$
une catégorie fibrée en groupoïdes. Soit $x_0$ un objet de $\mathcal{X}$
au-dessus d'un corps $k$ de type fini sur $S$. La catégorie de
prédéformations
(voir Théorie formelle des déformations,
définition \ref{formal-defos-definition-predeformation-category})
$$
\mathcal{F}_{\mathcal{X}, k, x_0}
\longrightarrow
\{\text{$S$-algèbres locales artiniennes de corps résiduel }k\}
$$
associée à $x_0$ sur $k$ joue un rôle important tout au long de ce chapitre.
Nous introduisons pour $\mathcal{X}$ la condition de Rim--Schlessinger (RS) et
montrons qu'elle assure que $\mathcal{F}_{\mathcal{X}, k, x_0}$ est une
catégorie de déformations, autrement dit que
$\mathcal{F}_{\mathcal{X}, k, x_0}$ satisfait elle-même (RS). Nous étudions la
façon dont $\mathcal{F}_{\mathcal{X}, k, x_0}$ varie lorsqu'on remplace $k$ par
une extension finie, ainsi que les espaces tangents.

\medskip\noindent
Nous étudions ensuite les objets formels $\xi = (\xi_n)$ de $\mathcal{X}$,
c'est-à-dire les systèmes projectifs d'objets situés au-dessus des quotients
$R/\mathfrak m^n$, où $R$ est une $S$-algèbre locale noethérienne complète
dont le corps résiduel est de type fini sur $S$. Cela revient à se donner un
objet formel de $\mathcal{F}_{\mathcal{X}, k, x_0}$ pour certains $x_0$ et $k$.
Un objet formel est dit effectif lorsqu'il provient d'un objet de
$\mathcal{X}$ au-dessus de $R$. Un objet formel de $\mathcal{X}$ est dit versel
lorsqu'il donne lieu à un objet formel versel de
$\mathcal{F}_{\mathcal{X}, k, x_0}$. Enfin, étant donnés un $S$-schéma $U$ de
type fini, un objet $x$ de $\mathcal{X}$ au-dessus de $U$ et un point fermé
$u_0 \in U$, nous disons que $x$ est versel en $u_0$ si l'objet formel induit
sur l'anneau local complété $\mathcal{O}_{U, u_0}^\wedge$ est versel.

\medskip\noindent
Après ces préliminaires, nous pouvons énoncer le célèbre théorème d'Artin :
notre $\mathcal{X}$ est un champ algébrique si les conditions suivantes sont
satisfaites :
\begin{enumerate}
\item $\mathcal{O}_{S, s}$ est un anneau G pour tout $s \in S$\footnote{Cela
signifie que $S$ est un schéma G ; voir
Propriétés, section \ref{properties-section-G-scheme}.},
\item $\Delta : \mathcal{X} \to \mathcal{X} \times \mathcal{X}$
est représentable par des espaces algébriques,
\item $\mathcal{X}$ est un champ pour la topologie étale,
\item $\mathcal{X}$ commute aux limites,
\item $\mathcal{X}$ satisfait (RS),
\item les espaces tangents et les espaces d'automorphismes infinitésimaux
des catégories de déformations $\mathcal{F}_{\mathcal{X}, k, x_0}$
sont de dimension finie,
\item les objets formels sont effectifs,
\item $\mathcal{X}$ satisfait l'ouverture de la versalité.
\end{enumerate}
C'est le lemme \ref{lemma-diagonal-representable} ; voir aussi la
proposition \ref{proposition-second-diagonal-representable}, qui en donne une
légère amélioration. Une proposition analogue caractérise les foncteurs
$F : (\Sch/S)_{fppf}^{opp} \to \textit{Sets}$ qui sont des espaces
algébriques ; voir la section \ref{section-algebraic-spaces}.

\medskip\noindent
Voici les grandes lignes de la démonstration du théorème d'Artin. Nous montrons
d'abord, à l'aide de (RS) et de la dimension finie des espaces tangents et
d'automorphismes, qu'il existe suffisamment d'objets formels versels ; voir par
exemple Théorie formelle des déformations, lemme
\ref{formal-defos-lemma-minimal-groupoid-in-functors-construction}. Par
hypothèse, ces objets formels sont effectifs. Les objets formels effectifs
peuvent être « approchés » par des objets $x$ au-dessus de $S$-schémas $U$ de
type fini ; voir le lemme \ref{lemma-approximate}. Cette approximation utilise
le fait que les anneaux locaux de $S$ sont des anneaux G et que $\mathcal{X}$
commute aux limites. C'est peut-être la partie la plus difficile de la
démonstration : elle repose sur la désingularisation générale de Néron pour
approcher des solutions formelles d'équations algébriques sur un anneau G
local noethérien par des solutions dans son hensélisé. L'ouverture de la
versalité permet ensuite, quitte à rétrécir $U$, de supposer que $x$ est versel
en tout point fermé de $U$. Nous montrons alors que $U \to \mathcal{X}$ est un
morphisme lisse. En prenant suffisamment de morphismes
$U \to \mathcal{X}$, nous obtenons un « atlas lisse » de $\mathcal{X}$, ce qui
montre que $\mathcal{X}$ est un champ algébrique.

\medskip\noindent
Lorsque l'on vérifie les axiomes d'Artin pour une catégorie $\mathcal{X}$
fibrée en groupoïdes, l'étape la plus difficile consiste souvent à établir
l'ouverture de la versalité. Dans la suite de cette discussion, supposons que
$\mathcal{X}/S$ satisfasse déjà les autres conditions ci-dessus. Nous
présentons deux méthodes permettant de montrer que $\mathcal{X}$ satisfait
l'ouverture de la versalité :
\begin{enumerate}
\item La première consiste à supposer une condition de Rim--Schlessinger plus
forte, notée (RS*), ainsi qu'une forme renforcée de l'effectivité formelle : on
demande essentiellement que les objets au-dessus de systèmes projectifs
d'épaississements soient effectifs. Sous ces hypothèses, l'ouverture de la
versalité est automatique ; voir le lemme
\ref{lemma-SGE-implies-openness-versality}. Observons que nous utilisons ici de
façon essentielle que $\mathcal{X}$ est définie sur la catégorie de tous les
schémas au-dessus de $S$, et non seulement sur celle des schémas noethériens.
\item La seconde, suivant Artin, consiste à munir $\mathcal{X}$ d'une théorie
des obstructions. Si cette théorie des obstructions « commute aux produits »
en un sens convenable, alors $\mathcal{X}$ satisfait l'ouverture de la
versalité ; voir le lemme \ref{lemma-get-openness-obstruction-theory}.
\end{enumerate}
Les théories des obstructions peuvent être axiomatisées de nombreuses manières,
et la littérature en contient effectivement beaucoup de variantes, souvent
adaptées à des champs de modules particuliers. Dans le lemme
\ref{lemma-dual-openness}, nous en présentons une variante fondée sur la
catégorie dérivée, qui apparaît souvent naturellement dans les calculs de
théorie des déformations.

\medskip\noindent
Dans la section \ref{section-algebraic-spaces-noetherian}, nous expliquons les
modifications nécessaires pour les foncteurs définis sur la catégorie
$(\textit{Noetherian}/S)_\etale$ des schémas localement noethériens au-dessus
de $S$.

\medskip\noindent
Dans la dernière section du chapitre, nous démontrons, comme application des
axiomes d'Artin, le théorème d'Artin sur l'existence des contractions ; voir la
section \ref{section-contractions}. En gros, ce théorème affirme que, si l'on se
donne un espace algébrique $X'$ séparé et de type fini sur $S$, un fermé
$T' \subset |X'|$ et une modification formelle
$$
\mathfrak{f} : X'_{/T'} \longrightarrow \mathfrak{X}
$$
où $\mathfrak{X}$ est un espace algébrique formel noethérien au-dessus de $S$,
il existe un morphisme propre $f : X' \to X$ qui « réalise la contraction ».
Cela signifie qu'il existe une identification $\mathfrak{X} = X_{/T}$ telle
que $\mathfrak{f} = f_{/T'} : X'_{/T'} \to X_{/T}$, où $T = f(T')$, et que de
plus $f$ est un isomorphisme au-dessus de $X \setminus T$. La démonstration
consiste à définir un foncteur $F$ sur la catégorie des schémas localement
noethériens au-dessus de $S$, puis à vérifier pour $F$ les axiomes d'Artin.
Curieusement, dans cette application, l'ouverture de la versalité n'est pas le
point le plus difficile : c'est la preuve que $F$ commute aux limites qui exige
de nombreux travaux et résultats préliminaires.




\section{Conventions}
\label{section-conventions}

\noindent
Les conventions employées dans ce chapitre sont les mêmes que dans le chapitre
sur les champs algébriques ; voir Champs algébriques, section
\ref{algebraic-section-conventions}. Dans ce chapitre, le schéma de base $S$
sera souvent localement noethérien, bien que nous répétions toujours cette
hypothèse dans l'énoncé des résultats.





\section{Catégories de prédéformations}
\label{section-predeformation-categories}

\noindent
Soit $S$ un schéma de base localement noethérien. Soit
$$
p : \mathcal{X} \longrightarrow (\Sch/S)_{fppf}
$$
une catégorie fibrée en groupoïdes. Soit $k$ un corps, et soit
$\Spec(k) \to S$ un morphisme de type fini (voir Morphismes, lemme
\ref{morphisms-lemma-point-finite-type}). Nous dirons parfois simplement que
{\it $k$ est un corps de type fini sur $S$}. Soit $x_0$ un objet de
$\mathcal{X}$ situé au-dessus de $\Spec(k)$. À partir de $S$, $\mathcal{X}$,
$k$ et $x_0$, nous allons construire une catégorie de prédéformations au sens
de Théorie formelle des déformations, définition
\ref{formal-defos-definition-predeformation-category}. Cette construction est
analogue à celle de Théorie formelle des déformations, remarque
\ref{formal-defos-remark-localize-cofibered-groupoid}.

\medskip\noindent
D'après Morphismes, lemme \ref{morphisms-lemma-point-finite-type}, nous pouvons
d'abord choisir un ouvert affine $\Spec(\Lambda) \subset S$ tel que
$\Spec(k) \to S$ se factorise par $\Spec(\Lambda)$ et que l'homomorphisme
d'anneaux associé $\Lambda \to k$ soit fini. On obtient ainsi la catégorie
$\mathcal{C}_\Lambda$ ; voir Théorie formelle des déformations, définition
\ref{formal-defos-definition-CLambda}. À équivalence canonique près, la
catégorie $\mathcal{C}_\Lambda$ ne dépend pas du choix de l'ouvert affine
$\Spec(\Lambda)$ de $S$. En effet, $\mathcal{C}_\Lambda$ est équivalente à
l'opposée de la catégorie des factorisations
\begin{equation}
\label{equation-factor}
\Spec(k) \to \Spec(A) \to S
\end{equation}
du morphisme structural pour lesquelles $A$ est un anneau local artinien et
$\Spec(k) \to \Spec(A)$ correspond à un homomorphisme d'anneaux $A \to k$
identifiant $k$ au corps résiduel de $A$.

\medskip\noindent
Notons $\mathcal{F} = \mathcal{F}_{\mathcal{X}, k, x_0}$ la catégorie dont
\begin{enumerate}
\item les objets sont les morphismes $x_0 \to x$ de $\mathcal{X}$ tels que
$p(x) = \Spec(A)$, avec $A$ anneau local artinien, et que
$p(x_0) \to p(x) \to S$ soit une factorisation comme en
(\ref{equation-factor}) ;
\item les morphismes $(x_0 \to x) \to (x_0 \to x')$ sont les diagrammes
commutatifs
$$
\xymatrix{
x & & x' \ar[ll] \\
& x_0 \ar[lu] \ar[ru]
}
$$
dans $\mathcal{X}$. (Remarquer le renversement des flèches.)
\end{enumerate}
Si $x_0 \to x$ est un objet de $\mathcal{F}$, l'écriture
$p(x) = \Spec(A)$ fournit un objet $A$ de $\mathcal{C}_\Lambda$. Nous dirons
souvent que $x_0 \to x$, ou que $x$, est au-dessus de $A$. Un morphisme de
$\mathcal{F}$ entre des objets $x_0 \to x$ au-dessus de $A$ et
$x_0 \to x'$ au-dessus de $A'$ correspond à un morphisme $x' \to x$ de
$\mathcal{X}$, donc à un morphisme
$p(x' \to x) : \Spec(A') \to \Spec(A)$, lequel correspond à son tour à un
homomorphisme d'anneaux $A \to A'$. Comme $\mathcal{X}$ est une catégorie
au-dessus de la catégorie des schémas sur $S$, l'homomorphisme $A \to A'$ est
un homomorphisme de $\Lambda$-algèbres. Nous obtenons donc un foncteur
\begin{equation}
\label{equation-predeformation-category}
p : \mathcal{F} = \mathcal{F}_{\mathcal{X}, k, x_0}
\longrightarrow
\mathcal{C}_\Lambda.
\end{equation}
Nous noterons $\mathcal{F}(A)$ la catégorie-fibre au-dessus d'un objet $A$ de
$\mathcal{C}_\Lambda$. Un objet de $\mathcal{F}(A)$ est simplement un
morphisme $x_0 \to x$ de $\mathcal{X}$ tel que $x$ soit au-dessus de
$\Spec(A)$ et que $x_0 \to x$ soit au-dessus de
$\Spec(k) \to \Spec(A)$.

\begin{lemma}
\label{lemma-predeformation-category}
Le foncteur $p : \mathcal{F} \to \mathcal{C}_\Lambda$ défini ci-dessus est
une catégorie de prédéformations.
\end{lemma}

\begin{proof}
Il faut montrer (a) que $\mathcal{F}$ est cofibrée en groupoïdes sur
$\mathcal{C}_\Lambda$ et (b) que $\mathcal{F}(k)$ est équivalente à une
catégorie ayant un seul objet et un seul morphisme.

\medskip\noindent
Démonstration de (a). Les catégories-fibres de $\mathcal{F}$ au-dessus de
$\mathcal{C}_\Lambda$ sont des groupoïdes, puisque celles de $\mathcal{X}$ le
sont. Soit $A \to A'$ un morphisme de $\mathcal{C}_\Lambda$, et soit
$x_0 \to x$ un objet de $\mathcal{F}(A)$. Comme $\mathcal{X}$ est fibrée en
groupoïdes, il existe un morphisme $x' \to x$ au-dessus de
$\Spec(A') \to \Spec(A)$. La composée $A \to A' \to k$ étant l'homomorphisme
donné $A \to k$, l'unicité des images réciproques à isomorphisme près montre
que l'image réciproque de $x'$ par $\Spec(k) \to \Spec(A')$ est $x_0$ ;
autrement dit, il existe un morphisme $x_0 \to x'$ au-dessus de
$\Spec(k) \to \Spec(A')$, compatible avec $x_0 \to x$ et $x' \to x$.
Ainsi, $\mathcal{F}$ admet des images directes. On conclut par le dual de
Catégories, lemme \ref{categories-lemma-fibred-groupoids}.

\medskip\noindent
Démonstration de (b). Si $A = k$, alors $\Spec(k) = \Spec(A)$. Puisque
$\mathcal{X}$ est fibrée en groupoïdes sur $(\Sch/S)_{fppf}$, pour tout objet
$x_0 \to x$ de $\mathcal{F}(k)$, le morphisme $x_0 \to x$ est un
isomorphisme. Tout objet de $\mathcal{F}(k)$ est donc isomorphe à
$x_0 \to x_0$, dont le seul endomorphisme dans $\mathcal{F}$ est évidemment
l'identité.
\end{proof}

\noindent
Soit $S$ un schéma de base localement noethérien. Soit
$F : \mathcal{X} \to \mathcal{Y}$ un $1$-morphisme entre catégories fibrées
en groupoïdes sur $(\Sch/S)_{fppf}$. Soit $k$ un corps de type fini sur $S$.
Soit $x_0$ un objet de $\mathcal{X}$ au-dessus de $\Spec(k)$. Posons
$y_0 = F(x_0)$ ; c'est un objet de $\mathcal{Y}$ au-dessus de $\Spec(k)$.
Alors $F$ induit un foncteur
\begin{equation}
\label{equation-functoriality}
F :
\mathcal{F}_{\mathcal{X}, k, x_0}
\longrightarrow
\mathcal{F}_{\mathcal{Y}, k, y_0}
\end{equation}
de catégories cofibrées sur $\mathcal{C}_\Lambda$. Plus précisément, à l'objet
$x_0 \to x$ de $\mathcal{F}_{\mathcal{X}, k, x_0}(A)$ nous associons l'objet
$F(x_0) \to F(x)$ de $\mathcal{F}_{\mathcal{Y}, k, y_0}(A)$.

\begin{lemma}
\label{lemma-formally-smooth-on-deformation-categories}
Soit $S$ un schéma localement noethérien. Soit
$F : \mathcal{X} \to \mathcal{Y}$ un $1$-morphisme de catégories fibrées en
groupoïdes sur $(\Sch/S)_{fppf}$. Supposons que l'une des conditions suivantes
soit satisfaite :
\begin{enumerate}
\item $F$ est formellement lisse sur les objets (Critères de représentabilité,
section \ref{criteria-section-formally-smooth}) ;
\item $F$ est représentable par des espaces algébriques et formellement lisse ;
\item $F$ est représentable par des espaces algébriques et lisse.
\end{enumerate}
Alors, pour tout corps $k$ de type fini sur $S$ et tout objet $x_0$ de
$\mathcal{X}$ au-dessus de $k$, le foncteur
(\ref{equation-functoriality}) est lisse au sens de Théorie formelle des
déformations, définition
\ref{formal-defos-definition-smooth-morphism}.
\end{lemma}

\begin{proof}
Le cas (1) résulte directement des définitions. L'hypothèse (2) implique (1)
d'après Critères de représentabilité, lemme
\ref{criteria-lemma-representable-by-spaces-formally-smooth}. L'hypothèse (3)
implique (2) d'après Compléments sur les morphismes d'espaces, lemme
\ref{spaces-more-morphisms-lemma-smooth-formally-smooth}
et le principe de Champs algébriques, lemme
\ref{algebraic-lemma-representable-transformations-property-implication}.
\end{proof}

\begin{lemma}
\label{lemma-fibre-product-deformation-categories}
Soit $S$ un schéma localement noethérien. Soit
$$
\xymatrix{
\mathcal{W} \ar[d] \ar[r] & \mathcal{Z} \ar[d] \\
\mathcal{X} \ar[r] & \mathcal{Y}
}
$$
un $2$-produit fibré de catégories fibrées en groupoïdes sur
$(\Sch/S)_{fppf}$. Soit $k$ un corps de type fini sur $S$, et soit $w_0$ un
objet de $\mathcal{W}$ au-dessus de $k$. Notons $x_0,z_0,y_0$ les images de
$w_0$ par les morphismes du diagramme. Alors
$$
\xymatrix{
\mathcal{F}_{\mathcal{W}, k, w_0} \ar[d] \ar[r] &
\mathcal{F}_{\mathcal{Z}, k, z_0} \ar[d] \\
\mathcal{F}_{\mathcal{X}, k, x_0} \ar[r] & \mathcal{F}_{\mathcal{Y}, k, y_0}
}
$$
est un produit fibré de catégories de prédéformations.
\end{lemma}

\begin{proof}
Cela résulte directement des définitions. Nous omettons les détails.
\end{proof}






\section{Sommes amalgamées et champs}
\label{section-pushouts}

\noindent
Dans cette section, nous montrons que les champs algébriques se comportent bien
vis-à-vis de certaines sommes amalgamées. Les résultats qui suivent valent sur
un schéma de base quelconque.

\medskip\noindent
Le lemme suivant reste vrai lorsque $Y$, $X'$, $X$ et $Y'$ sont des espaces
algébriques ; voir (insérer ici une référence ultérieure).

\begin{lemma}
\label{lemma-pushout}
\begin{slogan}
Les champs algébriques satisfont la condition (forte) de Rim--Schlessinger
\end{slogan}
Soit $S$ un schéma. Soit
$$
\xymatrix{
X \ar[r] \ar[d] & X' \ar[d] \\
Y \ar[r] & Y'
}
$$
un carré cocartésien dans la catégorie des schémas sur $S$, où $X \to X'$ est
un épaississement et $X \to Y$ est affine ; voir Compléments sur les
morphismes, lemme \ref{more-morphisms-lemma-pushout-along-thickening}. Soit
$\mathcal{Z}$ un champ algébrique sur $S$. Alors le foncteur entre
catégories-fibres
$$
\mathcal{Z}_{Y'}
\longrightarrow
\mathcal{Z}_Y \times_{\mathcal{Z}_X} \mathcal{Z}_{X'}
$$
est une équivalence de catégories.
\end{lemma}

\begin{proof}
Soit $y'$ un objet du membre de gauche. Le faisceau
$\mathit{Isom}(y', y')$ sur la catégorie des schémas au-dessus de $Y'$ est
représentable par un espace algébrique $I$ sur $Y'$ ; voir Champs algébriques,
lemme \ref{algebraic-lemma-representable-diagonal}. Le foncteur du lemme est
donc pleinement fidèle, puisque $Y'$ est la somme amalgamée aussi bien dans la
catégorie des espaces algébriques que dans celle des schémas ; voir Sommes
amalgamées d'espaces, lemme
\ref{spaces-pushouts-lemma-pushout-along-thickening-schemes}.

\medskip\noindent
Soit $(y,x',f)$ un objet du membre de droite, où
$f : y|_X \to x'|_X$ est un isomorphisme. Pour achever la démonstration, nous
devons construire un objet $y'$ de $\mathcal{Z}_{Y'}$ dont les restrictions à
$Y$ et $X'$ s'identifient respectivement à $y$ et $x'$, d'une manière
compatible avec $f$. Il suffit en fait de construire $y'$ localement pour la
topologie fppf sur $Y'$ ; voir Champs, lemme
\ref{stacks-lemma-characterize-essentially-surjective-when-ff}. Choisissons un
champ algébrique représentable $\mathcal{W}$ et un morphisme lisse surjectif
$\mathcal{W} \to \mathcal{Z}$. Alors
$$
(\Sch/Y)_{fppf} \times_{y, \mathcal{Z}} \mathcal{W}
\quad\text{et}\quad
(\Sch/X')_{fppf} \times_{x', \mathcal{Z}} \mathcal{W}
$$
sont des champs algébriques représentés par des espaces algébriques $V$ et
$U'$, lisses respectivement sur $Y$ et $X'$. L'isomorphisme $f$ induit un
isomorphisme $\varphi : V \times_Y X \to U' \times_{X'} X$ au-dessus de $X$.
D'après Sommes amalgamées d'espaces, lemmes
\ref{spaces-pushouts-lemma-pushout-along-thickening} et
\ref{spaces-pushouts-lemma-equivalence-categories-spaces-pushout-flat}
nous voyons que la somme amalgamée $V' = V \amalg_{V \times_Y X} U'$ est un
espace algébrique lisse sur $Y'$ dont les changements de base à $Y$ et $X'$
redonnent $V$ et $U'$ d'une manière compatible avec $\varphi$.

\medskip\noindent
Soit $W$ l'espace algébrique qui représente $\mathcal{W}$. Les projections
$V \to W$ et $U' \to W$ coïncident comme morphismes sur
$V \times_Y X \cong U' \times_{X'} X$ ; la propriété universelle de la somme
amalgamée détermine donc un morphisme d'espaces algébriques $V' \to W$.
Choisissons un schéma $Y_1'$ et un morphisme étale surjectif $Y_1' \to V'$.
Posons $Y_1 = Y \times_{Y'} Y_1'$, $X_1' = X' \times_{Y'} Y_1'$ et
$X_1 = X \times_{Y'} Y_1'$. La composée
$$
(\Sch/Y_1') \to (\Sch/V') \to (\Sch/W) = \mathcal{W} \to \mathcal{Z}
$$
correspond, par le lemme de $2$-Yoneda, à un objet $y_1'$ de $\mathcal{Z}$
au-dessus de $Y_1'$ dont les restrictions à $Y_1$ et $X_1'$ s'identifient à
$y|_{Y_1}$ et $x'|_{X_1'}$ d'une manière compatible avec $f|_{X_1}$. Nous
avons donc construit l'objet recherché localement pour la topologie lisse sur
$Y'$, ce qui achève la démonstration.
\end{proof}








\section{La condition de Rim--Schlessinger}
\label{section-RS}

\noindent
La définition suivante est motivée par le lemme \ref{lemma-pushout} et par
Théorie formelle des déformations, définition
\ref{formal-defos-definition-RS} et lemme
\ref{formal-defos-lemma-RS-2-categorical}.

\begin{definition}
\label{definition-RS}
Soit $S$ un schéma localement noethérien. Soit $\mathcal{Z}$ une catégorie
fibrée en groupoïdes sur $(\Sch/S)_{fppf}$. Nous disons que $\mathcal{Z}$
satisfait la {\it condition (RS)} si, pour tout carré cocartésien
$$
\xymatrix{
X \ar[r] \ar[d] & X' \ar[d] \\
Y \ar[r] & Y' = Y \amalg_X X'
}
$$
dans la catégorie des schémas sur $S$ tel que
\begin{enumerate}
\item $X$, $X'$, $Y$ et $Y'$ soient les spectres d'anneaux locaux artiniens ;
\item $X$, $X'$, $Y$ et $Y'$ soient de type fini sur $S$ ;
\item $X \to X'$ (et donc $Y \to Y'$) soit une immersion fermée,
\end{enumerate}
le foncteur entre catégories-fibres
$$
\mathcal{Z}_{Y'}
\longrightarrow
\mathcal{Z}_Y \times_{\mathcal{Z}_X} \mathcal{Z}_{X'}
$$
est une équivalence de catégories.
\end{definition}

\noindent
Si $A$ est un anneau local artinien de corps résiduel $k$, alors tout morphisme
$\Spec(A) \to S$ est affine et de type fini si et seulement si le morphisme
induit $\Spec(k) \to S$ est de type fini ; voir Morphismes, lemmes
\ref{morphisms-lemma-Artinian-affine} et
\ref{morphisms-lemma-artinian-finite-type}.

\begin{lemma}
\label{lemma-algebraic-stack-RS}
Soit $\mathcal{X}$ un champ algébrique sur une base localement noethérienne
$S$. Alors $\mathcal{X}$ satisfait (RS).
\end{lemma}

\begin{proof}
Cela résulte immédiatement des définitions et du lemme \ref{lemma-pushout}.
\end{proof}

\begin{lemma}
\label{lemma-fibre-product-RS}
Soit $S$ un schéma. Soient $p : \mathcal{X} \to \mathcal{Y}$ et
$q : \mathcal{Z} \to \mathcal{Y}$ des $1$-morphismes de catégories fibrées en
groupoïdes sur $(\Sch/S)_{fppf}$. Si $\mathcal{X}$, $\mathcal{Y}$ et
$\mathcal{Z}$ satisfont (RS), alors
$\mathcal{X} \times_\mathcal{Y} \mathcal{Z}$ satisfait également (RS).
\end{lemma}

\begin{proof}
L'argument est formel. Soit
$$
\xymatrix{
X \ar[r] \ar[d] & X' \ar[d] \\
Y \ar[r] & Y' = Y \amalg_X X'
}
$$
un diagramme comme dans la définition \ref{definition-RS}. Il faut montrer que
$$
(\mathcal{X} \times_{\mathcal{Y}} \mathcal{Z})_{Y'}
\longrightarrow
(\mathcal{X} \times_{\mathcal{Y}} \mathcal{Z})_Y
\times_{(\mathcal{X} \times_{\mathcal{Y}} \mathcal{Z})_X}
(\mathcal{X} \times_{\mathcal{Y}} \mathcal{Z})_{X'}
$$
est une équivalence. En développant la définition du $2$-produit fibré, ce
foncteur devient
\begin{equation}
\label{equation-RS-fibre-product}
\mathcal{X}_{Y'} \times_{\mathcal{Y}_{Y'}} \mathcal{Z}_{Y'}
\longrightarrow
(\mathcal{X}_Y \times_{\mathcal{Y}_Y} \mathcal{Z}_Y)
\times_{(\mathcal{X}_X \times_{\mathcal{Y}_X} \mathcal{Z}_X)}
(\mathcal{X}_{X'} \times_{\mathcal{Y}_{X'}} \mathcal{Z}_{X'}).
\end{equation}
Par hypothèse, chacun des foncteurs
$$
\mathcal{X}_{Y'} \to \mathcal{X}_Y \times_{\mathcal{Y}_Y} \mathcal{Z}_Y,
\quad
\mathcal{Y}_{Y'} \to \mathcal{X}_X \times_{\mathcal{Y}_X} \mathcal{Z}_X,
\quad
\mathcal{Z}_{Y'} \to
\mathcal{X}_{X'} \times_{\mathcal{Y}_{X'}} \mathcal{Z}_{X'}
$$
est une équivalence. Un objet du membre de droite de
(\ref{equation-RS-fibre-product}) est un système
$$
((x_Y, z_Y, \phi_Y), (x_{X'}, z_{X'}, \phi_{X'}), (\alpha, \beta)).
$$
Alors $(x_Y,x_{Y'},\alpha)$ est isomorphe à l'image d'un objet $x_{Y'}$ de
$\mathcal{X}_{Y'}$, et $(z_Y,z_{Y'},\beta)$ est isomorphe à l'image d'un objet
$z_{Y'}$ de $\mathcal{Z}_{Y'}$. Le couple de morphismes
$(\phi_Y,\phi_{X'})$ correspond à un morphisme $\psi$ entre les images de
$x_{Y'}$ et $z_{Y'}$ dans $\mathcal{Y}_{Y'}$. Ainsi,
$(x_{Y'},z_{Y'},\psi)$ est un objet du membre de gauche de
(\ref{equation-RS-fibre-product}) qui s'envoie sur l'objet donné du membre de
droite. Cela montre que (\ref{equation-RS-fibre-product}) est essentiellement
surjectif. Nous omettons la preuve de sa pleine fidélité.
\end{proof}





\section{Catégories de déformations}
\label{section-deformation-categories}

\noindent
Nous faisons correspondre les notations introduites ci-dessus à celles du
chapitre « Théorie formelle des déformations ».

\begin{lemma}
\label{lemma-deformation-category}
Soit $S$ un schéma localement noethérien. Soit $\mathcal{X}$ une catégorie
fibrée en groupoïdes sur $(\Sch/S)_{fppf}$ satisfaisant (RS). Pour tout corps
$k$ de type fini sur $S$ et tout objet $x_0$ de $\mathcal{X}$ au-dessus de
$k$, la catégorie de prédéformations
$p : \mathcal{F}_{\mathcal{X}, k, x_0} \to \mathcal{C}_\Lambda$
(\ref{equation-predeformation-category}) est une catégorie de déformations ;
voir Théorie formelle des déformations, définition
\ref{formal-defos-definition-deformation-category}.
\end{lemma}

\begin{proof}
Posons $\mathcal{F} = \mathcal{F}_{\mathcal{X}, k, x_0}$. Soient
$f_1 : A_1 \to A$ et $f_2 : A_2 \to A$ des homomorphismes d'anneaux dans
$\mathcal{C}_\Lambda$, avec $f_2$ surjectif. Il faut montrer que le foncteur
$$
\mathcal{F}(A_1 \times_A A_2)
\longrightarrow
\mathcal{F}(A_1) \times_{\mathcal{F}(A)} \mathcal{F}(A_2)
$$
est une équivalence ; voir Théorie formelle des déformations, lemme
\ref{formal-defos-lemma-RS-2-categorical}. Posons $X = \Spec(A)$,
$X' = \Spec(A_2)$, $Y = \Spec(A_1)$ et
$Y' = \Spec(A_1 \times_A A_2)$. Remarquons que
$Y' = Y \amalg_X X'$ dans la catégorie des schémas ; voir Compléments sur les
morphismes, lemme \ref{more-morphisms-lemma-pushout-along-thickening}. Dans le
diagramme de foncteurs entre catégories-fibres
$$
\xymatrix{
\mathcal{X}_{Y'} \ar[r] \ar[d] &
\mathcal{X}_Y \times_{\mathcal{X}_X} \mathcal{X}_{X'} \ar[d] \\
\mathcal{X}_{\Spec(k)} \ar@{=}[r] & \mathcal{X}_{\Spec(k)}
}
$$
la flèche horizontale supérieure est une équivalence d'après la définition
\ref{definition-RS}. Or $\mathcal{F}(B)$ est la catégorie des objets de
$\mathcal{X}_{\Spec(B)}$ munis d'une identification à $x_0$ au-dessus de $k$ ;
le résultat s'ensuit.
\end{proof}

\begin{remark}
\label{remark-deformation-category-implies}
Soit $S$ un schéma localement noethérien. Soit $\mathcal{X}$ fibrée en
groupoïdes sur $(\Sch/S)_{fppf}$. Soient $k$ un corps de type fini sur $S$ et
$x_0$ un objet de $\mathcal{X}$ au-dessus de $k$. Soit
$p : \mathcal{F} \to \mathcal{C}_\Lambda$ comme en
(\ref{equation-predeformation-category}). Si $\mathcal{F}$ est une catégorie
de déformations, c'est-à-dire si elle satisfait la condition (RS) de
Rim--Schlessinger, alors $\mathcal{F}$ satisfait les conditions (S1) et (S2) de
Schlessinger d'après Théorie formelle des déformations, lemme
\ref{formal-defos-lemma-RS-implies-S1-S2}. Soit
$\overline{\mathcal{F}}$ le foncteur des classes d'isomorphisme ; voir Théorie
formelle des déformations, remarques
\ref{formal-defos-remarks-cofibered-groupoids}
(\ref{formal-defos-item-associated-functor-isomorphism-classes}).
Alors $\overline{\mathcal{F}}$ satisfait lui aussi (S1) et (S2) ; voir Théorie
formelle des déformations, lemme
\ref{formal-defos-lemma-S1-S2-associated-functor}. Cela vaut en particulier
dans la situation du lemme \ref{lemma-deformation-category}.
\end{remark}




\section{Changement de corps}
\label{section-change-of-field}

\noindent
Cette section est l'analogue de Théorie formelle des déformations, section
\ref{formal-defos-section-change-of-field}. Comme il y est indiqué, pour
étudier le changement de corps, nous devons écrire
$\mathcal{C}_{\Lambda,k}$ au lieu de $\mathcal{C}_\Lambda$. Dans le lemme
suivant, nous utilisons la notation $\mathcal{F}_{l/k}$ introduite dans Théorie
formelle des déformations, situation
\ref{formal-defos-situation-change-of-fields}.

\begin{lemma}
\label{lemma-change-of-field}
Soit $S$ un schéma localement noethérien. Soit $\mathcal{X}$ une catégorie
fibrée en groupoïdes sur $(\Sch/S)_{fppf}$. Soient $k$ un corps de type fini
sur $S$ et $l/k$ une extension finie. Soit $x_0$ un objet de $\mathcal{F}$
au-dessus de $\Spec(k)$. Notons $x_{l,0}$ la restriction de $x_0$ à
$\Spec(l)$. Il existe alors un foncteur canonique
$$
(\mathcal{F}_{\mathcal{X}, k , x_0})_{l/k}
\longrightarrow
\mathcal{F}_{\mathcal{X}, l, x_{l, 0}}
$$
de catégories cofibrées en groupoïdes sur $\mathcal{C}_{\Lambda,l}$. Si
$\mathcal{X}$ satisfait (RS), ce foncteur est une équivalence.
\end{lemma}

\begin{proof}
Considérons une factorisation
$$
\Spec(l) \to \Spec(B) \to S
$$
comme en (\ref{equation-factor}). Par définition,
$$
(\mathcal{F}_{\mathcal{X}, k , x_0})_{l/k}(B) =
\mathcal{F}_{\mathcal{X}, k, x_0}(B \times_l k)
$$
voir Théorie formelle des déformations, situation
\ref{formal-defos-situation-change-of-fields}. Un objet de cette catégorie est
donc un morphisme $x_0 \to x$ de $\mathcal{X}$ situé au-dessus du morphisme
$\Spec(k) \to \Spec(B \times_l k)$. Après avoir choisi un foncteur image
réciproque pour $\mathcal{X}$, nous pouvons associer à $x_0 \to x$ le
morphisme $x_{l,0} \to x_B$, où $x_B$ est la restriction de $x$ à
$\Spec(B)$, par le morphisme $\Spec(B) \to \Spec(B \times_l k)$ issu de
$B \times_l k \subset B$. Cette construction est fonctorielle en $B$ et
compatible aux morphismes.

\medskip\noindent
Supposons maintenant que $\mathcal{X}$ satisfasse (RS). Considérons les
diagrammes
$$
\vcenter{
\xymatrix{
l & B \ar[l] \\
k \ar[u] & B \times_l k \ar[l] \ar[u]
}
}
\quad\text{and}\quad
\vcenter{
\xymatrix{
\Spec(l) \ar[d] \ar[r] & \Spec(B) \ar[d] \\
\Spec(k) \ar[r] & \Spec(B \times_l k)
}
}
$$
Le diagramme de gauche est un produit fibré d'anneaux. Celui de droite est un
carré cocartésien dans la catégorie des schémas ; voir Compléments sur les
morphismes, lemme \ref{more-morphisms-lemma-pushout-along-thickening}. Tous ces
schémas sont de type fini sur $S$ (voir les remarques qui suivent la définition
\ref{definition-RS}). La condition (RS) donne donc une équivalence de
catégories-fibres
$$
\mathcal{X}_{\Spec(B \times_l k)}
\longrightarrow
\mathcal{X}_{\Spec(k)}
\times_{\mathcal{X}_{\Spec(l)}}
\mathcal{X}_{\Spec(B)}
$$
Il s'ensuit que le foncteur défini ci-dessus induit une équivalence sur les
catégories-fibres. Il est donc une équivalence de catégories cofibrées en
groupoïdes, par le dual de Catégories, lemme
\ref{categories-lemma-equivalence-fibred-categories}.
\end{proof}








\section{Espaces tangents}
\label{section-tangent-spaces}

\noindent
Soit $S$ un schéma localement noethérien. Soit $\mathcal{X}$ une catégorie
fibrée en groupoïdes sur $(\Sch/S)_{fppf}$. Soit $k$ un corps de type fini
sur $S$ et soit $x_0$ un objet de $\mathcal{X}$ sur $k$.
Dans Théorie formelle des déformations, section
\ref{formal-defos-section-tangent-spaces}, nous avons défini l'{\it espace
tangent}
\begin{equation}
\label{equation-tangent-space}
T\mathcal{F}_{\mathcal{X}, k, x_0} =
\left\{
\begin{matrix}
\text{classes d'isomorphisme de morphismes}\\
x_0 \to x\text{ au-dessus de }\Spec(k) \to \Spec(k[\epsilon])
\end{matrix}
\right\}
\end{equation}
de la catégorie de prédéformations $\mathcal{F}_{\mathcal{X}, k, x_0}$.
Dans Théorie formelle des déformations, section
\ref{formal-defos-section-infinitesimal-automorphisms}, nous avons défini
\begin{equation}
\label{equation-infinitesimal-automorphisms}
\text{Inf}(\mathcal{F}_{\mathcal{X}, k, x_0}) =
\Ker\left(
\text{Aut}_{\Spec(k[\epsilon])}(x'_0) \to \text{Aut}_{\Spec(k)}(x_0)
\right)
\end{equation}
où $x_0'$ est l'image réciproque de $x_0$ sur $\Spec(k[\epsilon])$.
Si $\mathcal{X}$ satisfait la condition de Rim--Schlessinger (RS), alors
$T\mathcal{F}_{\mathcal{X}, k, x_0}$ est naturellement muni d'une structure
d'espace vectoriel sur $k$ d'après Théorie formelle des déformations, lemme
\ref{formal-defos-lemma-tangent-space-vector-space}
(les hypothèses sont satisfaites d'après le lemme
\ref{lemma-deformation-category} et la remarque
\ref{remark-deformation-category-implies}). De plus, Théorie formelle des
déformations, lemme \ref{formal-defos-lemma-infaut-vector-space}, montre que
$\text{Inf}(\mathcal{F}_{\mathcal{X}, k, x_0})$ possède une structure
naturelle d'espace vectoriel sur $k$ pour laquelle l'addition correspond à la
composition des automorphismes. Une condition naturelle consiste à demander
que ces espaces vectoriels soient de dimension finie.

\medskip\noindent
Le lemme suivant affirme qu'il en est ainsi lorsque
$\mathcal{X}$ est localement de type fini sur $S$ (voir Morphismes de champs,
section \ref{stacks-morphisms-section-finite-type}).

\begin{lemma}
\label{lemma-finite-dimension}
Soit $S$ un schéma localement noethérien. Supposons que
\begin{enumerate}
\item $\mathcal{X}$ soit un champ algébrique,
\item $U$ soit un schéma localement de type fini sur $S$, et
\item $(\Sch/U)_{fppf} \to \mathcal{X}$ soit un morphisme lisse et
surjectif.
\end{enumerate}
Alors, pour tout $\mathcal{F} = \mathcal{F}_{\mathcal{X}, k, x_0}$ comme dans
la section \ref{section-predeformation-categories}, l'espace tangent
$T\mathcal{F}$ et l'espace des automorphismes infinitésimaux
$\text{Inf}(\mathcal{F})$ sont de dimension finie sur $k$.
\end{lemma}

\begin{proof}
Posons $\mathcal{U} = (\Sch/U)_{fppf}$. D'après notre définition des champs
algébriques, le $1$-morphisme $\mathcal{U} \to \mathcal{X}$ est représentable
par des espaces algébriques. Par conséquent, le 2-produit fibré
$$
\mathcal{U}_{x_0} = (\Sch/\Spec(k))_{fppf} \times_\mathcal{X} \mathcal{U}
$$
est en particulier représentable par un espace algébrique $U_{x_0}$ sur
$\Spec(k)$. Le morphisme $U_{x_0} \to \Spec(k)$ est alors lisse et surjectif
(en particulier, $U_{x_0}$ est non vide). D'après Espaces sur les corps, lemme
\ref{spaces-over-fields-lemma-smooth-separable-closed-points-dense}, il existe
une extension finie $l/k$ et un point $\Spec(l) \to U_{x_0}$ au-dessus de
$k$. Nous avons
$$
(\mathcal{F}_{\mathcal{X}, k , x_0})_{l/k} =
\mathcal{F}_{\mathcal{X}, l, x_{l, 0}}
$$
d'après le lemme \ref{lemma-change-of-field} et le fait que $\mathcal{X}$
satisfait (RS). Nous voyons donc que
$$
T\mathcal{F} \otimes_k l \cong T\mathcal{F}_{\mathcal{X}, l, x_{l, 0}}
\quad\text{and}\quad
\text{Inf}(\mathcal{F}) \otimes_k l \cong
\text{Inf}(\mathcal{F}_{\mathcal{X}, l, x_{l, 0}})
$$
d'après Théorie formelle des déformations, lemmes
\ref{formal-defos-lemma-tangent-space-change-of-field} et
\ref{formal-defos-lemma-inf-aut-change-of-field}
(ceux-ci sont applicables d'après les lemmes
\ref{lemma-algebraic-stack-RS} et \ref{lemma-deformation-category}, et la
remarque \ref{remark-deformation-category-implies}). Il suffit donc de montrer
que $T\mathcal{F}_{\mathcal{X}, l, x_{l, 0}}$ et
$\text{Inf}(\mathcal{F}_{\mathcal{X}, l, x_{l, 0}})$ sont de dimension finie
sur $l$. Notons que $x_{l, 0}$ provient d'un point $u_0$ de $\mathcal{U}$ sur
$l$.

\medskip\noindent
Interrompons l'argument pour montrer que le lemme relatif aux automorphismes
infinitésimaux résulte du lemme relatif aux espaces tangents. Posons
$\mathcal{R} = \mathcal{U} \times_\mathcal{X} \mathcal{U}$.
Soit $r_0$ le point à valeurs dans $l$, $(u_0, u_0, \text{id}_{x_0})$, de
$\mathcal{R}$. En combinant le lemme
\ref{lemma-fibre-product-deformation-categories} et Théorie formelle des
déformations, lemme \ref{formal-defos-lemma-deformation-functor-diagonal},
nous voyons que
$$
\text{Inf}(\mathcal{F}_{\mathcal{X}, l, x_{l, 0}})
\subset
T\mathcal{F}_{\mathcal{R}, l, r_0}
$$
Notons que $\mathcal{R}$ est un champ algébrique ; voir Champs algébriques,
lemme \ref{algebraic-lemma-2-fibre-product-general}. De plus, $\mathcal{R}$ est
représentable par un espace algébrique $R$ lisse sur $U$ (par l'une ou l'autre
projection ; voir Champs algébriques, lemme
\ref{algebraic-lemma-stack-presentation}). Choisissons donc un schéma $U'$ et
un morphisme étale surjectif $U' \to R$. Alors $U'$ est lisse sur $U$, donc
localement de type fini sur $S$. Comme $(\Sch/U')_{fppf} \to \mathcal{R}$ est
surjectif et lisse, nous avons ramené la question au cas des espaces tangents.

\medskip\noindent
Le foncteur (\ref{equation-functoriality})
$$
\mathcal{F}_{\mathcal{U}, l, u_0}
\longrightarrow
\mathcal{F}_{\mathcal{X}, l, x_{l, 0}}
$$
est lisse d'après le lemme
\ref{lemma-formally-smooth-on-deformation-categories}. Le morphisme induit sur
les espaces tangents
$$
T\mathcal{F}_{\mathcal{U}, l, u_0}
\longrightarrow
T\mathcal{F}_{\mathcal{X}, l, x_{l, 0}}
$$
est $l$-linéaire (d'après Théorie formelle des déformations, lemme
\ref{formal-defos-lemma-k-linear-differential}) et surjectif (car les
morphismes lisses de catégories de prédéformations induisent des applications
surjectives sur les espaces tangents d'après Théorie formelle des déformations,
lemme \ref{formal-defos-lemma-smooth-morphism-essentially-surjective}). Il
suffit donc de montrer que l'espace tangent de l'espace de déformations associé
au champ algébrique représentable $\mathcal{U}$ au point $u_0$ est de dimension
finie. Soit $\Spec(R) \subset U$ un ouvert affine tel que
$u_0 : \Spec(l) \to U$ se factorise par $\Spec(R)$ et que
$\Spec(R) \to S$ se factorise par $\Spec(\Lambda) \subset S$. Soit
$\mathfrak m_R \subset R$ le noyau de l'homomorphisme de $\Lambda$-algèbres
$\varphi_0 : R \to l$ correspondant à $u_0$. Puisque $R$ est de type fini sur
l'anneau noethérien $\Lambda$, c'est un anneau noethérien. Par conséquent,
$\mathfrak m_R = (f_1, \ldots, f_n)$ est un idéal de type fini. Nous avons
$$
T\mathcal{F}_{\mathcal{U}, l, u_0}
=
\{\varphi : R \to l[\epsilon] \mid
\varphi \text{ est un homomorphisme de } \Lambda\text{-algèbres et }
\varphi \bmod \epsilon = \varphi_0\}
$$
Un élément du second membre est déterminé par ses valeurs en
$f_1, \ldots, f_n$ ; la dimension est donc au plus $n$, ce qui conclut.
Nous omettons certains détails.
\end{proof}

\begin{lemma}
\label{lemma-fibre-product-tangent-spaces}
Soit $S$ un schéma localement noethérien. Soient
$p : \mathcal{X} \to \mathcal{Y}$ et $q : \mathcal{Z} \to \mathcal{Y}$ des
$1$-morphismes de catégories fibrées en groupoïdes sur $(\Sch/S)_{fppf}$.
Supposons que $\mathcal{X}$, $\mathcal{Y}$ et $\mathcal{Z}$ satisfassent (RS).
Soit $k$ un corps de type fini sur $S$ et soit $w_0$ un objet de
$\mathcal{W} = \mathcal{X} \times_\mathcal{Y} \mathcal{Z}$ sur $k$. Notons
$x_0, y_0, z_0$ les objets de $\mathcal{X}, \mathcal{Y}, \mathcal{Z}$ obtenus
à partir de $w_0$. Il existe alors une suite exacte à $6$ termes
$$
\xymatrix{
0 \ar[r] &
\text{Inf}(\mathcal{F}_{\mathcal{W}, k, w_0}) \ar[r] &
\text{Inf}(\mathcal{F}_{\mathcal{X}, k, x_0}) \oplus
\text{Inf}(\mathcal{F}_{\mathcal{Z}, k, z_0}) \ar[r] &
\text{Inf}(\mathcal{F}_{\mathcal{Y}, k, y_0}) \ar[lld] \\
 &
T\mathcal{F}_{\mathcal{W}, k, w_0} \ar[r] &
T\mathcal{F}_{\mathcal{X}, k, x_0} \oplus
T\mathcal{F}_{\mathcal{Z}, k, z_0} \ar[r] &
T\mathcal{F}_{\mathcal{Y}, k, y_0}
}
$$
d'espaces vectoriels sur $k$.
\end{lemma}

\begin{proof}
D'après le lemme \ref{lemma-fibre-product-RS}, la catégorie $\mathcal{W}$
satisfait (RS), de sorte que l'énoncé a un sens. Pour le démontrer, appliquons
les lemmes \ref{lemma-fibre-product-deformation-categories} et
\ref{lemma-deformation-category}, ainsi que Théorie formelle des déformations,
lemme \ref{formal-defos-lemma-deformation-categories-fiber-product-morphisms}.
\end{proof}






\section{Objets formels}
\label{section-formal-objects}

\noindent
Dans cette section, nous transposons au cadre présent certaines des notions
déjà définies au chapitre « Théorie formelle des déformations ».
Dans la suite, nous dirons que « $R$ est une $S$-algèbre » pour indiquer
que $R$ est un anneau muni d'un morphisme de schémas $\Spec(R) \to S$.

\begin{definition}
\label{definition-formal-objects}
Soit $S$ un schéma localement noethérien. Soit
$p : \mathcal{X} \to (\Sch/S)_{fppf}$ une catégorie fibrée en groupoïdes.
\begin{enumerate}
\item Un {\it objet formel} $\xi = (R, \xi_n, f_n)$ de $\mathcal{X}$ consiste
en une $S$-algèbre locale noethérienne complète $R$, des objets $\xi_n$ de
$\mathcal{X}$ situés au-dessus de $\Spec(R/\mathfrak m_R^n)$ et des morphismes
$f_n : \xi_n \to \xi_{n + 1}$ de $\mathcal{X}$ situés au-dessus de
$\Spec(R/\mathfrak m^n) \to \Spec(R/\mathfrak m^{n + 1})$,
tels que $R/\mathfrak m$ soit un corps de type fini sur $S$.
\item Un {\it morphisme d'objets formels}
$a : \xi = (R, \xi_n, f_n) \to \eta = (T, \eta_n, g_n)$
est donné par des morphismes $a_n : \xi_n \to \eta_n$ tels que, pour tout
$n$, le diagramme
$$
\xymatrix{
\xi_n \ar[r]_{f_n} \ar[d]_{a_n} & \xi_{n + 1} \ar[d]^{a_{n + 1}} \\
\eta_n \ar[r]^{g_n} & \eta_{n + 1}
}
$$
est commutatif. En appliquant le foncteur $p$, nous obtenons une famille
compatible de morphismes
$\Spec(R/\mathfrak m_R^n) \to \Spec(T/\mathfrak m_T^n)$ et donc un morphisme
$a_0 : \Spec(R) \to \Spec(T)$ au-dessus de $S$. Nous disons que $a$ est
{\it au-dessus de} $a_0$.
\end{enumerate}
\end{definition}

\noindent
Nous obtenons ainsi une catégorie des objets formels de $\mathcal{X}$.

\begin{remark}
\label{remark-formal-objects-match}
Soit $S$ un schéma localement noethérien. Soit
$p : \mathcal{X} \to (\Sch/S)_{fppf}$ une catégorie fibrée en groupoïdes.
Soit $\xi = (R, \xi_n, f_n)$ un objet formel. Posons
$k = R/\mathfrak m$ et $x_0 = \xi_1$. L'objet formel $\xi$ définit un objet
formel $\xi$ de la catégorie de prédéformations
$\mathcal{F}_{\mathcal{X}, k, x_0}$. Cela résulte immédiatement de la
définition \ref{definition-formal-objects} ci-dessus, de Théorie formelle des
déformations, définition \ref{formal-defos-definition-formal-objects}, et de
notre construction de la catégorie de prédéformations
$\mathcal{F}_{\mathcal{X}, k, x_0}$ dans la section
\ref{section-predeformation-categories}.
\end{remark}

\noindent
Si $F : \mathcal{X} \to \mathcal{Y}$ est un $1$-morphisme de catégories
fibrées en groupoïdes sur $(\Sch/S)_{fppf}$, alors $F$ induit également un
foncteur entre les catégories d'objets formels.

\begin{lemma}
\label{lemma-smooth-lift-formal}
Soit $S$ un schéma localement noethérien. Soit
$F : \mathcal{X} \to \mathcal{Y}$ un $1$-morphisme de catégories fibrées en
groupoïdes sur $(\Sch/S)_{fppf}$. Soit $\eta = (R, \eta_n, g_n)$ un objet
formel de $\mathcal{Y}$ et soit $\xi_1$ un objet de $\mathcal{X}$ tel que
$F(\xi_1) \cong \eta_1$. Si $F$ est formellement lisse sur les objets (voir
Critères de représentabilité, section
\ref{criteria-section-formally-smooth}), alors il existe un objet formel
$\xi = (R, \xi_n, f_n)$ de $\mathcal{X}$ tel que $F(\xi) \cong \eta$.
\end{lemma}

\begin{proof}
Notons que chacun des morphismes
$\Spec(R/\mathfrak m^n) \to \Spec(R/\mathfrak m^{n + 1})$ est un
épaississement du premier ordre de schémas affines sur $S$. L'hypothèse sur
$F$ signifie donc que nous pouvons relever successivement $\xi_1$ en des
objets $\xi_2, \xi_3, \ldots$ de $\mathcal{X}$ munis d'isomorphismes
compatibles
$\eta_n|_{\Spec(R/\mathfrak m^{n - 1})} \cong \eta_{n - 1}$
et $F(\eta_n) \cong \xi_n$.
\end{proof}

\noindent
Soit $S$ un schéma localement noethérien. Soit
$p : \mathcal{X} \to (\Sch/S)_{fppf}$ une catégorie fibrée en groupoïdes.
Supposons que $x$ soit un objet de $\mathcal{X}$ sur $R$, où $R$ est une
$S$-algèbre locale noethérienne complète dont le corps résiduel est de type
fini sur $S$. Nous pouvons alors considérer le système des restrictions
$\xi_n = x|_{\Spec(R/\mathfrak m^n)}$, muni des morphismes naturels
$\xi_1 \to \xi_2 \to \ldots$ provenant de la transitivité de la restriction.
Ainsi, $\xi = (R, \xi_n, \xi_n \to \xi_{n + 1})$ est un objet formel de
$\mathcal{X}$. Cette construction est fonctorielle en l'objet $x$.
Nous obtenons donc un foncteur
\begin{equation}
\label{equation-approximation}
\left\{
\begin{matrix}
\text{objets }x\text{ de }\mathcal{X} \text{ tels que }p(x) = \Spec(R) \\
\text{où }R\text{ est locale noethérienne complète}\\
\text{avec }R/\mathfrak m\text{ de type fini sur }S
\end{matrix}
\right\}
\longrightarrow
\left\{
\begin{matrix}
\text{objets formels de }\mathcal{X}
\end{matrix}
\right\}
\end{equation}
Précisément, le membre de gauche est la sous-catégorie pleine de
$\mathcal{X}$ formée des objets indiqués et le membre de droite est la
catégorie des objets formels de $\mathcal{X}$ de la définition
\ref{definition-formal-objects}.

\begin{definition}
\label{definition-effective}
Soit $S$ un schéma localement noethérien. Soit $\mathcal{X}$ une catégorie
fibrée en groupoïdes sur $(\Sch/S)_{fppf}$. Un objet formel
$\xi = (R, \xi_n, f_n)$ de $\mathcal{X}$ est dit {\it effectif} s'il
appartient à l'image essentielle du foncteur
(\ref{equation-approximation}).
\end{definition}

\noindent
Si la catégorie fibrée en groupoïdes est un champ algébrique, alors tout
objet formel est effectif, comme il résulte du lemme suivant.

\begin{lemma}
\label{lemma-effective}
Soit $S$ un schéma localement noethérien. Soit $\mathcal{X}$ un champ
algébrique sur $S$. Le foncteur (\ref{equation-approximation}) est une
équivalence.
\end{lemma}

\begin{proof}
Cas I : $\mathcal{X}$ est représentable (par un schéma). Écrivons
$\mathcal{X} = (\Sch/X)_{fppf}$ pour un schéma $X$ sur $S$.
En développant les définitions, nous devons démontrer l'énoncé suivant : pour
une $S$-algèbre locale noethérienne complète $R$ telle que
$R/\mathfrak m$ soit de type fini sur $S$, l'application
$$
\Mor_S(\Spec(R), X) \longrightarrow \lim \Mor_S(\Spec(R/\mathfrak m^n), X)
$$
est bijective. Cela résulte d'Espaces formels, lemme
\ref{formal-spaces-lemma-map-into-scheme}.

\medskip\noindent
Cas II. $\mathcal{X}$ est représentable par un espace algébrique. Écrivons
que $\mathcal{X}$ est représentable par $X$. Il nous faut de nouveau montrer
que
$$
\Mor_S(\Spec(R), X) \longrightarrow \lim \Mor_S(\Spec(R/\mathfrak m^n), X)
$$
est bijective pour $R$ comme ci-dessus. C'est Espaces formels, lemme
\ref{formal-spaces-lemma-map-into-algebraic-space}.

\medskip\noindent
Cas III : cas général d'un champ algébrique. Remarquons en général que les
membres de gauche et de droite de (\ref{equation-approximation}) sont des
catégories fibrées en groupoïdes sur la catégorie des schémas affines sur
$S$ qui sont les spectres d'anneaux locaux noethériens complets dont le corps
résiduel est de type fini sur $S$. Nous verrons aussi dans la démonstration
ci-dessous qu'ils forment des champs pour une certaine topologie sur cette
catégorie.

\medskip\noindent
Montrons d'abord la pleine fidélité. Soit $R$ une $S$-algèbre locale
noethérienne complète telle que $k = R/\mathfrak m$ soit de type fini sur
$S$. Soient $x, x'$ des objets de $\mathcal{X}$ sur $R$. Comme
$\mathcal{X}$ est un champ algébrique, $\mathit{Isom}(x, x')$ est
représentable par un espace algébrique $I$ sur $\Spec(R)$ ; voir Champs
algébriques, lemme \ref{algebraic-lemma-representable-diagonal}.
L'application du cas II à $I$ sur $\Spec(R)$ implique immédiatement que
(\ref{equation-approximation}) est pleinement fidèle sur les catégories-fibres
au-dessus de $\Spec(R)$. Le foncteur est donc pleinement fidèle d'après
Catégories, lemme \ref{categories-lemma-equivalence-fibred-categories}.

\medskip\noindent
Surjectivité essentielle. Soit $\xi = (R, \xi_n, f_n)$ un objet formel de
$\mathcal{X}$. Choisissons un schéma $U$ sur $S$ et un morphisme lisse
surjectif $f : (\Sch/U)_{fppf} \to \mathcal{X}$. Pour tout $n$, considérons
le produit fibré
$$
(\Sch/\Spec(R/\mathfrak m^n))_{fppf}
\times_{\xi_n, \mathcal{X}, f}
(\Sch/U)_{fppf}
$$
Par hypothèse, il est représentable par un espace algébrique $V_n$ lisse et
surjectif sur $\Spec(R/\mathfrak m^n)$. Les morphismes
$f_n : \xi_n \to \xi_{n + 1}$ induisent des carrés cartésiens
$$
\xymatrix{
V_{n + 1} \ar[d] & V_n \ar[d] \ar[l] \\
\Spec(R/\mathfrak m^{n + 1}) & \Spec(R/\mathfrak m^n) \ar[l]
}
$$
 d'espaces algébriques. D'après Espaces sur les corps, lemme
\ref{spaces-over-fields-lemma-smooth-separable-closed-points-dense}, il existe
une extension séparable finie $k'/k$ et un point
$v'_1 : \Spec(k') \to V_1$ au-dessus de $k$. Soit $R \subset R'$ l'extension
étale finie dont l'extension des corps résiduels est $k'/k$ (elle existe et est
unique d'après Algèbre, lemmes
\ref{algebra-lemma-henselian-cat-finite-etale} et
\ref{algebra-lemma-complete-henselian}). En appliquant le critère infinitésimal
de lissité (voir Compléments sur les morphismes d'espaces, lemme
\ref{spaces-more-morphisms-lemma-smooth-formally-smooth}) à
$V_n \to \Spec(R/\mathfrak m^n)$ pour $n = 2, 3, 4, \ldots$, nous pouvons
trouver successivement des morphismes
$v'_n : \Spec(R'/(\mathfrak m')^n) \to V_n$ au-dessus de
$\Spec(R/\mathfrak m^n)$, s'insérant dans des diagrammes commutatifs
$$
\xymatrix{
\Spec(R'/(\mathfrak m')^{n + 1}) \ar[d]_{v'_{n + 1}} &
\Spec(R'/(\mathfrak m')^n) \ar[d]^{v'_n} \ar[l] \\
V_{n + 1} & V_n \ar[l]
}
$$
En composant avec les morphismes de projection $V_n \to U$, nous obtenons un
système compatible de morphismes
$u'_n : \Spec(R'/(\mathfrak m')^n) \to U$. D'après le cas I, la famille
$(u'_n)$ provient d'un unique morphisme $u' : \Spec(R') \to U$. Notons $x'$
l'objet de $\mathcal{X}$ sur $\Spec(R')$ obtenu en appliquant le $1$-morphisme
$f$ à $u'$. Par construction, il existe un morphisme d'objets formels
$$
(\ref{equation-approximation})(x') =
(R', x'|_{\Spec(R'/(\mathfrak m')^n)}, \ldots)
\longrightarrow
(R, \xi_n, f_n)
$$
au-dessus de $\Spec(R') \to \Spec(R)$. Notons que $R' \otimes_R R'$ est un
produit fini de spectres d'anneaux locaux noethériens complets auxquels
s'applique la présente discussion. Notons
$p_0, p_1 : \Spec(R' \otimes_R R') \to \Spec(R')$ les deux projections. Par la
pleine fidélité démontrée ci-dessus, il existe un isomorphisme canonique
$\varphi : p_0^*x' \to p_1^*x'$, car nous disposons de tels isomorphismes sur
$\Spec((R' \otimes_R R')/\mathfrak m^n(R' \otimes_R R'))$. Nous omettons de
vérifier que l'isomorphisme $\varphi$ satisfait la condition de cocycle (voir
Champs, définition \ref{stacks-definition-descent-data}). Puisque
$\{\Spec(R') \to \Spec(R)\}$ est un recouvrement fppf, nous concluons que
$x'$ descend en un objet $x$ de $\mathcal{X}$ sur $\Spec(R)$. Nous omettons de
vérifier que $x_n$ est la restriction de $x$ à
$\Spec(R/\mathfrak m^n)$.
\end{proof}

\begin{lemma}
\label{lemma-fibre-product-effective}
Soit $S$ un schéma. Soient $p : \mathcal{X} \to \mathcal{Y}$ et
$q : \mathcal{Z} \to \mathcal{Y}$ des $1$-morphismes de catégories fibrées
en groupoïdes sur $(\Sch/S)_{fppf}$. Si le foncteur
(\ref{equation-approximation}) est une équivalence pour
$\mathcal{X}$, $\mathcal{Y}$ et $\mathcal{Z}$, alors il en est une pour
$\mathcal{X} \times_\mathcal{Y} \mathcal{Z}$.
\end{lemma}

\begin{proof}
Les membres de gauche et de droite de (\ref{equation-approximation}) pour
$\mathcal{X} \times_\mathcal{Y} \mathcal{Z}$ sont simplement les
$2$-produits fibrés des membres de gauche et de droite de
(\ref{equation-approximation}) pour $\mathcal{X}$ et $\mathcal{Z}$ au-dessus
de $\mathcal{Y}$. Le résultat en découle, puisque la formation des
$2$-produits fibrés est compatible avec les équivalences de catégories ; voir
Catégories, lemme \ref{categories-lemma-equivalence-2-fibre-product}.
\end{proof}






\section{Approximation}
\label{section-approximation}

\noindent
Une idée fondamentale de Michael Artin est que l'on peut approcher les objets
d'un champ qui commute aux limites. Plus précisément, étant donné un objet
$x$ du champ sur un anneau local noethérien complet, on peut trouver un objet
$x_A$ sur un anneau algébrique qui est « proche » de $x$. Ici, un anneau
algébrique désigne une $S$-algèbre de type fini, et proche signifie proche pour
la topologie adique. Nous en donnons dans cette section une forme simple mais
générale.

\medskip\noindent
Pour formuler le résultat, il nous faut rassembler quelques définitions
provenant de différents endroits du projet Stacks. Tout d'abord, dans
Critères de représentabilité, section
\ref{criteria-section-limit-preserving}, nous avons introduit la propriété
{\it de commuter aux limites sur les objets} pour les $1$-morphismes de
catégories fibrées en groupoïdes sur la catégorie des schémas. Dans Compléments
d'algèbre, définition \ref{more-algebra-definition-G-ring}, nous avons défini
la notion d'{\it anneau G}. Soit $S$ un schéma localement noethérien. Soit $A$
une $S$-algèbre. Nous disons que $A$ est {\it de type fini sur $S$}, ou que
c'est une {\it $S$-algèbre de type fini}, si $\Spec(A) \to S$ est de type
fini. Dans ce cas, $A$ est un anneau noethérien. Enfin, pour un anneau $A$ et
un idéal $I$, nous notons $\text{Gr}_I(A) = \bigoplus I^n/I^{n + 1}$.

\begin{lemma}
\label{lemma-approximate}
Soit $S$ un schéma localement noethérien. Soit
$p : \mathcal{X} \to (\Sch/S)_{fppf}$ une catégorie fibrée en groupoïdes.
Soit $x$ un objet de $\mathcal{X}$ situé au-dessus de $\Spec(R)$, où $R$ est
un anneau local noethérien complet dont le corps résiduel $k$ est de type fini
sur $S$. Soit $s \in S$ l'image de $\Spec(k) \to S$. Supposons que
(a) $\mathcal{O}_{S, s}$ soit un anneau G et que (b) $p$ commute aux limites
sur les objets. Alors, pour tout entier $N \geq 1$, il existe
\begin{enumerate}
\item une $S$-algèbre de type fini $A$,
\item un idéal maximal $\mathfrak m_A \subset A$,
\item un objet $x_A$ de $\mathcal{X}$ sur $\Spec(A)$,
\item un $S$-isomorphisme $R/\mathfrak m_R^N \cong A/\mathfrak m_A^N$,
\item un isomorphisme
$x|_{\Spec(R/\mathfrak m_R^N)} \cong x_A|_{\Spec(A/\mathfrak m_A^N)}$
compatible à (4), et
\item un isomorphisme
$\text{Gr}_{\mathfrak m_R}(R) \cong \text{Gr}_{\mathfrak m_A}(A)$
de $k$-algèbres graduées.
\end{enumerate}
\end{lemma}

\begin{proof}
Choisissons un ouvert affine $\Spec(\Lambda) \subset S$ tel que $k$ soit une
$\Lambda$-algèbre finie ; voir Morphismes, lemme
\ref{morphisms-lemma-point-finite-type}. Nous pouvons remplacer, et remplaçons,
$S$ par $\Spec(\Lambda)$.

\medskip\noindent
Nous pouvons écrire $R$ comme une colimite filtrante $R = \colim C_j$, où
chaque $C_j$ est une $\Lambda$-algèbre de type fini (voir Algèbre, lemme
\ref{algebra-lemma-ring-colimit-fp}). D'après l'hypothèse (b), l'objet $x$ est
isomorphe à la restriction d'un objet sur l'une des $C_j$. Nous pouvons donc
choisir une $\Lambda$-algèbre de type fini $C$, un homomorphisme de
$\Lambda$-algèbres $C \to R$ et un objet $x_C$ de $\mathcal{X}$ sur
$\Spec(C)$ tels que $x = x_C|_{\Spec(R)}$. Le choix de $C$ est un artifice de
tenue des comptes dont on pourrait se passer. Pour la suite, écrivons
$C = \Lambda[y_1, \ldots, y_u]/(f_1, \ldots, f_v)$ et notons
$\overline{a}_i \in R$ l'image de $y_i$ par l'homomorphisme $C \to R$.
Posons $\mathfrak m_C = C \cap \mathfrak m_R$.

\medskip\noindent
Choisissons une surjection de $\Lambda$-algèbres
$\Lambda[x_1, \ldots, x_s] \to k$ et notons $\mathfrak m'$ son noyau.
La propriété universelle des anneaux de polynômes permet de la relever en un
homomorphisme de $\Lambda$-algèbres
$\Lambda[x_1, \ldots, x_s] \to R$. Ajoutons quelques variables (c'est-à-dire
augmentons un peu $s$) dont les images engendrent $\mathfrak m_R$. Ceci fait,
$\Lambda[x_1, \ldots, x_s] \to R/\mathfrak m_R^2$ est surjectif. Par suite,
\begin{equation}
\label{equation-surjection}
P = \Lambda[x_1, \ldots, x_s]_{\mathfrak m'}^\wedge \longrightarrow R
\end{equation}
est un homomorphisme surjectif d'anneaux locaux noethériens complets ; voir par
exemple Théorie formelle des déformations, lemme
\ref{formal-defos-lemma-surjective-cotangent-space}.

\medskip\noindent
Choisissons des relèvements $a_i \in P$ des $\overline{a}_i$ obtenus
ci-dessus. Choisissons des générateurs $b_1, \ldots, b_r \in P$ du noyau de
(\ref{equation-surjection}). Choisissons $c_{ji} \in P$ tels que
$$
f_j(a_1, \ldots, a_u) = \sum c_{ji} b_i
$$
dans $P$, ce qui est possible d'après les choix déjà faits. Choisissons des
générateurs
$$
k_1, \ldots, k_t \in
\Ker(P^{\oplus r} \xrightarrow{(b_1, \ldots, b_r)} P)
$$
et écrivons $k_i = (k_{i1}, \ldots, k_{ir})$ et $K = (k_{ij})$, de sorte que
$$
P^{\oplus t} \xrightarrow{K}
P^{\oplus r} \xrightarrow{(b_1, \ldots, b_r)}
P \to R \to 0
$$
soit une suite exacte de $P$-modules. En particulier,
$\sum k_{ij} b_j = 0$. Quitte à augmenter $N$, nous pouvons supposer que
$N - 1$ convient dans le lemme d'Artin--Rees pour les deux premiers
morphismes de cette suite exacte (voir Compléments d'algèbre, section
\ref{more-algebra-section-artin-rees}, pour la terminologie).

\medskip\noindent
D'après l'hypothèse,
$\mathcal{O}_{S, s} = \Lambda_{\Lambda \cap \mathfrak m'}$ est un anneau G.
Par conséquent, d'après Compléments d'algèbre, proposition
\ref{more-algebra-proposition-finite-type-over-G-ring}, l'anneau
$\Lambda[x_1, \ldots, x_s]_{\mathfrak m'}$ est un anneau $G$. D'après
Lissage des homomorphismes d'anneaux, théorème
\ref{smoothing-theorem-approximation-property-variant}, il existe donc un
homomorphisme d'anneaux étale
$$
\Lambda[x_1, \ldots, x_s]_{\mathfrak m'} \to B,
$$
un idéal maximal $\mathfrak m_B$ de $B$ au-dessus de $\mathfrak m'$ et des
éléments $a'_i, b'_i, c'_{ij}, k'_{ij} \in B'$ tels que
\begin{enumerate}
\item $\kappa(\mathfrak m') = \kappa(\mathfrak m_B)$, ce qui implique que
$\Lambda[x_1, \ldots, x_s]_{\mathfrak m'} \subset B_{\mathfrak m_B}
\subset P$, et $P$ s'identifie au complété de $B$ en $\mathfrak m_B$ ; voir
la remarque précédant Lissage des homomorphismes d'anneaux, théorème
\ref{smoothing-theorem-approximation-property-variant},
\item $a_i - a'_i, b_i - b'_i, c_{ij} - c'_{ij}, k_{ij} - k'_{ij} \in
(\mathfrak m')^N P$, et
\item $f_j(a'_1, \ldots, a'_u) = \sum c'_{ji} b'_i$ et
$\sum k'_{ij}b'_j = 0$.
\end{enumerate}
Posons $A = B/(b'_1, \ldots, b'_r)$ et notons $\mathfrak m_A$ l'image de
$\mathfrak m_B$ dans $A$. (Notons que $A$ est essentiellement de type fini
sur $\Lambda$ ; à la fin de la démonstration, nous expliquerons comment
obtenir un $A$ qui soit de type fini sur $\Lambda$.) Il existe un
homomorphisme d'anneaux $C \to A$ défini par $y_i \mapsto a'_i$, car les $a'_i$
satisfont les équations voulues modulo $(b'_1, \ldots, b'_r)$. Notons que
$A/\mathfrak m_A^N = R/\mathfrak m_R^N$ comme quotients de $P = B^\wedge$,
d'après la propriété (2) ci-dessus. Posons $x_A = x_C|_{\Spec(A)}$. Puisque
les homomorphismes
$$
C \to A \to A/\mathfrak m_A^N \cong R/\mathfrak m_R^N
\quad\text{et}\quad
C \to R \to R/\mathfrak m_R^N
$$
sont égaux, nous voyons que $x_A$ et $x$ coïncident modulo
$\mathfrak m_R^N$ par l'isomorphisme
$A/\mathfrak m_A^N = R/\mathfrak m_R^N$. Nous avons ainsi démontré les
propriétés (1) -- (5) de l'énoncé du lemme. Pour démontrer (6), notons que
$$
P^{\oplus t} \xrightarrow{K}
P^{\oplus r} \xrightarrow{(b_1, \ldots, b_r)}
P
\quad\text{et}\quad
P^{\oplus t} \xrightarrow{K'}
P^{\oplus r} \xrightarrow{(b'_1, \ldots, b'_r)}
P
$$
sont deux complexes de $P$-modules congrus modulo $(\mathfrak m')^N$, le
premier étant exact. D'après notre choix de $N$ ci-dessus et Compléments
d'algèbre, lemme \ref{more-algebra-lemma-approximate-complex-graded}, nous
voyons que $R = P/(b_1, \ldots, b_r)$ et
$P/(b'_1, \ldots, b'_r) = B^\wedge/(b'_1, \ldots, b'_r) = A^\wedge$
ont des algèbres graduées associées isomorphes, ce qu'il fallait démontrer.

\medskip\noindent
Ce dernier paragraphe de la démonstration sert à résoudre la difficulté que
$A$ est essentiellement de type fini sur $S$, mais pas encore de type fini.
La construction ci-dessus donne $A = B/(b'_1, \ldots, b'_r)$ et
$\mathfrak m_A \subset A$, avec $B$ étale sur
$\Lambda[x_1, \ldots, x_s]_{\mathfrak m'}$. Ainsi, $A$ est de type fini sur
l'anneau noethérien $\Lambda[x_1, \ldots, x_s]_{\mathfrak m'}$. Nous pouvons
donc écrire $A = (A_0)_{\mathfrak m'}$ pour une
$\Lambda[x_1, \ldots, x_n]$-algèbre de type fini $A_0$. Alors
$A = \colim (A_0)_f$, où
$f \in \Lambda[x_1, \ldots, x_n] \setminus \mathfrak m'$ ; voir Algèbre,
lemme \ref{algebra-lemma-localization-colimit}. Comme
$p : \mathcal{X} \to (\Sch/S)_{fppf}$ commute aux limites sur les objets,
nous voyons que $x_A$ provient d'un objet $x_{(A_0)_f}$ sur
$\Spec((A_0)_f)$ pour un $f$ comme ci-dessus. Après avoir remplacé $A$ par
$(A_0)_f$, $x_A$ par $x_{(A_0)_f}$ et $\mathfrak m_A$ par
$(A_0)_f \cap \mathfrak m_A$, la démonstration est achevée.
\end{proof}





\section{Commutation aux limites}
\label{section-limits}

\noindent
Le morphisme $p : \mathcal{X} \to (\Sch/S)_{fppf}$ commute aux limites sur les
objets, au sens de Critères de représentabilité, section
\ref{criteria-section-limit-preserving}, si le foncteur de la définition
ci-dessous est essentiellement surjectif. Toutefois, l'exemple d'Exemples,
section \ref{examples-section-limit-preserving}, montre que cela n'équivaut
pas à commuter aux limites.

\begin{definition}
\label{definition-limit-preserving}
Soit $S$ un schéma. Soit $\mathcal{X}$ une catégorie fibrée en groupoïdes sur
$(\Sch/S)_{fppf}$. Nous disons que $\mathcal{X}$ {\it commute aux limites} si,
pour tout schéma affine $T$ sur $S$ qui est une limite $T = \lim T_i$ d'un
système projectif filtrant de schémas affines $T_i$ sur $S$, nous avons une
équivalence
$$
\colim \mathcal{X}_{T_i} \longrightarrow \mathcal{X}_T
$$
de catégories-fibres.
\end{definition}

\noindent
Explicitons ce que cela signifie. Premièrement, pour des objets $x, y$ de
$\mathcal{X}$ sur $T_i$, nous devons avoir
$$
\Mor_{\mathcal{X}_T}(x|_T, y|_T) =
\colim_{i' \geq i} \Mor_{\mathcal{X}_{T_{i'}}}(x|_{T_{i'}}, y|_{T_{i'}})
$$
et, deuxièmement, tout objet de $\mathcal{X}_T$ doit être isomorphe à la
restriction d'un objet sur $T_i$ pour un certain $i$. Notons que la première
condition signifie que les préfaisceaux $\mathit{Isom}_\mathcal{X}(x, y)$
(voir Champs, définition \ref{stacks-definition-mor-presheaf}) commutent aux
limites.

\begin{lemma}
\label{lemma-fibre-product-limit-preserving}
Soit $S$ un schéma. Soient $p : \mathcal{X} \to \mathcal{Y}$ et
$q : \mathcal{Z} \to \mathcal{Y}$ des $1$-morphismes de catégories fibrées
en groupoïdes sur $(\Sch/S)_{fppf}$.
\begin{enumerate}
\item Si $\mathcal{X} \to (\Sch/S)_{fppf}$ et
$\mathcal{Z} \to (\Sch/S)_{fppf}$ commutent aux limites sur les objets et si
$\mathcal{Y}$ commute aux limites, alors
$\mathcal{X} \times_\mathcal{Y} \mathcal{Z} \to (\Sch/S)_{fppf}$ commute aux
limites sur les objets.
\item Si $\mathcal{X}$, $\mathcal{Y}$ et $\mathcal{Z}$ commutent aux limites,
alors $\mathcal{X} \times_\mathcal{Y} \mathcal{Z}$ commute aussi aux limites.
\end{enumerate}
\end{lemma}

\begin{proof}
C'est formel. Démontrons (1). Soit $T = \lim_{i \in I} T_i$ la limite
filtrante de schémas affines $T_i$ sur $S$. Nous allons montrer que le
foncteur $\colim \mathcal{X}_{T_i} \to \mathcal{X}_T$ est essentiellement
surjectif. Rappelons qu'un objet du produit fibré au-dessus de $T$ est un
quadruplet $(T, x, z, \alpha)$, où $x$ est un objet de $\mathcal{X}$ situé
au-dessus de $T$, où $z$ est un objet de $\mathcal{Z}$ situé au-dessus de $T$
et où $\alpha : p(x) \to q(z)$ est un morphisme de la catégorie-fibre de
$\mathcal{Y}$ au-dessus de $T$. D'après les hypothèses sur $\mathcal{X}$ et
$\mathcal{Z}$, il existe un $i$ et des objets $x_i$ et $z_i$ sur $T_i$ tels
que $x_i|_T \cong T$ et $z_i|_T \cong z$. Alors $\alpha$ correspond à un
isomorphisme $p(x_i)|_T \to q(z_i)|_T$ qui provient d'un isomorphisme
$\alpha_{i'} : p(x_i)|_{T_{i'}} \to q(z_i)|_{T_{i'}}$, d'après notre hypothèse
sur $\mathcal{Y}$. Après avoir remplacé
$i$ par $i'$, $x_i$ par $x_i|_{T_{i'}}$ et $z_i$ par $z_i|_{T_{i'}}$, nous
voyons que $(T_i, x_i, z_i, \alpha_i)$ est un objet du produit fibré au-dessus
de $T_i$ dont la restriction est un objet isomorphe à
$(T, x, z, \alpha)$ sur $T$, comme voulu.

\medskip\noindent
Nous omettons les arguments montrant que
$\colim \mathcal{X}_{T_i} \to \mathcal{X}_T$ est pleinement fidèle dans (2).
\end{proof}

\begin{lemma}
\label{lemma-limit-preserving-algebraic-space}
Soit $S$ un schéma. Soit $\mathcal{X}$ un champ algébrique sur $S$.
Les conditions suivantes sont équivalentes :
\begin{enumerate}
\item $\mathcal{X}$ est un champ en ensemblöïdes et
$\mathcal{X} \to (\Sch/S)_{fppf}$ commute aux limites sur les objets ;
\item $\mathcal{X}$ est un champ en ensemblöïdes et commute aux limites ;
\item $\mathcal{X}$ est représentable par un espace algébrique localement de
présentation finie.
\end{enumerate}
\end{lemma}

\begin{proof}
Sous chacune des trois hypothèses, $\mathcal{X}$ est représentable par un
espace algébrique $X$ sur $S$ ; voir Champs algébriques, proposition
\ref{algebraic-proposition-algebraic-stack-no-automorphisms}. Il est clair que
(1) et (2) sont équivalentes, car un foncteur entre ensemblöïdes est une
équivalence si et seulement s'il est surjectif sur les classes d'isomorphisme.
Enfin, (1) et (3) sont équivalentes d'après Limites d'espaces, proposition
\ref{spaces-limits-proposition-characterize-locally-finite-presentation}.
\end{proof}

\begin{lemma}
\label{lemma-diagonal}
Soit $S$ un schéma. Soit $\mathcal{X}$ une catégorie fibrée en groupoïdes sur
$(\Sch/S)_{fppf}$. Supposons que
$\Delta : \mathcal{X} \to \mathcal{X} \times \mathcal{X}$ soit représentable
par des espaces algébriques et que $\mathcal{X}$ commute aux limites. Alors
$\Delta$ est localement de type fini.
\end{lemma}

\begin{proof}
Appliquons Critères de représentabilité, lemme
\ref{criteria-lemma-check-property-limit-preserving}. Soit $V$ un schéma
affine $V$ localement de présentation finie sur $S$, et soit $\theta$ un objet
de $\mathcal{X} \times \mathcal{X}$ sur $V$. Soit $F_\theta$ un espace
algébrique représentant
$\mathcal{X} \times_{\Delta, \mathcal{X} \times \mathcal{X}, \theta}
(\Sch/V)_{fppf}$, et soit $f_\theta : F_\theta \to V$ le morphisme canonique
(voir Champs algébriques, section
\ref{algebraic-section-morphisms-representable-by-algebraic-spaces}). Il
suffit de montrer que $F_\theta \to V$ possède les propriétés correspondantes.
D'après les lemmes \ref{lemma-fibre-product-limit-preserving} et
\ref{lemma-limit-preserving-algebraic-space}, le morphisme
$F_\theta \to S$ est localement de présentation finie. Il en résulte que
$F_\theta \to V$ est localement de type fini d'après Morphismes d'espaces,
lemme \ref{spaces-morphisms-lemma-permanence-finite-type}.
\end{proof}






\section{Versalité}
\label{section-versality}

\noindent
Dans la section précédente, nous avons expliqué comment approcher des objets
sur des anneaux locaux complets par des objets algébriques. Mais, pour montrer
qu'un champ $\mathcal{X}$ est un champ algébrique, nous devons trouver des
$1$-morphismes lisses de schémas vers $\mathcal{X}$. Comme nous ne supposerons
pas a priori que $\mathcal{X}$ possède une diagonale représentable, nous ne
pouvons même pas parler de morphismes lisses vers $\mathcal{X}$. Empruntant
plutôt la terminologie de la théorie des déformations, nous allons introduire
les objets versels.

\begin{definition}
\label{definition-versal-formal-object}
Soit $S$ un schéma localement noethérien. Soit
$p : \mathcal{X} \to (\Sch/S)_{fppf}$ une catégorie fibrée en groupoïdes.
Soit $\xi = (R, \xi_n, f_n)$ un objet formel. Posons $k = R/\mathfrak m$ et
$x_0 = \xi_1$. Nous dirons que $\xi$ est {\it versel} si $\xi$, considéré
comme objet formel de $\mathcal{F}_{\mathcal{X}, k, x_0}$ (remarque
\ref{remark-formal-objects-match}), est versel au sens de Théorie formelle des
déformations, définition \ref{formal-defos-definition-versal}.
\end{definition}

\noindent
Explicitons brièvement cette condition. Avec les notations de la définition,
supposons donnés des morphismes $\xi_1 = x_0 \to y \to z$ de $\mathcal{X}$
situés au-dessus d'immersions fermées
$\Spec(k) \to \Spec(A) \to \Spec(B)$
où $A, B$ sont des anneaux locaux artiniens de corps résiduel $k$.
Supposons donnés un $n \geq 1$ et un diagramme commutatif
$$
\vcenter{
\xymatrix{
& y \ar[ld] \\
\xi_n & \xi_1 \ar[u] \ar[l]
}
}
\quad\text{au-dessus de}\quad
\vcenter{
\xymatrix{
& \Spec(A) \ar[ld] \\
\Spec(R/\mathfrak m^n) & \Spec(k) \ar[u] \ar[l]
}
}
$$
La versalité signifie que, pour toutes données comme ci-dessus, il existe un
$m \geq n$ et un diagramme commutatif
$$
\vcenter{
\xymatrix{
& & z \ar[lldd] \\
& & y \ar[ld] \ar[u] \\
\xi_m & \xi_n \ar[l] & \xi_1 \ar[u] \ar[l]
}
}
\quad\text{au-dessus de}
\vcenter{
\xymatrix{
& & \Spec(B) \ar[lldd] \\
& & \Spec(A) \ar[ld] \ar[u] \\
\Spec(R/\mathfrak m^m) &
\Spec(R/\mathfrak m^n) \ar[l] &
\Spec(k) \ar[u] \ar[l]
}
}
$$
Comparer avec Théorie formelle des déformations, remarque
\ref{formal-defos-remark-versal-object}.

\medskip\noindent
Soit $S$ un schéma localement noethérien. Soit $U$ un schéma sur $S$ dont le
morphisme structural $U \to S$ est localement de type fini. Soit
$u_0 \in U$ un point de type fini de $U$ ; voir Morphismes, définition
\ref{morphisms-definition-finite-type-point}. Posons $k = \kappa(u_0)$.
Notons que la composée $\Spec(k) \to S$ est elle aussi de type fini ; voir
Morphismes, lemme \ref{morphisms-lemma-composition-finite-type}. Soit
$p : \mathcal{X} \to (\Sch/S)_{fppf}$ une catégorie fibrée en groupoïdes.
Soit $x$ un objet de $\mathcal{X}$ situé au-dessus de $U$. Notons $x_0$
l'image réciproque de $x$ par $u_0$. D'après le $2$-lemme de Yoneda, $x$
correspond à un $1$-morphisme
$$
x : (\Sch/U)_{fppf} \longrightarrow \mathcal{X},
$$
voir Champs algébriques, section \ref{algebraic-section-2-yoneda}. Nous
obtenons un morphisme de catégories de prédéformations
\begin{equation}
\label{equation-hat-x}
\hat x :
\mathcal{F}_{(\Sch/U)_{fppf}, k, u_0}
\longrightarrow
\mathcal{F}_{\mathcal{X}, k, x_0},
\end{equation}
au-dessus de $\mathcal{C}_\Lambda$ ; voir (\ref{equation-functoriality}).

\begin{definition}
\label{definition-versal}
Soit $S$ un schéma localement noethérien.
Soit $\mathcal{X}$ une catégorie fibrée en groupoïdes sur $(\Sch/S)_{fppf}$.
Soit $U$ un schéma localement de type fini sur $S$.
Soit $x$ un objet de $\mathcal{X}$ situé au-dessus de $U$.
Soit $u_0$ un point de type fini de $U$.
Nous disons que $x$ est {\it versel} en $u_0$ si le morphisme $\hat x$
(\ref{equation-hat-x}) est lisse ; voir Théorie formelle des déformations,
définition \ref{formal-defos-definition-smooth-morphism}.
\end{definition}

\noindent
Cette définition concorde avec notre notion de versalité pour les objets
formels de $\mathcal{X}$.

\begin{lemma}
\label{lemma-versality-matches}
Reprenons les notations de la définition \ref{definition-versal}.
Soit $R = \mathcal{O}_{U, u_0}^\wedge$. Soit $\xi$ l'objet formel de
$\mathcal{X}$ sur $R$ associé à $x|_{\Spec(R)}$ ; voir
(\ref{equation-approximation}). Alors
$$
x\text{ est versel en }u_0
\Leftrightarrow
\xi\text{ est versel}
$$
\end{lemma}

\begin{proof}
Notons que $\mathcal{O}_{U, u_0}$ est une $S$-algèbre locale noethérienne de
corps résiduel $k$. Ainsi, $R = \mathcal{O}_{U, u_0}^\wedge$ est un objet de
$\mathcal{C}_\Lambda^\wedge$ ; voir Théorie formelle des déformations,
définition \ref{formal-defos-definition-completion-CLambda}.
Rappelons que $\xi$ est versel si
$\underline{\xi} : \underline{R}|_{\mathcal{C}_\Lambda} \to
\mathcal{F}_{\mathcal{X}, k, x_0}$
est lisse, et que $x$ est versel en $u_0$ si
$\hat x : \mathcal{F}_{(\Sch/U)_{fppf}, k, u_0}
\to \mathcal{F}_{\mathcal{X}, k, x_0}$ est lisse. Il existe une
identification de catégories de prédéformations
$$
\underline{R}|_{\mathcal{C}_\Lambda}
=
\mathcal{F}_{(\Sch/U)_{fppf}, k, u_0},
$$
voir Théorie formelle des déformations, remarque
\ref{formal-defos-remark-formal-objects-yoneda}, pour les notations. En effet,
pour une $S$-algèbre locale artinienne $A$ dont le corps résiduel est identifié
à $k$, nous avons
$$
\Mor_{\mathcal{C}_\Lambda^\wedge}(R, A) =
\{\varphi \in \Mor_S(\Spec(A), U) \mid \varphi|_{\Spec(k)} = u_0\}
$$
En développant les définitions, le lecteur vérifie que l'application obtenue
$$
\underline{R}|_{\mathcal{C}_\Lambda} =
\mathcal{F}_{(\Sch/U)_{fppf}, k, u_0}
\xrightarrow{\hat x}
\mathcal{F}_{\mathcal{X}, k, x_0},
$$
est égale à $\underline{\xi}$, ce qui démontre le lemme.
\end{proof}

\noindent
Voici une vérification de cohérence.

\begin{lemma}
\label{lemma-versal-implies-smooth}
Soit $S$ un schéma localement noethérien. Soit $f : U \to V$ un morphisme de
schémas localement de type fini sur $S$. Soit $u_0 \in U$ un point de type
fini. Les conditions suivantes sont équivalentes :
\begin{enumerate}
\item $f$ est lisse en $u_0$ ;
\item $f$, considéré comme un objet de $(\Sch/V)_{fppf}$ sur $U$, est versel
en $u_0$.
\end{enumerate}
\end{lemma}

\begin{proof}
C'est une reformulation de Compléments sur les morphismes, lemme
\ref{more-morphisms-lemma-lifting-along-artinian-at-point}.
\end{proof}

\noindent
Cette notion se comporte bien vis-à-vis des extensions de corps.

\begin{lemma}
\label{lemma-versal-change-of-field}
Soient $S$, $\mathcal{X}$, $U$, $x$, $u_0$ comme dans la définition
\ref{definition-versal}. Soit $l$ un corps et soit
$u_{l, 0} : \Spec(l) \to U$ un morphisme d'image $u_0$ tel que
$l/k = \kappa(u_0)$ soit fini. Posons $x_{l, 0} = x_0|_{\Spec(l)}$.
Si $\mathcal{X}$ satisfait (RS) et si $x$ est versel en $u_0$, alors
$$
\mathcal{F}_{(\Sch/U)_{fppf}, l, u_{l, 0}}
\longrightarrow
\mathcal{F}_{\mathcal{X}, l, x_{l, 0}}
$$
est lisse.
\end{lemma}

\begin{proof}
Notons que $(\Sch/U)_{fppf}$ satisfait (RS) d'après le lemme
\ref{lemma-algebraic-stack-RS}. Le foncteur du lemme est donc le foncteur
$$
(\mathcal{F}_{(\Sch/U)_{fppf}, k , u_0})_{l/k}
\longrightarrow
(\mathcal{F}_{\mathcal{X}, k , x_0})_{l/k}
$$
associé à $\hat x$ ; voir le lemme \ref{lemma-change-of-field}. Le résultat
découle donc de Théorie formelle des déformations, lemme
\ref{formal-defos-lemma-change-of-fields-smooth}.
\end{proof}

\noindent
Le lemme suivant est une autre vérification de cohérence. Il signifie à peu
près que, si $x$ est versel en $u_0$ comme dans la définition
\ref{definition-versal}, alors $x$, considéré comme un morphisme de $U$ vers
$\mathcal{X}$, est lisse après tout changement de base par un schéma.

\begin{lemma}
\label{lemma-base-change-versal}
Soient $S$, $\mathcal{X}$, $U$, $x$, $u_0$ comme dans la définition
\ref{definition-versal}. Supposons que
\begin{enumerate}
\item $\Delta : \mathcal{X} \to \mathcal{X} \times \mathcal{X}$
soit représentable par des espaces algébriques ;
\item $\Delta$ soit localement de type fini (par exemple si $\mathcal{X}$
commute aux limites) ;
\item $\mathcal{X}$ satisfasse (RS).
\end{enumerate}
Soit $V$ un schéma localement de type fini sur $S$ et soit $y$ un objet de
$\mathcal{X}$ sur $V$. Formons le $2$-produit fibré
$$
\xymatrix{
\mathcal{Z} \ar[r] \ar[d] & (\Sch/U)_{fppf} \ar[d]^x \\
(\Sch/V)_{fppf} \ar[r]^y & \mathcal{X}
}
$$
Soit $Z$ l'espace algébrique représentant $\mathcal{Z}$ et soit
$z_0 \in |Z|$ un point de type fini situé au-dessus de $u_0$. Si $x$ est
versel en $u_0$, alors le morphisme $Z \to V$ est lisse en $z_0$.
\end{lemma}

\begin{proof}
(La remarque entre parenthèses dans l'énoncé résulte du lemme
\ref{lemma-diagonal}.) Notons que $Z$ existe d'après l'hypothèse (1) et Champs
algébriques, lemme \ref{algebraic-lemma-representable-diagonal}. D'après
l'hypothèse (2), $Z \to V \times_S U$ est localement de type fini.
Choisissons un schéma $W$, un point fermé $w_0 \in W$ et un morphisme étale
$W \to Z$ envoyant $w_0$ sur $z_0$ ; voir Morphismes d'espaces, définition
\ref{spaces-morphisms-definition-finite-type-point}. Alors $W$ est localement
de type fini sur $S$ et $w_0$ est un point de type fini de $W$. Posons
$l = \kappa(z_0)$. Notons $z_{l, 0}$, $v_{l, 0}$, $u_{l, 0}$ et $x_{l, 0}$
les objets de $\mathcal{Z}$, $(\Sch/V)_{fppf}$, $(\Sch/U)_{fppf}$ et
$\mathcal{X}$ sur $\Spec(l)$ obtenus par image réciproque à
$\Spec(l) = w_0$. Considérons
$$
\xymatrix{
\mathcal{F}_{(\Sch/W)_{fppf}, l, w_0} \ar[r] &
\mathcal{F}_{\mathcal{Z}, l, z_{l, 0}} \ar[d] \ar[r] &
\mathcal{F}_{(\Sch/U)_{fppf}, l, u_{l, 0}} \ar[d] \\
& \mathcal{F}_{(\Sch/V)_{fppf}, l, v_{l, 0}} \ar[r] &
\mathcal{F}_{\mathcal{X}, l, x_{l, 0}}
}
$$
D'après le lemme \ref{lemma-fibre-product-deformation-categories}, le carré
est un produit fibré de catégories de prédéformations. D'après le lemme
\ref{lemma-versal-change-of-field}, la flèche verticale de droite est lisse.
D'après Théorie formelle des déformations, lemme
\ref{formal-defos-lemma-smooth-properties}, la flèche verticale de gauche est
lisse. D'après le lemme
\ref{lemma-formally-smooth-on-deformation-categories}, la flèche horizontale
de gauche est lisse. Nous en concluons que l'application
$$
\mathcal{F}_{(\Sch/W)_{fppf}, l, w_0} \to
\mathcal{F}_{(\Sch/V)_{fppf}, l, v_{l, 0}}
$$
est lisse d'après Théorie formelle des déformations, lemme
\ref{formal-defos-lemma-smooth-properties}. Nous en concluons que
$W \to V$ est lisse en $w_0$ d'après Compléments sur les morphismes, lemme
\ref{more-morphisms-lemma-lifting-along-artinian-at-point}. Cela signifie
exactement que $Z \to V$ est lisse en $z_0$, et la démonstration est achevée.
\end{proof}

\noindent
Reformulons le résultat d'approximation en termes d'objets versels.

\begin{lemma}
\label{lemma-approximate-versal}
Soit $S$ un schéma localement noethérien. Soit
$p : \mathcal{X} \to (\Sch/S)_{fppf}$ une catégorie fibrée en groupoïdes.
Soit $\xi = (R, \xi_n, f_n)$ un objet formel de $\mathcal{X}$ tel que
$\xi_1$ soit situé au-dessus de $\Spec(k) \to S$, d'image $s \in S$.
Supposons que
\begin{enumerate}
\item $\xi$ soit versel ;
\item $\xi$ soit effectif ;
\item $\mathcal{O}_{S, s}$ soit un anneau G ;
\item $p : \mathcal{X} \to (\Sch/S)_{fppf}$ commute aux limites sur les
objets.
\end{enumerate}
Alors il existe un morphisme de type fini $U \to S$, un point de type fini
$u_0 \in U$ de corps résiduel $k$ et un objet $x$ de $\mathcal{X}$ sur $U$,
tels que $x$ soit versel en $u_0$ et que
$x|_{\Spec(\mathcal{O}_{U, u_0}/\mathfrak m_{u_0}^n)} \cong \xi_n$.
\end{lemma}

\begin{proof}
Choisissons un objet $x_R$ de $\mathcal{X}$ situé au-dessus de $\Spec(R)$ dont
l'objet formel associé est $\xi$. Posons $N = 2$ et appliquons le lemme
\ref{lemma-approximate}. Nous obtenons $A, \mathfrak m_A, x_A, \ldots$.
Soit $\eta = (A^\wedge, \eta_n, g_n)$ l'objet formel associé à
$x_A|_{\Spec(A^\wedge)}$. Nous avons un diagramme
$$
\vcenter{
\xymatrix{
& \eta \ar[d] \\
\xi \ar[r] \ar@{..>}[ru] & \xi_2 = \eta_2
}
}
\quad\text{au-dessus de}\quad
\vcenter{
\xymatrix{
& A^\wedge \ar[d] \\
R \ar[r] \ar@{..>}[ru] & R/\mathfrak m_R^2 = A/\mathfrak m_A^2
}
}
$$
La versalité de $\xi$ signifie exactement que nous pouvons trouver les
flèches pointillées dans les diagrammes, car nous pouvons trouver
successivement des morphismes $\xi \to \eta_3$, $\xi \to \eta_4$, et ainsi de
suite, d'après Théorie formelle des déformations, remarque
\ref{formal-defos-remark-versal-object}. L'homomorphisme d'anneaux correspondant
$R \to A^\wedge$ est surjectif d'après Théorie formelle des déformations,
lemme \ref{formal-defos-lemma-surjective-cotangent-space}. D'autre part,
nous avons
$\dim_k \mathfrak m_R^n/\mathfrak m_R^{n + 1} =
\dim_k \mathfrak m_A^n/\mathfrak m_A^{n + 1}$ pour tout $n$, par construction.
Ainsi, $R/\mathfrak m_R^n$ et $A/\mathfrak m_A^n$ ont la même longueur
(finie) comme $\Lambda$-modules, par additivité de la longueur et d'après
Théorie formelle des déformations, lemme \ref{formal-defos-lemma-length}.
Il s'ensuit que $R/\mathfrak m_R^n \to A/\mathfrak m_A^n$ est un isomorphisme
pour tout $n$, donc que $R \to A^\wedge$ est un isomorphisme. Ainsi, $\eta$
est isomorphe à un objet versel et est donc lui-même versel. D'après le lemme
\ref{lemma-versality-matches}, nous concluons que $x_A$ est versel au point
$u_0$ de $U = \Spec(A)$ correspondant à $\mathfrak m_A$.
\end{proof}

\begin{example}
\label{example-approximate-versal-implies}
Dans cet exemple, nous montrons que l'anneau local $\mathcal{O}_{S, s}$ doit
être un anneau G pour que le résultat du lemme
\ref{lemma-approximate-versal} soit vrai. Soit en effet $\Lambda$ un anneau
noethérien et soit $\mathfrak m$ un idéal maximal de $\Lambda$. Posons
$R = \Lambda_\mathfrak m^\wedge$. Soit $\Lambda \to C \to R$ une
factorisation avec $C$ de type fini sur $\Lambda$. Posons
$S = \Spec(\Lambda)$, $U = S \setminus \{\mathfrak m\}$ et
$S' = U \amalg \Spec(C)$. Considérons le foncteur
$F : (\Sch/S)_{fppf}^{opp} \to \textit{Sets}$ défini par la règle
$$
F(T) = 
\left\{
\begin{matrix}
* & \text{si }T \to S\text{ se factorise par }S' \\
\emptyset & \text{sinon}
\end{matrix}
\right.
$$
Soit $\mathcal{X} = \mathcal{S}_F$ la catégorie fibrée en ensembles associée
à $F$ ; voir Champs algébriques, section \ref{algebraic-section-split}.
Alors $\mathcal{X} \to (\Sch/S)_{fppf}$ commute aux limites sur les objets et
il existe un objet formel effectif et versel $\xi$ sur $R$. Ainsi, si la
conclusion du lemme \ref{lemma-approximate-versal} vaut pour $\mathcal{X}$,
il existe un homomorphisme d'anneaux de type fini $\Lambda \to A$ et un idéal
maximal $\mathfrak m_A$ au-dessus de $\mathfrak m$ tels que
\begin{enumerate}
\item $\kappa(\mathfrak m) = \kappa(\mathfrak m_A)$ ;
\item $\Lambda \to A$ et $\mathfrak m_A$ satisfont la condition (4)
d'Algèbre, lemme \ref{algebra-lemma-smooth-test-artinian} ;
\item il existe un homomorphisme de $\Lambda$-algèbres $C \to A$.
\end{enumerate}
Ainsi, $\Lambda \to A$ est lisse en $\mathfrak m_A$ d'après le lemme cité.
En prenant une tranche de $A$, nous pouvons supposer que $\Lambda \to A$ est
étale en $\mathfrak m_A$ ; voir par exemple Compléments sur les morphismes,
lemme \ref{more-morphisms-lemma-slice-smooth}, ou raisonner directement.
Écrivons $C = \Lambda[y_1, \ldots, y_n]/(f_1, \ldots, f_m)$. Alors
$C \to R$ correspond à une solution dans $R$ du système d'équations
$f_1 = \ldots = f_m = 0$ ; voir Lissage des homomorphismes d'anneaux, section
\ref{smoothing-section-approximation-G-rings}. Ainsi, si la conclusion du
lemme \ref{lemma-approximate-versal} vaut pour tout $\mathcal{X}$ comme
ci-dessus, tout système d'équations qui admet une solution dans $R$ en admet
une dans l'hensélisé de $\Lambda_{\mathfrak m}$. Autrement dit, la propriété
d'approximation vaut pour $\Lambda_{\mathfrak m}^h$. Cela implique que
$\Lambda_{\mathfrak m}^h$ est un anneau G (insérer ici un renvoi futur ; voir
aussi la discussion dans Lissage des homomorphismes d'anneaux, section
\ref{smoothing-section-introduction}), ce qui implique à son tour que
$\Lambda_{\mathfrak m}$ est un anneau G.
\end{example}






\section{Ouverture de la versalité}
\label{section-openness-versality}

\noindent
Nous en venons maintenant à l'ouverture de la versalité.

\begin{definition}
\label{definition-openness-versality}
Soit $S$ un schéma localement noethérien.
\begin{enumerate}
\item Soit $\mathcal{X}$ une catégorie fibrée en groupoïdes sur
$(\Sch/S)_{fppf}$. Nous disons que $\mathcal{X}$ satisfait l'{\it ouverture
de la versalité} si, pour tout schéma $U$ localement de type fini sur $S$,
tout objet $x$ de $\mathcal{X}$ sur $U$ et tout point de type fini
$u_0 \in U$ tel que $x$ soit versel en $u_0$, il existe un voisinage ouvert
$u_0 \in U' \subset U$ tel que $x$ soit versel en tout point de type fini de
$U'$.
\item Soit $f : \mathcal{X} \to \mathcal{Y}$ un $1$-morphisme de catégories
fibrées en groupoïdes sur $(\Sch/S)_{fppf}$. Nous disons que $f$ satisfait
l'{\it ouverture de la versalité} si, pour tout schéma $U$ localement de type
fini sur $S$ et tout objet $y$ de $\mathcal{Y}$ sur $U$, l'ouverture de la
versalité vaut pour
$(\Sch/U)_{fppf} \times_\mathcal{Y} \mathcal{X}$.
\end{enumerate}
\end{definition}

\noindent
L'ouverture de la versalité est souvent la condition la plus difficile à
vérifier. L'exemple suivant montre néanmoins qu'il est nécessaire de l'exiger.

\begin{example}
\label{example-versality}
Soit $k$ un corps et posons $\Lambda = k[s, t]$. Considérons le foncteur
$F : \Lambda\text{-algèbres} \longrightarrow \textit{Sets}$
défini par la règle
$$
F(A) =
\left\{
\begin{matrix}
* & \text{s'il existe }f_1, \ldots, f_n \in A\text{ tels que } \\
  & A = (s, t, f_1, \ldots, f_n)\text{ et } f_i s = 0\ \forall i \\
\emptyset & \text{sinon}
\end{matrix}
\right.
$$
Géométriquement, $F(A) = *$ signifie qu'il existe un voisinage ouvert
quasi-compact $W$ de $V(s, t) \subset \Spec(A)$ tel que $s|_W = 0$.
Soit $\mathcal{X} \subset (\Sch/\Spec(\Lambda))_{fppf}$ la sous-catégorie
pleine formée des schémas $T$ qui possèdent un recouvrement ouvert affine
$T = \bigcup \Spec(A_j)$ tel que $F(A_j) = *$ pour tout $j$. Alors
$\mathcal{X}$ satisfait [0], [1], [2], [3] et [4], mais pas [5]. En effet,
sur $U = \Spec(k[s, t]/(s))$, il existe un objet $x$ qui est versel en
$u_0 = (s, t)$ mais en aucun autre point. Nous omettons les détails.
\end{example}

\noindent
Soit $S$ un schéma localement noethérien. Soit
$f : \mathcal{X} \to \mathcal{Y}$ un $1$-morphisme de catégories fibrées en
groupoïdes sur $(\Sch/S)_{fppf}$. Considérons la propriété suivante :
\begin{equation}
\label{equation-smooth}
\begin{matrix}
\text{pour tout corps }k\text{ de type fini sur }S
\text{ et tout }x_0 \in \Ob(\mathcal{X}_{\Spec(k)})\text{ l'application}\\
\text{application }
\mathcal{F}_{\mathcal{X}, k, x_0} \to \mathcal{F}_{\mathcal{Y}, k, f(x_0)}
\text{ de catégories de prédéformations est lisse}
\end{matrix}
\end{equation}
Formulons quelques lemmes autour de cette notion. Relions-la d'abord à la
versalité (et à son ouverture).

\begin{lemma}
\label{lemma-versal-smooth}
Soit $S$ un schéma localement noethérien. Soit $\mathcal{X}$ une catégorie
fibrée en groupoïdes sur $(\Sch/S)_{fppf}$. Soit $U$ un schéma localement de
type fini sur $S$. Soit $x$ un objet de $\mathcal{X}$ sur $U$. Supposons que
$x$ soit versel en tout point de type fini de $U$ et que $\mathcal{X}$
satisfasse (RS). Alors $x : (\Sch/U)_{fppf} \to \mathcal{X}$ satisfait
(\ref{equation-smooth}).
\end{lemma}

\begin{proof}
Soit $\Spec(l) \to U$ un morphisme, où $l$ est de type fini sur $S$. Alors
l'image $u_0 \in U$ est un point de type fini de $U$ et
$l/\kappa(u_0)$ est une extension finie ; voir la discussion dans Morphismes,
section \ref{morphisms-section-points-finite-type}. Nous voyons donc que
$\mathcal{F}_{(\Sch/U)_{fppf}, l, u_{l, 0}} \to
\mathcal{F}_{\mathcal{X}, l, x_{l, 0}}$
est lisse d'après le lemme \ref{lemma-versal-change-of-field}.
\end{proof}

\begin{lemma}
\label{lemma-composition-smooth}
Soit $S$ un schéma localement noethérien. Soient
$f : \mathcal{X} \to \mathcal{Y}$ et $g : \mathcal{Y} \to \mathcal{Z}$ des
$1$-morphismes composables de catégories fibrées en groupoïdes sur
$(\Sch/S)_{fppf}$. Si $f$ et $g$ satisfont (\ref{equation-smooth}), alors
$g \circ f$ la satisfait aussi.
\end{lemma}

\begin{proof}
Cela résulte formellement de Théorie formelle des déformations, lemme
\ref{formal-defos-lemma-smooth-properties}.
\end{proof}

\begin{lemma}
\label{lemma-base-change-smooth}
Soit $S$ un schéma localement noethérien. Soient
$f : \mathcal{X} \to \mathcal{Y}$ et $\mathcal{Z} \to \mathcal{Y}$ des
$1$-morphismes de catégories fibrées en groupoïdes sur $(\Sch/S)_{fppf}$.
Si $f$ satisfait (\ref{equation-smooth}), alors la projection
$\mathcal{X} \times_\mathcal{Y} \mathcal{Z} \to \mathcal{Z}$ la satisfait aussi.
\end{lemma}

\begin{proof}
Cela résulte immédiatement du lemme
\ref{lemma-fibre-product-deformation-categories} et de Théorie formelle des
déformations, lemme
\ref{formal-defos-lemma-smooth-properties}.
\end{proof}

\begin{lemma}
\label{lemma-smooth-smooth}
Soit $S$ un schéma localement noethérien. Soit
$f : \mathcal{X} \to \mathcal{Y}$ un $1$-morphisme de catégories fibrées en
groupoïdes sur $(\Sch/S)_{fppf}$. Si $f$ est formellement lisse sur les objets,
alors $f$ satisfait (\ref{equation-smooth}). Si $f$ est représentable par des
espaces algébriques et lisse, alors $f$ satisfait (\ref{equation-smooth}).
\end{lemma}

\begin{proof}
Il s'agit d'une reformulation du lemme
\ref{lemma-formally-smooth-on-deformation-categories}.
\end{proof}

\begin{lemma}
\label{lemma-implies-smooth}
Soit $S$ un schéma localement noethérien. Soit
$f : \mathcal{X} \to \mathcal{Y}$ un $1$-morphisme de catégories fibrées en
groupoïdes sur $(\Sch/S)_{fppf}$. Supposons que
\begin{enumerate}
\item $f$ soit représentable par des espaces algébriques ;
\item $f$ satisfasse (\ref{equation-smooth}) ;
\item $\mathcal{X} \to (\Sch/S)_{fppf}$ commute aux limites sur les objets ;
\item $\mathcal{Y}$ commute aux limites.
\end{enumerate}
Alors $f$ est lisse.
\end{lemma}

\begin{proof}
L'ingrédient essentiel de la démonstration est Compléments sur les morphismes,
lemme \ref{more-morphisms-lemma-lifting-along-artinian-at-point}, qui affirme
(presque) qu'un morphisme de schémas de type fini sur $S$ satisfaisant
(\ref{equation-smooth}) est un morphisme lisse. Les autres arguments de la
démonstration relèvent essentiellement de la tenue des comptes.

\medskip\noindent
Soit $V$ un schéma sur $S$ et soit $y$ un objet de $\mathcal{Y}$ sur $V$.
Soit $Z$ un espace algébrique représentant le $2$-produit fibré
$\mathcal{Z} = \mathcal{X} \times_{f, \mathcal{X}, y} (\Sch/V)_{fppf}$.
Nous devons montrer que le morphisme de projection $Z \to V$ est lisse ; voir
Champs algébriques, définition
\ref{algebraic-definition-relative-representable-property}. En fait, il
suffit de le faire lorsque $V$ est un schéma affine localement de présentation
finie sur $S$ ; voir Critères de représentabilité, lemme
\ref{criteria-lemma-check-property-limit-preserving}. Alors
$(\Sch/V)_{fppf}$ commute aux limites d'après le lemme
\ref{lemma-limit-preserving-algebraic-space}. Par conséquent, $Z \to S$ est
localement de présentation finie d'après les lemmes
\ref{lemma-fibre-product-limit-preserving} et
\ref{lemma-limit-preserving-algebraic-space}. Choisissons un schéma $W$ et un
morphisme étale surjectif $W \to Z$. Alors $W$ est localement de présentation
finie sur $S$.

\medskip\noindent
Puisque $f$ satisfait (\ref{equation-smooth}), il en est de même de
$\mathcal{Z} \to (\Sch/V)_{fppf}$ ; voir le lemme
\ref{lemma-base-change-smooth}. Ensuite,
$(\Sch/W)_{fppf} \to \mathcal{Z}$ satisfait (\ref{equation-smooth}) d'après le
lemme \ref{lemma-smooth-smooth}. Ainsi, la composée
$$
(\Sch/W)_{fppf} \to \mathcal{Z} \to (\Sch/V)_{fppf}
$$
satisfait (\ref{equation-smooth}) d'après le lemme
\ref{lemma-composition-smooth}. Compléments sur les morphismes, lemme
\ref{more-morphisms-lemma-lifting-along-artinian-at-point}, montre que la
composée $W \to Z \to V$ est lisse en tout point de type fini $w_0$ de $W$.
Comme le lieu de lissité est ouvert, nous concluons que $W \to V$ est un
morphisme lisse de schémas d'après Morphismes, lemme
\ref{morphisms-lemma-enough-finite-type-points}. Nous concluons donc, par
définition, que $Z \to V$ est un morphisme lisse d'espaces algébriques.
\end{proof}

\noindent
Le lemme ci-dessous décrit l'usage que nous ferons de l'ouverture de la
versalité.

\begin{lemma}
\label{lemma-get-smooth}
Soit $S$ un schéma localement noethérien. Soit
$p : \mathcal{X} \to (\Sch/S)_{fppf}$ une catégorie fibrée en groupoïdes.
Soit $k$ un corps de type fini sur $S$ et soit $x_0$ un objet de
$\mathcal{X}$ sur $\Spec(k)$, d'image $s \in S$. Supposons que
\begin{enumerate}
\item $\Delta : \mathcal{X} \to \mathcal{X} \times \mathcal{X}$ soit
représentable par des espaces algébriques ;
\item $\mathcal{X}$ satisfasse les axiomes [1], [2], [3] (voir la section
\ref{section-axioms}) ;
\item tout objet formel de $\mathcal{X}$ soit effectif ;
\item l'ouverture de la versalité vaille pour $\mathcal{X}$ ;
\item $\mathcal{O}_{S, s}$ soit un anneau G.
\end{enumerate}
Alors il existe un morphisme de type fini $U \to S$ et un objet $x$ de
$\mathcal{X}$ sur $U$ tels que
$$
x : (\Sch/U)_{fppf} \longrightarrow \mathcal{X}
$$
soit lisse et qu'il existe un point de type fini $u_0 \in U$ de corps
résiduel $k$ tel que $x|_{u_0} \cong x_0$.
\end{lemma}

\begin{proof}
D'après l'axiome [2], le lemme \ref{lemma-deformation-category} et la remarque
\ref{remark-deformation-category-implies}, la catégorie
$\mathcal{F}_{\mathcal{X}, k, x_0}$ satisfait (S1) et (S2). Comme son espace
tangent est aussi de dimension finie d'après l'axiome [3], Théorie formelle
des déformations, lemme \ref{formal-defos-lemma-versal-object-existence},
montre que $\mathcal{F}_{\mathcal{X}, k, x_0}$ possède un objet formel versel
$\xi$. L'hypothèse (3) affirme que $\xi$ est effectif. D'après l'axiome [1] et
le lemme \ref{lemma-approximate-versal}, il existe un morphisme de type fini
$U \to S$, un objet $x$ de $\mathcal{X}$ sur $U$ et un point de type fini
$u_0$ de $U$, de corps résiduel $k$, tels que $x$ soit versel en $u_0$ et que
$x|_{\Spec(k)} \cong x_0$. Par l'ouverture de la versalité, nous pouvons
rétrécir $U$ et supposer que $x$ est versel en tout point de type fini de $U$.
Nous affirmons que
$$
x : (\Sch/U)_{fppf} \longrightarrow \mathcal{X}
$$
est lisse, ce qui démontre le lemme. En effet, d'après le lemme
\ref{lemma-versal-smooth}, $x$ satisfait (\ref{equation-smooth}), puis le
lemme \ref{lemma-implies-smooth} achève la démonstration.
\end{proof}






\section{Axiomes}
\label{section-axioms}

\noindent
Soit $S$ un schéma localement noethérien. Soit
$p : \mathcal{X} \to (\Sch/S)_{fppf}$ une catégorie fibrée en groupoïdes.
Voici les axiomes que nous considérerons pour $\mathcal{X}$.
\begin{enumerate}
\item[{[-1]}] une condition ensembliste\footnote{La condition est la suivante :
le supremum de tous les cardinaux
$|\Ob(\mathcal{X}_{\Spec(k)})/\cong|$ et
$|\text{Flèches}(\mathcal{X}_{\Spec(k)})|$, où $k$ parcourt les corps de type
fini sur $S$, est $\leq$ au cardinal d'un certain objet de
$(\Sch/S)_{fppf}$.} que peuvent ignorer les lecteurs qui ne s'intéressent pas
aux questions ensemblistes ;
\item[{[0]}] $\mathcal{X}$ est un champ en groupoïdes pour la topologie étale ;
\item[{[1]}] $\mathcal{X}$ commute aux limites ;
\item[{[2]}] $\mathcal{X}$ satisfait la condition de Rim--Schlessinger (RS) ;
\item[{[3]}] les espaces $T\mathcal{F}_{\mathcal{X}, k, x_0}$ et
$\text{Inf}(\mathcal{F}_{\mathcal{X}, k, x_0})$
sont de dimension finie pour tous $k$ et $x_0$ ; voir
(\ref{equation-tangent-space}) et
(\ref{equation-infinitesimal-automorphisms}) ;
\item[{[4]}] le foncteur (\ref{equation-approximation}) est une équivalence ;
\item[{[5]}] $\mathcal{X}$ et
$\Delta : \mathcal{X} \to \mathcal{X} \times \mathcal{X}$ satisfont
l'ouverture de la versalité.
\end{enumerate}










\section{Axiomes pour les foncteurs}
\label{section-axioms-functors}

\noindent
Soit $S$ un schéma. Soit
$F : (\Sch/S)_{fppf}^{opp} \to \textit{Sets}$ un foncteur. Notons
$\mathcal{X} = \mathcal{S}_F$ la catégorie fibrée en ensembles associée à
$F$ ; voir Champs algébriques, section \ref{algebraic-section-split}. Dans
cette section, nous traduisons les résultats précédents appliqués à
$\mathcal{X}$ en énoncés sur $F$.

\medskip\noindent
Soit $S$ un schéma localement noethérien. Soit
$F : (\Sch/S)_{fppf}^{opp} \to \textit{Sets}$ un foncteur. Soit $k$ un corps
de type fini sur $S$. Soit $x_0 \in F(\Spec(k))$. La catégorie de
prédéformations associée (\ref{equation-predeformation-category}) correspond
au foncteur
$$
F_{k, x_0} : \mathcal{C}_\Lambda \longrightarrow \textit{Sets},
\quad
A \longmapsto \{x \in F(\Spec(A)) \mid x|_{\Spec(k)} = x_0 \}.
$$
Rappelons que nous ne distinguons pas les catégories cofibrées en ensembles
sur $\mathcal{C}_\Lambda$ des foncteurs
$\mathcal{C}_\Lambda \to \textit{Sets}$ ; voir Théorie formelle des
déformations, remarques
\ref{formal-defos-remarks-cofibered-groupoids}
(\ref{formal-defos-item-convention-cofibered-sets}).
Étant donnée une transformation de foncteurs $a : F \to G$, en posant
$y_0 = a(x_0)$, nous obtenons un morphisme
$$
F_{k, x_0} \longrightarrow G_{k, y_0}
$$
voir (\ref{equation-functoriality}). Le lemme
\ref{lemma-formally-smooth-on-deformation-categories} affirme que, si
$a : F \to G$ est formellement lisse (au sens de Compléments sur les
morphismes d'espaces, définition
\ref{spaces-more-morphisms-definition-formally-smooth-etale-unramified}), alors
$F_{k, x_0} \longrightarrow G_{k, y_0}$ est lisse au sens de Théorie formelle
des déformations, remarque
\ref{formal-defos-remark-compare-smooth-schlessinger}.

\medskip\noindent
Le lemme \ref{lemma-pushout} affirme que, si $Y' = Y \amalg_X X'$ dans la
catégorie des schémas sur $S$, où $X \to X'$ est un épaississement et
$X \to Y$ est affine, alors l'application
$$
F(Y \amalg_X X') \to F(Y) \times_{F(X)} F(X')
$$
est bijective, pourvu que $F$ soit un espace algébrique. Nous disons qu'un
foncteur général $F$ satisfait la {\it condition de Rim--Schlessinger}, ou que
$F$ {\it satisfait (RS)}, si, pour toute somme amalgamée
$Y' = Y \amalg_X X'$ où $Y, X, X'$ sont les spectres d'anneaux locaux
artiniens de type fini sur $S$, l'application
$$
F(Y \amalg_X X') \to F(Y) \times_{F(X)} F(X')
$$
est bijective. Ainsi, tout espace algébrique satisfait (RS).

\medskip\noindent
Le lemme \ref{lemma-deformation-category} affirme que, pour tout foncteur $F$
qui satisfait (RS), tous les $F_{k, x_0}$ sont des foncteurs de déformations
au sens de Théorie formelle des déformations, définition
\ref{formal-defos-definition-deformation-category}, c'est-à-dire qu'ils
satisfont (RS) au sens de Théorie formelle des déformations, remarque
\ref{formal-defos-remark-compare-schlessinger-H4}.
En particulier, l'espace tangent
$$
TF_{k, x_0} = \{x \in F(\Spec(k[\epsilon])) \mid x|_{\Spec(k)} = x_0\}
$$
est muni d'une structure d'espace vectoriel sur $k$ d'après Théorie formelle
des déformations, lemme
\ref{formal-defos-lemma-tangent-space-vector-space}.

\medskip\noindent
Le lemme \ref{lemma-finite-dimension} affirme qu'un espace algébrique $F$
localement de type fini sur $S$ donne des foncteurs de déformations
$F_{k, x_0}$ dont les espaces tangents $TF_{k, x_0}$ sont de dimension finie.

\medskip\noindent
Un {\it objet formel}\footnote{C'est ce qu'Artin appelle une déformation
formelle.} $\xi = (R, \xi_n)$ de $F$ consiste en une $S$-algèbre locale
noethérienne complète $R$ dont le corps résiduel est de type fini sur $S$,
ainsi qu'en des éléments $\xi_n \in F(\Spec(R/\mathfrak m^n))$ tels que
$\xi_{n + 1}|_{\Spec(R/\mathfrak m^n)} = \xi_n$. Un objet formel $\xi$
définit un objet formel $\xi$ de $F_{R/\mathfrak m, \xi_1}$. Nous disons que
$\xi$ est {\it versel} si et seulement s'il est versel au sens de Théorie
formelle des déformations, définition \ref{formal-defos-definition-versal}.
Un objet formel $\xi = (R, \xi_n)$ est dit {\it effectif} s'il existe un
$x \in F(\Spec(R))$ tel que $\xi_n = x|_{\Spec(R/\mathfrak m^n)}$ pour tout
$n \geq 1$. Le lemme \ref{lemma-effective} affirme que, si $F$ est un espace
algébrique, tout objet formel est effectif.

\medskip\noindent
Soit $U$ un schéma localement de type fini sur $S$ et soit $x \in F(U)$.
Soit $u_0 \in U$ un point de type fini. Nous disons que $x$ est versel en
$u_0$ si et seulement si
$\xi = (\mathcal{O}_{U, u_0}^\wedge,
x|_{\Spec(\mathcal{O}_{U, u_0}/\mathfrak m_{u_0}^n)})$
est un objet formel versel au sens décrit ci-dessus.

\medskip\noindent
Soit $S$ un schéma localement noethérien. Soit
$F : (\Sch/S)_{fppf}^{opp} \to \textit{Sets}$ un foncteur.
Voici les axiomes que nous considérerons pour $F$.
\begin{enumerate}
\item[{[-1]}] une condition ensembliste\footnote{La condition est la suivante :
le supremum de tous les cardinaux $|F(\Spec(k))|$, où $k$ parcourt les corps
de type fini sur $S$, est $\leq$ au cardinal d'un certain objet de
$(\Sch/S)_{fppf}$.} que peuvent ignorer les lecteurs qui ne s'intéressent pas
aux questions ensemblistes ;
\item[{[0]}] $F$ est un faisceau pour la topologie étale ;
\item[{[1]}] $F$ commute aux limites ;
\item[{[2]}] $F$ satisfait la condition de Rim--Schlessinger (RS) ;
\item[{[3]}] tout espace tangent $TF_{k, x_0}$ est de dimension finie ;
\item[{[4]}] tout objet formel est effectif ;
\item[{[5]}] $F$ satisfait l'ouverture de la versalité.
\end{enumerate}
Ici, {\it commute aux limites} désigne la notion définie dans Limites
d'espaces, définition \ref{spaces-limits-definition-locally-finite-presentation},
et {\it l'ouverture de la versalité} signifie ceci : pour tout schéma $U$
localement de type fini sur $S$, tout $x \in F(U)$ et tout point de type fini
$u_0 \in U$ tel que $x$ soit versel en $u_0$, il existe un voisinage ouvert
$u_0 \in U' \subset U$ tel que $x$ soit versel en tout point de type fini de
$U'$.






\section{Espaces algébriques}
\label{section-algebraic-spaces}

\noindent
Voici notre premier résultat principal sur les espaces algébriques.

\begin{proposition}
\label{proposition-spaces-diagonal-representable}
Soit $S$ un schéma localement noethérien. Soit
$F : (\Sch/S)_{fppf}^{opp} \to \textit{Sets}$ un foncteur. Supposons que
\begin{enumerate}
\item $\Delta : F \to F \times F$ soit représentable par des espaces
algébriques ;
\item $F$ satisfasse les axiomes [-1], [0], [1], [2], [3], [4], [5] (voir la
section \ref{section-axioms-functors}) ;
\item $\mathcal{O}_{S, s}$ soit un anneau G pour tout point de type fini $s$
de $S$.
\end{enumerate}
Alors $F$ est un espace algébrique.
\end{proposition}

\begin{proof}
Le lemme \ref{lemma-get-smooth} s'applique à $F$. Nous choisissons donc, pour
tout corps de type fini $k$ sur $S$ et tout $x_0 \in F(\Spec(k))$, un schéma
affine $U_{k, x_0}$ de type fini sur $S$ et un morphisme lisse
$U_{k, x_0} \to F$ tels qu'il existe un point de type fini
$u_{k, x_0} \in U_{k, x_0}$ de corps résiduel $k$ tel que $x_0$ soit l'image
de $u_{k, x_0}$.
Alors
$$
U = \coprod\nolimits_{k, x_0} U_{k, x_0} \longrightarrow F
$$
est lisse\footnote{Remarque ensembliste : ce coproduit est isomorphe à un
objet de $(\Sch/S)_{fppf}$, car l'axiome [-1] borne l'ensemble d'indices ;
voir Ensembles, lemme \ref{sets-lemma-what-is-in-it}.}. Pour achever la
démonstration, il suffit de montrer que cette application est surjective ;
voir Bootstrap, lemme \ref{bootstrap-lemma-spaces-etale-smooth-cover} (c'est
ici que nous utilisons l'axiome [0]). D'après Critères de représentabilité,
lemme \ref{criteria-lemma-check-property-limit-preserving}, il suffit de
montrer que $U \times_F V \to V$ est surjectif pour les $V \to F$ où $V$ est
un schéma affine localement de présentation finie sur $S$. Comme
$U \times_F V \to V$ est lisse, son image est ouverte. Il suffit donc de
montrer que l'image de $U \times_F V \to V$ contient tous les points de type
fini de $V$ ; voir Morphismes, lemme
\ref{morphisms-lemma-enough-finite-type-points}. Soit $v_0 \in V$ un point de
type fini. Alors $k = \kappa(v_0)$ est un corps de type fini sur $S$. Notons
$x_0$ la composée $\Spec(k) \xrightarrow{v_0} V \to F$. Alors
$(u_{k, x_0}, v_0) : \Spec(k) \to U \times_F V$ est un point qui s'envoie sur
$v_0$, ce qui conclut.
\end{proof}

\begin{lemma}
\label{lemma-monomorphism}
Soit $S$ un schéma localement noethérien. Soit $a : F \to G$ une
transformation de foncteurs $(\Sch/S)_{fppf}^{opp} \to \textit{Sets}$.
Supposons que
\begin{enumerate}
\item $a$ soit injective ;
\item $F$ satisfasse les axiomes [0], [1], [2], [4] et [5] ;
\item $\mathcal{O}_{S, s}$ soit un anneau G pour tout point de type fini $s$
de $S$ ;
\item $G$ soit un espace algébrique localement de type fini sur $S$.
\end{enumerate}
Alors $F$ est un espace algébrique.
\end{lemma}

\begin{proof}
D'après le lemme \ref{lemma-finite-dimension}, le foncteur $G$ satisfait [3].
Comme $F \to G$ est injectif, nous en concluons que $F$ satisfait aussi [3].
De plus, l'injectivité de $F \to G$ montre que, pour des schémas $U$, $V$ et
des morphismes $U \to F$ et $V \to F$, nous avons
$U \times_F V = U \times_G V$. Ainsi, $\Delta : F \to F \times F$ est
représentable (par des schémas), puisque cela vaut par hypothèse pour $G$.
La proposition \ref{proposition-spaces-diagonal-representable} s'applique
donc\footnote{La condition ensembliste [-1] vaut pour $F$, puisqu'elle vaut
pour $G$. Nous omettons les détails.}.
\end{proof}












\section{Champs algébriques}
\label{section-algebraic-stacks}

\noindent
La proposition \ref{proposition-second-diagonal-representable} est notre
premier résultat principal sur les champs algébriques.

\begin{lemma}
\label{lemma-diagonal-representable}
Soit $S$ un schéma localement noethérien. Soit
$p : \mathcal{X} \to (\Sch/S)_{fppf}$ une catégorie fibrée en groupoïdes.
Supposons que
\begin{enumerate}
\item $\Delta : \mathcal{X} \to \mathcal{X} \times \mathcal{X}$
soit représentable par des espaces algébriques ;
\item $\mathcal{X}$ satisfasse les axiomes [-1], [0], [1], [2], [3] (voir la
section \ref{section-axioms}) ;
\item tout objet formel de $\mathcal{X}$ soit effectif ;
\item $\mathcal{X}$ satisfasse l'ouverture de la versalité ;
\item $\mathcal{O}_{S, s}$ soit un anneau G pour tout point de type fini $s$
de $S$.
\end{enumerate}
Alors $\mathcal{X}$ est un champ algébrique.
\end{lemma}

\begin{proof}
Le lemme \ref{lemma-get-smooth} s'applique à $\mathcal{X}$. Nous choisissons
donc, pour tout corps de type fini $k$ sur $S$ et toute classe d'isomorphisme
d'objets $x_0 \in \Ob(\mathcal{X}_{\Spec(k)})$, un schéma affine
$U_{k, x_0}$ de type fini sur $S$ et un morphisme lisse
$(\Sch/U_{k, x_0})_{fppf} \to \mathcal{X}$ tels qu'il existe un point de type
fini $u_{k, x_0} \in U_{k, x_0}$ de corps résiduel $k$ tel que $x_0$ soit
l'image de $u_{k, x_0}$. Alors
$$
(\Sch/U)_{fppf} \to \mathcal{X},
\quad\text{avec}\quad
U = \coprod\nolimits_{k, x_0} U_{k, x_0}
$$
est lisse\footnote{Remarque ensembliste : ce coproduit est isomorphe à un
objet de $(\Sch/S)_{fppf}$, car l'axiome [-1] borne l'ensemble d'indices ;
voir Ensembles, lemme \ref{sets-lemma-what-is-in-it}.}. Pour achever la
démonstration, il suffit de montrer que cette application est surjective ;
voir Critères de représentabilité, lemme \ref{criteria-lemma-stacks-etale}
(c'est ici que nous utilisons l'axiome [0]). D'après Critères de
représentabilité, lemme
\ref{criteria-lemma-check-property-limit-preserving}, il suffit de montrer que
$(\Sch/U)_{fppf} \times_\mathcal{X} (\Sch/V)_{fppf} \to (\Sch/V)_{fppf}$
est surjectif pour les $y : (\Sch/V)_{fppf} \to \mathcal{X}$ où $V$ est un
schéma affine localement de présentation finie sur $S$. D'après l'hypothèse
(1), le produit fibré
$(\Sch/U)_{fppf} \times_\mathcal{X} (\Sch/V)_{fppf}$ est représentable
est représentable par un espace algébrique $W$. Alors $W \to V$ est lisse,
donc son image est ouverte. Il suffit ainsi de montrer que l'image de
$W \to V$ contient tous les points de type fini de $V$ ; voir Morphismes,
lemme \ref{morphisms-lemma-enough-finite-type-points}. Soit $v_0 \in V$ un
point de type fini. Alors $k = \kappa(v_0)$ est un corps de type fini sur $S$.
Notons $x_0 = y|_{\Spec(k)}$ l'image réciproque de $y$ par $v_0$. Alors
$(u_{k, x_0}, v_0)$ donne un morphisme $\Spec(k) \to W$ dont la composée avec
$W \to V$ est $v_0$, ce qui conclut.
\end{proof}

\begin{proposition}
\label{proposition-second-diagonal-representable}
Soit $S$ un schéma localement noethérien. Soit
$p : \mathcal{X} \to (\Sch/S)_{fppf}$ une catégorie fibrée en groupoïdes.
Supposons que
\begin{enumerate}
\item $\Delta_\Delta : \mathcal{X} \to
\mathcal{X} \times_{\mathcal{X} \times \mathcal{X}} \mathcal{X}$
soit représentable par des espaces algébriques ;
\item $\mathcal{X}$ satisfasse les axiomes [-1], [0], [1], [2], [3], [4] et
[5] (voir la section \ref{section-axioms}) ;
\item $\mathcal{O}_{S, s}$ soit un anneau G pour tout point de type fini $s$
de $S$.
\end{enumerate}
Alors $\mathcal{X}$ est un champ algébrique.
\end{proposition}

\begin{proof}
Montrons d'abord que
$\Delta : \mathcal{X} \to \mathcal{X} \times \mathcal{X}$ est représentable
par des espaces algébriques. Il suffit pour cela de montrer que
$$
\mathcal{Y} =
\mathcal{X} \times_{\Delta, \mathcal{X} \times \mathcal{X}, y} (\Sch/V)_{fppf}
$$
est représentable par un espace algébrique pour tout schéma affine $V$
localement de présentation finie sur $S$ et tout objet $y$ de
$\mathcal{X} \times \mathcal{X}$ sur $V$ ; voir Critères de représentabilité,
lemme
\ref{criteria-lemma-check-representable-limit-preserving}\footnote{La
condition ensembliste de Critères de représentabilité, lemme
\ref{criteria-lemma-check-representable-limit-preserving}
sera satisfaite : le cardinal de l'espace algébrique $Y$ représentant
$\mathcal{Y}$ est convenablement borné. En effet, $Y \to S$ sera localement
de type fini et $Y$ satisfera l'axiome [-1]. Nous omettons les détails.}.
Notons que $\mathcal{Y}$ est fibrée en ensemblöïdes (Champs, lemme
\ref{stacks-lemma-isom-as-2-fibre-product}) et soit
$Y : (\Sch/S)_{fppf}^{opp} \to \textit{Sets}$,
$T \mapsto \Ob(\mathcal{Y}_T)/\cong$, le foncteur des classes d'isomorphisme.
Nous allons appliquer la proposition
\ref{proposition-spaces-diagonal-representable} pour voir que $Y$ est un
espace algébrique.

\medskip\noindent
Notons que
$\Delta_\mathcal{Y} : \mathcal{Y} \to \mathcal{Y} \times \mathcal{Y}$
(et donc aussi $Y \to Y \times Y$) est représentable par des espaces
algébriques d'après la condition (1) et Critères de représentabilité, lemme
\ref{criteria-lemma-second-diagonal}. Notons que $Y$ est un faisceau pour la
topologie étale d'après Champs, lemmes
\ref{stacks-lemma-stack-in-setoids-characterize} et
\ref{stacks-lemma-2-fibre-product-gives-stack-in-setoids} ; autrement dit,
l'axiome [0] est satisfait. De plus, $Y$ commute aux limites d'après le lemme
\ref{lemma-fibre-product-limit-preserving} ; nous avons donc [1]. Notons que
$Y$ satisfait (RS), c'est-à-dire l'axiome [2], d'après les lemmes
\ref{lemma-algebraic-stack-RS} et \ref{lemma-fibre-product-RS}. L'axiome [3]
pour $Y$ résulte des lemmes \ref{lemma-finite-dimension} et
\ref{lemma-fibre-product-tangent-spaces}. L'axiome [4] résulte des lemmes
\ref{lemma-effective} et \ref{lemma-fibre-product-effective}. L'axiome [5]
pour $Y$ résulte directement de l'ouverture de la versalité pour
$\Delta_\mathcal{X}$, qui fait partie de l'axiome [5] pour $\mathcal{X}$.
Toutes les hypothèses de la proposition
\ref{proposition-spaces-diagonal-representable} sont donc satisfaites, et $Y$
est un espace algébrique.

\medskip\noindent
Le lemme \ref{lemma-diagonal-representable} montre alors que $\mathcal{X}$ est
un champ algébrique.
\end{proof}





\section{Condition forte de Rim--Schlessinger}
\label{section-RS-star}

\noindent
Dans la suite de ce chapitre, la version strictement plus forte suivante des
conditions de Rim--Schlessinger jouera un rôle important.

\begin{definition}
\label{definition-RS-star}
Soit $S$ un schéma. Soit $\mathcal{X}$ une catégorie fibrée en groupoïdes sur
$(\Sch/S)_{fppf}$. Nous disons que $\mathcal{X}$ satisfait la
{\it condition (RS*)} si, pour tout diagramme de produit fibré
$$
\xymatrix{
B' \ar[r] & B \\
A' = A \times_B B' \ar[u] \ar[r] & A \ar[u]
}
$$
de $S$-algèbres, où $B' \to B$ est surjectif de noyau de carré nul, le
foncteur de catégories-fibres
$$
\mathcal{X}_{\Spec(A')}
\longrightarrow
\mathcal{X}_{\Spec(A)} \times_{\mathcal{X}_{\Spec(B)}} \mathcal{X}_{\Spec(B')}
$$
est une équivalence de catégories.
\end{definition}

\noindent
Faisons quelques observations. Pour $A \to B \leftarrow B'$ comme dans la
définition \ref{definition-RS-star},
\begin{enumerate}
\item nous avons $\Spec(A') = \Spec(A) \amalg_{\Spec(B)} \Spec(B')$ dans la
catégorie des schémas ; voir Compléments sur les morphismes, lemme
\ref{more-morphisms-lemma-pushout-along-thickening} ;
\item si $\mathcal{X}$ est un champ algébrique, alors $\mathcal{X}$ satisfait
(RS*) d'après le lemme \ref{lemma-algebraic-stack-RS-star}.
\end{enumerate}
Si $S$ est localement noethérien, alors
\begin{enumerate}
\item[(3)] si $A$, $B$, $B'$ sont de type fini sur $S$ et si $B$ est fini sur
$A$, alors $A'$ est de type fini sur $S$\footnote{Si $\Spec(A)$ s'envoie dans
un ouvert affine de $S$, cela résulte de Compléments d'algèbre, lemme
\ref{more-algebra-lemma-fibre-product-finite-type}. Le cas général résulte de
Compléments d'algèbre, lemme \ref{more-algebra-lemma-diagram-localize}.} ;
\item[(4)] si $\mathcal{X}$ satisfait (RS*), alors $\mathcal{X}$ satisfait
(RS), car (RS) couvre exactement les cas de (RS*) où $A$, $B$, $B'$ sont des
anneaux locaux artiniens.
\end{enumerate}

\begin{lemma}
\label{lemma-algebraic-stack-RS-star}
Soit $\mathcal{X}$ un champ algébrique sur une base $S$. Alors $\mathcal{X}$
satisfait (RS*).
\end{lemma}

\begin{proof}
Cela résulte du lemme \ref{lemma-pushout} ; voir les remarques qui suivent la
définition \ref{definition-RS-star}.
\end{proof}

\begin{lemma}
\label{lemma-fibre-product-RS-star}
Soit $S$ un schéma. Soient $p : \mathcal{X} \to \mathcal{Y}$ et
$q : \mathcal{Z} \to \mathcal{Y}$ des $1$-morphismes de catégories fibrées
en groupoïdes sur $(\Sch/S)_{fppf}$. Si $\mathcal{X}$, $\mathcal{Y}$ et
$\mathcal{Z}$ satisfont (RS*), alors
$\mathcal{X} \times_\mathcal{Y} \mathcal{Z}$ satisfait aussi (RS*).
\end{lemma}

\begin{proof}
La démonstration est exactement la même que celle du lemme
\ref{lemma-fibre-product-RS}.
\end{proof}








\section{Versalité et générisations}
\label{section-generalize-versality}

\noindent
Nous démontrons que la versalité se conserve par générisation pour les champs
qui satisfont (RS*) et commutent aux limites. Nous suggérons d'omettre cette
section en première lecture.

\begin{lemma}
\label{lemma-single-point}
Soit $S$ un schéma localement noethérien. Soit $\mathcal{X}$ une catégorie
fibrée en groupoïdes sur $(\Sch/S)_{fppf}$ satisfaisant (RS*). Soit $x$ un
objet de $\mathcal{X}$ sur un schéma affine $U$ de type fini sur $S$. Soit
$u \in U$ un point de type fini tel que $x$ ne soit pas versel en $u$. Alors
il existe un morphisme $x \to y$ de $\mathcal{X}$ situé au-dessus de
$U \to T$ et satisfaisant aux conditions suivantes :
\begin{enumerate}
\item le morphisme $U \to T$ est un épaississement du premier ordre ;
\item nous avons une suite exacte courte
$$
0 \to \kappa(u) \to \mathcal{O}_T \to \mathcal{O}_U \to 0
$$
\item il n'existe {\bf aucun} couple $(W, \alpha)$ formé d'un voisinage ouvert
$W \subset T$ de $u$ et d'un morphisme $\beta : y|_W \to x$ tels que la
composée
$$
x|_{U \cap W} \xrightarrow{\text{restriction de }x \to y}
y|_W \xrightarrow{\beta} x
$$
soit le morphisme canonique $x|_{U \cap W} \to x$.
\end{enumerate}
\end{lemma}

\begin{proof}
Soit $R = \mathcal{O}_{U, u}^\wedge$. Soit $k = \kappa(u)$ le corps résiduel
de $R$. Soit $\xi$ l'objet formel de $\mathcal{X}$ sur $R$ associé à $x$.
Puisque $x$ n'est pas versel en $u$, $\xi$ n'est pas versel ; voir le lemme
\ref{lemma-versality-matches}. D'après la discussion qui suit la définition
\ref{definition-versal-formal-object}, cela signifie qu'il existe des
morphismes $\xi_1 \to x_A \to x_B$ de $\mathcal{X}$ situés au-dessus
d'immersions fermées $\Spec(k) \to \Spec(A) \to \Spec(B)$, où $A, B$ sont des
anneaux locaux artiniens de corps résiduel $k$, un $n \geq 1$ et un diagramme
commutatif
$$
\vcenter{
\xymatrix{
& x_A \ar[ld] \\
\xi_n & \xi_1 \ar[u] \ar[l]
}
}
\quad\text{au-dessus de}\quad
\vcenter{
\xymatrix{
& \Spec(A) \ar[ld] \\
\Spec(R/\mathfrak m^n) & \Spec(k) \ar[u] \ar[l]
}
}
$$
tels qu'il n'existe {\bf aucun} $m \geq n$ ni diagramme commutatif
$$
\vcenter{
\xymatrix{
& & x_B \ar[lldd] \\
& & x_A \ar[ld] \ar[u] \\
\xi_m & \xi_n \ar[l] & \xi_1 \ar[u] \ar[l]
}
}
\quad\text{au-dessus de}
\vcenter{
\xymatrix{
& & \Spec(B) \ar[lldd] \\
& & \Spec(A) \ar[ld] \ar[u] \\
\Spec(R/\mathfrak m^m) &
\Spec(R/\mathfrak m^n) \ar[l] &
\Spec(k) \ar[u] \ar[l]
}
}
$$
Nous pouvons de plus supposer que $B \to A$ est une petite extension,
c'est-à-dire que le noyau $I$ de la surjection $B \to A$ est isomorphe à $k$
comme $A$-module. Cela résulte de Théorie formelle des déformations, remarque
\ref{formal-defos-remark-versal-object}. Posons simplement
$$
T = U \amalg_{\Spec(A)} \Spec(B)
$$
D'après la propriété (RS*), il existe $y$ sur $T$ dont la restriction à
$\Spec(B)$ est $x_B$ et dont la restriction à $U$ est $x$ (ce qui donne la
flèche $x \to y$ située au-dessus de $U \to T$). Pour achever la
démonstration, vérifions les conditions (1), (2) et (3).

\medskip\noindent
Par construction de la somme amalgamée, nous avons un diagramme commutatif
$$
\xymatrix{
0 \ar[r] &
I \ar[r] &
B \ar[r] &
A \ar[r] &
0 \\
0 \ar[r] &
I \ar[r] \ar[u] &
\Gamma(T, \mathcal{O}_T) \ar[r] \ar[u] &
\Gamma(U, \mathcal{O}_U) \ar[r] \ar[u] &
0
}
$$
à lignes exactes. Cela démontre immédiatement (1) et (2). Pour achever la
démonstration, raisonnons par l'absurde. Supposons disposer d'un couple
$(W, \beta)$ comme dans (3). Puisque $\Spec(B) \to T$ se factorise par $W$,
nous obtenons le morphisme
$$
x_B \to y|_W \xrightarrow{\beta} x
$$
Comme $B$ est local artinien de corps résiduel $k = \kappa(u)$, nous voyons
que $x_B \to x$ est situé au-dessus d'un morphisme $\Spec(B) \to U$ qui se
factorise par $\Spec(\mathcal{O}_{U, u}/\mathfrak m_u^m)$ pour un certain
$m \geq n$. Autrement dit, $x_B \to x$ se factorise par $\xi_m$, ce qui donne
une application $x_B \to \xi_m$. La condition de compatibilité portant sur le
morphisme $\alpha$ dans la condition (3) signifie que
$$
\xymatrix{
x_B \ar[d] & x_A \ar[d] \ar[l] \\
\xi_m & \xi_n \ar[l]
}
$$
est commutatif. C'est la contradiction recherchée.
\end{proof}

\begin{lemma}
\label{lemma-generalization-versality}
Soit $S$ un schéma localement noethérien. Soit $\mathcal{X}$ une catégorie
fibrée en groupoïdes sur $(\Sch/S)_{fppf}$. Supposons que
\begin{enumerate}
\item $\Delta : \mathcal{X} \to \mathcal{X} \times \mathcal{X}$ soit
représentable par des espaces algébriques ;
\item $\mathcal{X}$ satisfasse (RS*) ;
\item $\mathcal{X}$ commute aux limites.
\end{enumerate}
Soit $x$ un objet de $\mathcal{X}$ sur un schéma $U$ de type fini sur $S$.
Soit $u \leadsto u_0$ une spécialisation de points de type fini de $U$ telle
que $x$ soit versel en $u_0$. Alors $x$ est versel en $u$.
\end{lemma}

\begin{proof}
Après avoir rétréci $U$, nous pouvons supposer que $U$ est affine et que $U$ s'envoie
dans un ouvert affine $\Spec(\Lambda)$ de $S$. Si $x$ n'est pas versel en $u$,
choisissons $x \to y$ situé au-dessus de $U \to T$ comme dans le lemme
\ref{lemma-single-point}. Écrivons $U = \Spec(R_0)$ et $T = \Spec(R)$. Le
morphisme $U \to T$ correspond à un homomorphisme d'anneaux surjectif
$R \to R_0$ dont le noyau est un idéal de carré nul. D'après l'hypothèse (3),
$y$ provient d'un objet $x'$ sur $U' = \Spec(R')$, pour une
$\Lambda$-sous-algèbre de type fini $R' \subset R$. Quitte à agrandir $R'$,
nous pouvons supposer, et supposons, que $R' \to R_0$ est surjectif, de sorte
que $U \subset U'$ est un épaississement du premier ordre. Nous avons donc
$$
x \to y \to x'
\text{ au-dessus de }
U \to T \to U'
$$
D'après l'hypothèse (1), il existe un espace algébrique $Z$ sur $S$ représentant
$$
(\Sch/U)_{fppf} \times_{x, \mathcal{X}, x'} (\Sch/U')_{fppf}
$$
voir Champs algébriques, lemme \ref{algebraic-lemma-representable-diagonal}.
Par construction des $2$-produits fibrés, un $V$-point de $Z$
correspond à un triplet $(a, a', \alpha)$ formé de morphismes
$a : V \to U$, $a' : V \to U'$ et d'un morphisme
$\alpha : a^*x \to (a')^*x'$. Nous obtenons un diagramme commutatif
$$
\xymatrix{
U \ar[rd] \ar[rdd] \ar[rrd] \\
& Z \ar[r]_{p'} \ar[d]^p & U' \ar[d] \\
& U \ar[r] & S
}
$$
Le morphisme $i : U \to Z$ provient de l'isomorphisme $x \to x'|_U$. Soit
$z_0 = i(u_0) \in Z$. D'après le lemme \ref{lemma-base-change-versal}, le
morphisme $Z \to U'$ est lisse en $z_0$. En remplaçant $U$ par un voisinage
ouvert affine de $u_0$, $U'$ par l'ouvert correspondant et $Z$ par
l'intersection des images réciproques de ces ouverts par $p$ et $p'$, nous
arrivons à la situation où $Z \to U'$ est lisse le long de $i(U)$. Puisque
$u \leadsto u_0$, le point $u$ appartient à cet ouvert. La condition (3) du
lemme \ref{lemma-single-point} est manifestement préservée quand on rétrécit
$U$ (tous les schémas $U$, $T$, $U'$ ont le même espace topologique
sous-jacent). Comme $U \to U'$ est un épaississement du premier ordre de
schémas affines, nous pouvons choisir un morphisme $i' : U' \to Z$ tel que
$p' \circ i' = \text{id}_{U'}$ et dont la restriction à $U$ soit $i$
(Compléments sur les morphismes d'espaces, lemme
\ref{spaces-more-morphisms-lemma-smooth-formally-smooth}). En prenant l'image
réciproque du morphisme universel $p^*x \to (p')^*x'$ par $i'$, nous obtenons
un morphisme
$$
x' \to x
$$
situé au-dessus de $p \circ i' : U' \to U$ tel que la composée
$$
x \to x' \to x
$$
soit l'identité. Rappelons que nous avons $y \to x'$ situé au-dessus du
morphisme $T \to U'$. Par composition, nous obtenons un morphisme $y \to x$
dont l'existence contredit la condition (3) du lemme
\ref{lemma-single-point}. Cette contradiction achève la démonstration.
\end{proof}






\section{Effectivité formelle forte}
\label{section-strong-formal-effectiveness}

\noindent
Dans cette section, nous montrons comment une version forte de l'effectivité
des objets formels entraîne l'ouverture de la versalité. La démonstration de
\cite[Theorem 1.1]{Bhatt-Algebraize} montre que les espaces algébriques
quasi-compacts et quasi-séparés satisfont à l'effectivité formelle forte
étudiée dans la remarque \ref{remark-strong-effectiveness}. En outre, la
théorie que nous développons n'est pas vide : nous l'utiliserons plus loin
pour établir l'ouverture de la versalité pour le champ des faisceaux cohérents
et pour les espaces de modules de complexes ; voir Quot, théorèmes
\ref{quot-theorem-coherent-algebraic-general} et
\ref{quot-theorem-complexes-algebraic}.

\begin{lemma}
\label{lemma-infinite-sequence-pre}
Soit $S$ un schéma localement noethérien. Soit $\mathcal{X}$ une catégorie
fibrée en groupoïdes sur $(\Sch/S)_{fppf}$ qui possède (RS*).
Soit $x$ un objet de $\mathcal{X}$ au-dessus d'un schéma affine $U$ de type
fini sur $S$. Soient $u_n \in U$, $n \geq 1$, des points de type fini tels que
(a) il n'existe aucune spécialisation $u_n \leadsto u_m$ pour $n \not = m$ et
(b) $x$ ne soit pas versel en $u_n$ pour tout $n$. Il existe alors des morphismes
$$
x \to x_1 \to x_2 \to \ldots
\quad\text{dans }\mathcal{X}\text{ situés au-dessus de }\quad
U \to U_1 \to U_2 \to \ldots
$$
sur $S$, tels que
\begin{enumerate}
\item pour tout $n$, le morphisme $U \to U_n$ soit un épaississement du
premier ordre ;
\item pour tout $n$, on ait une suite exacte courte
$$
0 \to \kappa(u_n) \to \mathcal{O}_{U_n} \to \mathcal{O}_{U_{n - 1}} \to 0
$$
où $U_0 = U$ pour $n = 1$ ;
\item pour tout $n$, il n'existe {\bf pas} de paire $(W, \alpha)$ formée
d'un voisinage ouvert $W \subset U_n$ de $u_n$ et d'un morphisme
$\alpha : x_n|_W \to x$ tels que la composée
$$
x|_{U \cap W} \xrightarrow{\text{restriction de }x \to x_n}
x_n|_W \xrightarrow{\alpha} x
$$
soit le morphisme canonique $x|_{U \cap W} \to x$.
\end{enumerate}
\end{lemma}

\begin{proof}
Puisqu'il n'y a aucune spécialisation entre les points $u_n$ (et qu'en
particulier les $u_n$ sont deux à deux distincts), nous pouvons, pour tout
$n$, trouver un ouvert $U' \subset U$ tel que $u_n \in U'$ et
$u_i \not \in U'$ pour $i = 1, \ldots, n - 1$. D'après le lemme
\ref{lemma-single-point}, pour tout $n \geq 1$, nous pouvons trouver
$$
x \to y_n
\quad\text{dans }\mathcal{X}\text{ situé au-dessus de}\quad
U \to T_n
$$
tel que
\begin{enumerate}
\item le morphisme $U \to T_n$ soit un épaississement du premier ordre ;
\item on ait une suite exacte courte
$$
0 \to \kappa(u_n) \to \mathcal{O}_{T_n} \to \mathcal{O}_U \to 0
$$
\item il n'existe {\bf pas} de paire $(W, \alpha)$ formée d'un voisinage
ouvert $W \subset T_n$ de $u_n$ et d'un morphisme
$\beta : y_n|_W \to x$ tels que la composée
$$
x|_{U \cap W} \xrightarrow{\text{restriction de }x \to y_n}
y_n|_W \xrightarrow{\beta} x
$$
soit le morphisme canonique $x|_{U \cap W} \to x$.
\end{enumerate}
Nous pouvons donc définir par récurrence
$$
U_1 = T_1, \quad
U_{n + 1} = U_n \amalg_U T_{n + 1}
$$
En posant $x_1 = y_1$ et en utilisant (RS*), nous obtenons par récurrence
$x_{n + 1}$ au-dessus de $U_{n + 1}$, dont les restrictions sont $x_n$
au-dessus de $U_n$ et $y_{n + 1}$ au-dessus de $T_{n + 1}$.
La propriété (1) pour $U \to U_n$ résulte de la construction de la somme
amalgamée dans Compléments sur les morphismes, lemme
\ref{more-morphisms-lemma-pushout-along-thickening}.
De même, la propriété (2) pour $U_n$ résulte de la propriété (2) pour $T_n$
par la construction de la somme amalgamée. Après restriction au voisinage
ouvert $U'$ de $u_n$ considéré plus haut, la propriété (3) pour $(U_n, x_n)$
résulte de la propriété (3) pour $(T_n, y_n)$, simplement parce que les
sous-schémas ouverts correspondants de $T_n$ et de $U_n$ sont isomorphes.
Nous omettons quelques détails.
\end{proof}

\begin{remark}[Effectivité forte]
\label{remark-strong-effectiveness}
Soit $S$ un schéma localement noethérien. Soit $\mathcal{X}$ une catégorie
fibrée en groupoïdes sur $(\Sch/S)_{fppf}$. Supposons donnés
\begin{enumerate}
\item un ouvert affine $\Spec(\Lambda) \subset S$ ;
\item un système projectif $(R_n)$ de $\Lambda$-algèbres dont les morphismes
de transition sont surjectifs et ont des noyaux localement nilpotents ;
\item un système $(\xi_n)$ d'objets de $\mathcal{X}$ situé au-dessus du
système $(\Spec(R_n))$.
\end{enumerate}
Dans cette situation, posons $R = \lim R_n$. On dit que $(\xi_n)$ est
{\it effectif} s'il existe un objet $\xi$ de $\mathcal{X}$ au-dessus de
$\Spec(R)$ dont les restrictions aux $\Spec(R_n)$ redonnent le système
$(\xi_n)$.
\end{remark}

\noindent
Il n'est pas vrai que tout champ algébrique $\mathcal{X}$ sur $S$ satisfasse
à un axiome d'effectivité forte de la forme suivante : tout système $(\xi_n)$
comme dans la remarque \ref{remark-strong-effectiveness} est effectif. Un
exemple est donné dans Exemples, section
\ref{examples-section-non-formal-effectiveness}.

\begin{lemma}
\label{lemma-SGE-implies-openness-versality}
Soit $S$ un schéma localement noethérien. Soit $\mathcal{X}$ une catégorie
fibrée en groupoïdes sur $(\Sch/S)_{fppf}$. Supposons que
\begin{enumerate}
\item $\Delta : \mathcal{X} \to \mathcal{X} \times \mathcal{X}$ soit
représentable par des espaces algébriques ;
\item $\mathcal{X}$ possède (RS*) ;
\item $\mathcal{X}$ commute aux limites ;
\item les systèmes $(\xi_n)$ comme dans la remarque
\ref{remark-strong-effectiveness}, où $\Ker(R_m \to R_n)$ est un idéal de
carré nul pour tous $m \geq n$, soient effectifs.
\end{enumerate}
Alors $\mathcal{X}$ satisfait à l'ouverture de la versalité.
\end{lemma}

\begin{proof}
Choisissons un schéma $U$ localement de type fini sur $S$, un point de type
fini $u_0$ de $U$ et un objet $x$ de $\mathcal{X}$ au-dessus de $U$ tel que
$x$ soit versel en $u_0$. Après avoir rétréci $U$, nous pouvons supposer que
$U$ est affine et que $U$ s'envoie dans un ouvert affine $\Spec(\Lambda)$ de $S$.
Soit $E \subset U$ l'ensemble des points de type fini $u$ tels que $x$ ne
soit pas versel en $u$. D'après le lemme
\ref{lemma-generalization-versality}, si $u \in E$, alors $u_0$ n'est pas une
spécialisation de $u$. Si l'ouverture de la versalité n'est pas satisfaite,
alors $u_0$ appartient à l'adhérence $\overline{E}$ de $E$. D'après
Propriétés, lemme \ref{properties-lemma-countable-dense-subset}, nous pouvons
choisir un sous-ensemble dénombrable $E' \subset E$ ayant la même adhérence
que $E$. D'après Propriétés, lemme \ref{properties-lemma-maximal-points}, nous
pouvons supposer qu'il n'y a aucune spécialisation entre les points de $E'$.
Remarquons que $E'$ doit être infini (dénombrable), puisque, comme nous venons
de le voir, $u_0$ n'est la spécialisation d'aucun point de $E'$. Nous pouvons
donc écrire $E' = \{u_1, u_2, u_3, \ldots\}$ ; il n'y a aucune
spécialisation entre les $u_i$, et $u_0$ appartient à l'adhérence de $E'$.

\medskip\noindent
Choisissons $x \to x_1 \to x_2 \to \ldots$ situé au-dessus de
$U \to U_1 \to U_2 \to \ldots$ comme dans le lemme
\ref{lemma-infinite-sequence-pre}. Écrivons $U_n = \Spec(R_n)$ et
$U = \Spec(R_0)$. Posons $R = \lim R_n$. Remarquons que $R \to R_0$ est
surjectif et que son noyau est un idéal de carré nul. Par l'hypothèse (4),
nous obtenons $\xi$ au-dessus de $\Spec(R)$ dont le changement de base à
$R_n$ est $x_n$. Par l'hypothèse (3), $\xi$ provient d'un objet $\xi'$ au-dessus
de $U' = \Spec(R')$ pour une certaine sous-$\Lambda$-algèbre de type fini
$R' \subset R$. En agrandissant $R'$, nous pouvons et voulons supposer que
$R' \to R_0$ est surjectif, de sorte que $U \subset U'$ soit un
épaississement du premier ordre. Nous avons donc maintenant
$$
x \to x_1 \to x_2 \to \ldots \to \xi'
\text{ situé au-dessus de }
U \to U_1 \to U_2 \to \ldots \to U'
$$
Par l'hypothèse (1), il existe un espace algébrique $Z$ sur $S$ représentant
$$
(\Sch/U)_{fppf} \times_{x, \mathcal{X}, \xi'} (\Sch/U')_{fppf}
$$
voir Champs algébriques, lemme
\ref{algebraic-lemma-representable-diagonal}. Par construction des produits
$2$-fibrés, un $T$-point de $Z$ correspond à un triplet
$(a, a', \alpha)$ formé de morphismes $a : T \to U$, $a' : T \to U'$ et
d'un morphisme $\alpha : a^*x \to (a')^*\xi'$. Nous obtenons un diagramme
commutatif
$$
\xymatrix{
U \ar[rd] \ar[rdd] \ar[rrd] \\
& Z \ar[r]_{p'} \ar[d]^p & U' \ar[d] \\
& U \ar[r] & S
}
$$
Le morphisme $i : U \to Z$ provient de l'isomorphisme $x \to \xi'|_U$.
Soit $z_0 = i(u_0) \in Z$. D'après le lemme
\ref{lemma-base-change-versal}, nous voyons que $Z \to U'$ est lisse en
$z_0$. En remplaçant $U$ par un voisinage ouvert affine de $u_0$, $U'$ par
l'ouvert correspondant et $Z$ par l'intersection des images réciproques de
ces ouverts par $p$ et $p'$, nous nous ramenons au cas où $Z \to U'$ est
lisse le long de $i(U)$. Remarquons que cela impose aussi de remplacer $u_n$
par une sous-suite, à savoir celle des indices pour lesquels $u_n$ appartient
à l'ouvert. En outre, la condition (3) du lemme
\ref{lemma-infinite-sequence-pre} est manifestement conservée lorsqu'on
rétrécit $U$ (tous les schémas $U$, $U_n$, $U'$ ont le même espace
topologique sous-jacent). Puisque $U \to U'$ est un épaississement du premier
ordre de schémas affines, nous pouvons choisir un morphisme $i' : U' \to Z$
tel que $p' \circ i' = \text{id}_{U'}$ et dont la restriction à $U$ soit $i$
(Compléments sur les morphismes d'espaces, lemme
\ref{spaces-more-morphisms-lemma-smooth-formally-smooth}).
En prenant l'image réciproque du morphisme universel
$p^*x \to (p')^*\xi'$ par $i'$, nous obtenons un morphisme
$$
\xi' \to x
$$
situé au-dessus de $p \circ i' : U' \to U$ et tel que la composée
$$
x \to \xi' \to x
$$
soit l'identité. Rappelons que nous avons $x_1 \to \xi'$ situé au-dessus du
morphisme $U_1 \to U'$. Par composition, nous obtenons un morphisme
$x_1 \to x$ dont l'existence contredit la condition (3) du lemme
\ref{lemma-infinite-sequence-pre}. Cette contradiction achève la démonstration.
\end{proof}

\begin{remark}
\label{remark-trade-openness-versality-diagonal-with-strong-effectiveness}
On peut déduire l'ouverture de la versalité de la diagonale d'une catégorie
fibrée en groupoïdes d'un axiome d'effectivité formelle forte. Soit $S$ un
schéma localement noethérien. Soit $\mathcal{X}$ une catégorie fibrée en
groupoïdes sur $(\Sch/S)_{fppf}$. Supposons que
\begin{enumerate}
\item $\Delta_\Delta : \mathcal{X} \to
\mathcal{X} \times_{\mathcal{X} \times \mathcal{X}} \mathcal{X}$
soit représentable par des espaces algébriques ;
\item $\mathcal{X}$ possède (RS*) ;
\item $\mathcal{X}$ commute aux limites ;
\item pour tout système projectif $(R_n)$ de $S$-algèbres comme dans la
remarque \ref{remark-strong-effectiveness}, où $\Ker(R_m \to R_n)$ est un
idéal de carré nul pour tous $m \geq n$, le foncteur
$$
\mathcal{X}_{\Spec(\lim R_n)} \longrightarrow
\lim_n \mathcal{X}_{\Spec(R_n)}
$$
soit pleinement fidèle.
\end{enumerate}
Alors $\Delta : \mathcal{X} \to \mathcal{X} \times \mathcal{X}$ satisfait
à l'ouverture de la versalité. Cela résulte de l'application du lemme
\ref{lemma-SGE-implies-openness-versality} aux produits fibrés de la forme
$\mathcal{X} \times_{\Delta, \mathcal{X} \times \mathcal{X}, y}
(\Sch/V)_{fppf}$ pour tout schéma affine $V$ localement de présentation finie
sur $S$ et tout objet $y$ de $\mathcal{X} \times \mathcal{X}$ au-dessus de
$V$. Si nous en avons un jour besoin, nous transformerons cette remarque en
lemme et en donnerons une démonstration détaillée.
\end{remark}











\section{Déformations infinitésimales}
\label{section-inf}

\noindent
Dans cette section, nous étudions une généralisation de la notion d'espace
tangent introduite dans la section \ref{section-tangent-spaces}. Pour le faire
de manière judicieuse, nous empruntons quelques notations à Théorie formelle
des déformations, sections
\ref{formal-defos-section-tangent-spaces-functors},
\ref{formal-defos-section-lifts}, et
\ref{formal-defos-section-infinitesimal-automorphisms}.

\medskip\noindent
Soit $S$ un schéma. Soit $\mathcal{X}$ une catégorie fibrée en groupoïdes sur
$(\Sch/S)_{fppf}$. Étant donnés un homomorphisme $A' \to A$ de $S$-algèbres
et un objet $x$ de $\mathcal{X}$ au-dessus de $\Spec(A)$, nous notons
$\textit{Lift}(x, A')$ la catégorie des relèvements de $x$ à $\Spec(A')$.
Un objet de $\textit{Lift}(x, A')$ est un morphisme $x \to x'$ de
$\mathcal{X}$ situé au-dessus de $\Spec(A) \to \Spec(A')$, et les morphismes
de $\textit{Lift}(x, A')$ sont définis comme des diagrammes commutatifs.
L'ensemble des classes d'isomorphisme de $\textit{Lift}(x, A')$ est noté
$\text{Lift}(x, A')$. Voir Théorie formelle des déformations, définition
\ref{formal-defos-definition-lifts} et remarque
\ref{formal-defos-remark-omit-arrow}. Si $A' \to A$ est surjectif et de noyau
localement nilpotent, nous appelons un élément $x'$ de
$\text{Lift}(x, A')$ une {\it déformation (infinitésimale)} de $x$. Dans ce
cas, le {\it groupe des automorphismes infinitésimaux de $x'$ au-dessus de
$x$} est le noyau
$$
\text{Inf}(x'/x) =
\Ker\left(
\text{Aut}_{\mathcal{X}_{\Spec(A')}}(x') \to
\text{Aut}_{\mathcal{X}_{\Spec(A)}}(x)\right)
$$
Remarquons qu'un élément de $\text{Inf}(x'/x)$ revient au même qu'un relèvement
de $\text{id}_x$ au-dessus de $\Spec(A')$ pour (la catégorie fibrée en
ensembles associée à) $\mathit{Aut}_\mathcal{X}(x')$. Comparer avec Théorie
formelle des déformations, définition
\ref{formal-defos-definition-relative-infinitesimal-auts} et Théorie formelle
des déformations, remarque
\ref{formal-defos-remark-infaut-lifting-equalities}.

\medskip\noindent
Si $M$ est un $A$-module, nous notons $A[M]$ la $A$-algèbre dont le
$A$-module sous-jacent est $A \oplus M$ et dont la multiplication est donnée
par $(a, m) \cdot (a', m') = (aa', am' + a'm)$. Lorsque $M = A$, c'est
l'algèbre des nombres duaux sur $A$, que nous notons, comme il est d'usage,
$A[\epsilon]$. Il existe un morphisme de $A$-algèbres $A[M] \to A$. L'image
réciproque de $x$ sur $\Spec(A[M])$ est appelée la {\it déformation triviale}
de $x$ à $\Spec(A[M])$.

\begin{lemma}
\label{lemma-functoriality}
Soit $S$ un schéma. Soit $f : \mathcal{X} \to \mathcal{Y}$ un $1$-morphisme
de catégories fibrées en groupoïdes sur $(\Sch/S)_{fppf}$. Soit
$$
\xymatrix{
B' \ar[r] & B \\
A' \ar[u] \ar[r] & A \ar[u]
}
$$
un diagramme commutatif de $S$-algèbres. Soit $x$ un objet de $\mathcal{X}$
au-dessus de $\Spec(A)$, soit $y$ un objet de $\mathcal{Y}$ au-dessus de
$\Spec(B)$, et soit $\phi : f(x)|_{\Spec(B)} \to y$ un morphisme de
$\mathcal{Y}$ au-dessus de $\Spec(B)$. Il existe alors un foncteur canonique
$$
\textit{Lift}(x, A') \longrightarrow \textit{Lift}(y, B')
$$
entre les catégories de relèvements, induit par $f$ et $\phi$. La construction
est, de manière évidente, compatible aux compositions de $1$-morphismes de
catégories fibrées en groupoïdes.
\end{lemma}

\begin{proof}
Le lemme se démontre de lui-même.
\end{proof}

\noindent
Soit $S$ un schéma de base. Soit $\mathcal{X}$ une catégorie fibrée en
groupoïdes sur $(\Sch/S)_{fppf}$. Nous définissons une catégorie dont les
objets sont les paires $(x, A' \to A)$ telles que
\begin{enumerate}
\item $A' \to A$ soit une surjection de $S$-algèbres dont le noyau est un
idéal de carré nul ;
\item $x$ soit un objet de $\mathcal{X}$ situé au-dessus de $\Spec(A)$.
\end{enumerate}
Un morphisme $(y, B' \to B) \to (x, A' \to A)$ est donné par un diagramme
commutatif
$$
\xymatrix{
B' \ar[r] & B \\
A' \ar[u] \ar[r] & A \ar[u]
}
$$
de $S$-algèbres et par un morphisme $x|_{\Spec(B)} \to y$ au-dessus de
$\Spec(B)$. Appelons-la la catégorie des {\it situations de déformation}.

\begin{lemma}
\label{lemma-properties-lift-RS-star}
Soit $S$ un schéma. Soit $\mathcal{X}$ une catégorie fibrée en groupoïdes sur
$(\Sch/S)_{fppf}$. Supposons que $\mathcal{X}$ satisfasse à la condition
(RS*). Soit $A$ une $S$-algèbre et soit $x$ un objet de $\mathcal{X}$
au-dessus de $\Spec(A)$.
\begin{enumerate}
\item Il existe un foncteur $A$-linéaire
$\text{Inf}_x : \text{Mod}_A \to \text{Mod}_A$
tel que, pour toute situation de déformation $(x, A' \to A)$ et tout
relèvement $x'$, il existe un isomorphisme
$\text{Inf}_x(I) \to \text{Inf}(x'/x)$, où $I = \Ker(A' \to A)$.
\item Il existe un foncteur $A$-linéaire
$T_x : \text{Mod}_A \to \text{Mod}_A$
tel que
\begin{enumerate}
\item pour tout $M$ dans $\text{Mod}_A$, il existe une bijection
$T_x(M) \to \text{Lift}(x, A[M])$ ;
\item pour toute situation de déformation $(x, A' \to A)$, il existe une
action
$$
T_x(I) \times \text{Lift}(x, A') \to \text{Lift}(x, A')
$$
où $I = \Ker(A' \to A)$. Elle est simplement transitive si
$\text{Lift}(x, A') \not = \emptyset$.
\end{enumerate}
\end{enumerate}
\end{lemma}

\begin{proof}
Nous définissons $\text{Inf}_x$ comme le foncteur
$$
\text{Mod}_A \longrightarrow \textit{Sets},\quad
M \longrightarrow
\text{Inf}(x'_M/x) = \text{Lift}(\text{id}_x, A[M])
$$
qui envoie $M$ sur le groupe des automorphismes infinitésimaux de la
déformation triviale $x'_M$ de $x$ à $\Spec(A[M])$, ou, de manière
équivalente, sur le groupe des relèvements de $\text{id}_x$ dans
$\mathit{Aut}_\mathcal{X}(x'_M)$. Nous définissons $T_x$ comme le foncteur
$$
\text{Mod}_A \longrightarrow \textit{Sets},\quad
M \longrightarrow \text{Lift}(x, A[M])
$$
des classes d'isomorphisme des déformations infinitésimales de $x$ à
$\Spec(A[M])$. Nous appliquons Théorie formelle des déformations, lemme
\ref{formal-defos-lemma-linear-functor}, à $\text{Inf}_x$ et à $T_x$. Ce lemme
s'applique, puisque (RS*) nous dit que
$$
\textit{Lift}(x, A[M \times N]) =
\textit{Lift}(x, A[M]) \times \textit{Lift}(x, A[N])
$$
comme catégories (et que les déformations triviales se correspondent aussi).

\medskip\noindent
Soit $(x, A' \to A)$ une situation de déformation. Considérons le morphisme
d'anneaux $g : A' \times_A A' \to A[I]$ défini par la règle
$g(a_1, a_2) = \overline{a_1} \oplus a_2 - a_1$. Il existe un isomorphisme
$$
A' \times_A A' \longrightarrow A' \times_A A[I]
$$
donné par $(a_1, a_2) \mapsto (a_1, g(a_1, a_2))$. Cet isomorphisme commute
aux projections sur le premier facteur $A'$, et donc aux projections sur
$A$. En appliquant ainsi (RS*) deux fois, nous obtenons des équivalences de
catégories
\begin{align*}
\textit{Lift}(x, A') \times \textit{Lift}(x, A')
& =
\textit{Lift}(x, A' \times_A A') \\
& =
\textit{Lift}(x, A' \times_A A[I]) \\
& =
\textit{Lift}(x, A') \times \textit{Lift}(x, A[I])
\end{align*}
En utilisant ces applications et la projection sur le dernier facteur du
dernier produit, nous voyons que nous obtenons des « applications différence »
$$
\text{Inf}(x'/x) \times  \text{Inf}(x'/x)
\longrightarrow
\text{Inf}_x(I)
\quad\text{et}\quad
\text{Lift}(x, A') \times \text{Lift}(x, A')
\longrightarrow
T_x(I)
$$
Ces applications différence satisfont à la règle de transitivité
« $(x'_1 - x'_2) + (x'_2 - x'_3) = x'_1 - x'_3$ », car
$$
\xymatrix{
A' \times_A A' \times_A A'
\ar[rrrrr]_-{(a_1, a_2, a_3) \mapsto (g(a_1, a_2), g(a_2, a_3))}
\ar[rrrrrd]_{(a_1, a_2, a_3) \mapsto g(a_1, a_3)} & & & & &
A[I] \times_A A[I] = A[I \times I] \ar[d]^{+} \\
& & & & & A[I]
}
$$
est commutatif. En inversant la suite d'équivalences ci-dessus, nous obtenons
une action libre et transitive pourvu que $\text{Inf}(x'/x)$, resp.\
$\text{Lift}(x, A')$, soit non vide. Remarquons que $\text{Inf}(x'/x)$ est
toujours non vide, puisque c'est un groupe.
\end{proof}

\begin{remark}[Fonctorialité]
\label{remark-functoriality}
Sous les hypothèses et avec les notations du lemme
\ref{lemma-properties-lift-RS-star}, supposons que $A \to B$ soit un morphisme
d'anneaux et que $y = x|_{\Spec(B)}$. Soient $M \in \text{Mod}_A$,
$N \in \text{Mod}_B$, et soit $M \to N$ une application $A$-linéaire. Il
existe alors des applications canoniques
$\text{Inf}_x(M) \to \text{Inf}_y(N)$ et $T_x(M) \to T_y(N)$, simplement
parce qu'il existe un foncteur d'image réciproque
$$
\textit{Lift}(x, A[M]) \to \textit{Lift}(y, B[N])
$$
provenant du morphisme d'anneaux $A[M] \to B[N]$. De même, étant donné un
morphisme de situations de déformation
$(y, B' \to B) \to (x, A' \to A)$, nous obtenons un foncteur d'image
réciproque $\textit{Lift}(x, A') \to \textit{Lift}(y, B')$. Puisque la
construction de l'action, de l'addition et de la multiplication scalaire sur
$\text{Inf}_x$ et $T_x$ n'utilise que des morphismes dans les catégories de
relèvements (voir la démonstration de Théorie formelle des déformations, lemme
\ref{formal-defos-lemma-linear-functor}), nous voyons que les constructions
ci-dessus sont fonctorielles. Autrement dit, nous obtenons des applications
$A$-linéaires
$$
\text{Inf}_x(M) \to \text{Inf}_y(N)
\quad\text{et}\quad
T_x(M) \to T_y(N)
$$
telles que les diagrammes
$$
\vcenter{
\xymatrix{
\text{Inf}_y(J) \ar[r] & \text{Inf}(y'/y) \\
\text{Inf}_x(I) \ar[r] \ar[u] & \text{Inf}(x'/x) \ar[u]
}
}
\quad\text{et}\quad
\vcenter{
\xymatrix{
T_y(J) \times \text{Lift}(y, B') \ar[r] & \text{Lift}(y, B') \\
T_x(I) \times \text{Lift}(x, A') \ar[r] \ar[u] & \text{Lift}(x, A') \ar[u]
}
}
$$
soient commutatifs. Ici $I = \Ker(A' \to A)$, $J = \Ker(B' \to B)$,
$x'$ est un relèvement de $x$ à $A'$ (qui peut ne pas toujours exister) et
$y' = x'|_{\Spec(B')}$.
\end{remark}

\begin{remark}[Automorphismes]
\label{remark-automorphisms}
Sous les hypothèses et avec les notations du lemme
\ref{lemma-properties-lift-RS-star}, soient $x', x''$ des relèvements de $x$
à $A'$. Nous avons alors une application de composition
$$
\text{Inf}(x'/x) \times
\Mor_{\textit{Lift}(x, A')}(x', x'') \times \text{Inf}(x''/x)
\longrightarrow
\Mor_{\textit{Lift}(x, A')}(x', x'').
$$
Puisque $\textit{Lift}(x, A')$ est un groupoïde, si
$\Mor_{\textit{Lift}(x, A')}(x', x'')$ est non vide, ceci définit une action
à gauche simplement transitive de $\text{Inf}(x'/x)$ sur
$\Mor_{\textit{Lift}(x, A')}(x', x'')$ et une action à droite simplement
transitive de $\text{Inf}(x''/x)$. Or le lemme dit que
$\text{Inf}(x'/x) = \text{Inf}_x(I) = \text{Inf}(x''/x)$. Nous affirmons que
les deux actions ainsi décrites coïncident par ces identifications. En effet,
ou bien $x' \not \cong x''$, auquel cas l'affirmation est claire, ou bien
$x' \cong x''$ et, dans ce cas, nous pouvons supposer que $x'' = x'$ ; le
résultat découle alors de la commutativité de $\text{Inf}(x'/x)$. En
particulier, nous obtenons une action bien définie
$$
\text{Inf}_x(I) \times \Mor_{\textit{Lift}(x, A')}(x', x'')
\longrightarrow
\Mor_{\textit{Lift}(x, A')}(x', x'')
$$
qui est simplement transitive dès que
$\Mor_{\textit{Lift}(x, A')}(x', x'')$ est non vide.
\end{remark}

\begin{remark}
\label{remark-short-exact-sequence-thickenings}
Soit $S$ un schéma. Soit $\mathcal{X}$ une catégorie fibrée en groupoïdes sur
$(\Sch/S)_{fppf}$. Soit $A$ une $S$-algèbre. On dispose d'une notion de
{\it suite exacte courte}
$$
(x, A_1' \to A) \to (x, A_2' \to A) \to (x, A_3' \to A)
$$
de situations de déformation : nous demandons que les applications
correspondantes entre les noyaux $I_i = \Ker(A_i' \to A)$ donnent une suite
exacte courte
$$
0 \to I_3 \to I_2 \to I_1 \to 0
$$
de $A$-modules. Remarquons que, dans ce cas, l'application $A_3' \to A_1'$ se
factorise par $A$ ; il existe donc un isomorphisme canonique
$A_1' = A[I_1]$.
\end{remark}

\begin{lemma}
\label{lemma-ses-inf-and-T}
Soit $S$ un schéma. Soient $p : \mathcal{X} \to \mathcal{Y}$ et
$q : \mathcal{Z} \to \mathcal{Y}$ des $1$-morphismes de catégories fibrées en
groupoïdes sur $(\Sch/S)_{fppf}$. Supposons que $\mathcal{X}$,
$\mathcal{Y}$, $\mathcal{Z}$ satisfassent à (RS*). Soit $A$ une $S$-algèbre
et soit $w$ un objet de
$\mathcal{W} = \mathcal{X} \times_\mathcal{Y} \mathcal{Z}$ au-dessus de $A$.
Notons $x, y, z$ les objets de $\mathcal{X}, \mathcal{Y}, \mathcal{Z}$ que
l'on déduit de $w$. Pour tout $A$-module $M$, il existe une suite exacte à
$6$ termes
$$
\xymatrix{
0 \ar[r] &
\text{Inf}_w(M) \ar[r] &
\text{Inf}_x(M) \oplus \text{Inf}_z(M) \ar[r] &
\text{Inf}_y(M) \ar[lld] \\
 &
T_w(M) \ar[r] &
T_x(M) \oplus T_z(M) \ar[r] &
T_y(M)
}
$$
de $A$-modules.
\end{lemma}

\begin{proof}
D'après le lemme \ref{lemma-fibre-product-RS-star}, nous voyons que
$\mathcal{W}$ satisfait à (RS*) ; ainsi, $T_w(M)$ et $\text{Inf}_w(M)$ sont
définis. Les flèches horizontales sont définies à l'aide de la fonctorialité
du lemme \ref{lemma-functoriality}.

\medskip\noindent
Définition de l'application « bord »
$\delta : \text{Inf}_y(M) \to T_w(M)$. Choisissons des isomorphismes
$p(x) \to y$ et $y \to q(z)$ tels que
$w = (x, z, p(x) \to y \to q(z))$ dans la description du produit $2$-fibré
de Catégories, lemme \ref{categories-lemma-2-product-fibred-categories}, et,
plus précisément, Catégories, lemme
\ref{categories-lemma-2-product-categories-over-C}. Notons
$x', y', z', w'$ les déformations triviales de $x, y, z, w$ au-dessus de
$A[M]$. Par image réciproque, nous obtenons des isomorphismes $y' \to p(x')$
et $q(z') \to y'$. Un élément $\alpha \in \text{Inf}_y(M)$ revient au même
qu'un automorphisme $\alpha : y' \to y'$ au-dessus de $A[M]$ dont la
restriction à $y$ au-dessus de $A$ est l'identité. Ainsi, en posant
$$
\delta(\alpha) =
(x', z', p(x') \to y' \xrightarrow{\alpha} y' \to q(z'))
$$
on obtient un objet de $T_w(M)$. C'est une application de $A$-modules d'après
Théorie formelle des déformations, lemme
\ref{formal-defos-lemma-morphism-linear-functors}.

\medskip\noindent
La suite de la démonstration est exactement la même que celle de Théorie
formelle des déformations, lemme
\ref{formal-defos-lemma-deformation-categories-fiber-product-morphisms}.
\end{proof}

\begin{remark}[Compatibilité avec les espaces tangents précédents]
\label{remark-compare-deformation-spaces}
Soit $S$ un schéma localement noethérien. Soit $\mathcal{X}$ une catégorie
fibrée en groupoïdes sur $(\Sch/S)_{fppf}$. Supposons que $\mathcal{X}$
possède (RS*). Soit $k$ un corps de type fini sur $S$ et soit $x_0$ un objet
de $\mathcal{X}$ au-dessus de $\Spec(k)$. Nous avons alors des égalités
d'espaces vectoriels sur $k$
$$
T\mathcal{F}_{\mathcal{X}, k, x_0} = T_{x_0}(k)
\quad\text{et}\quad
\text{Inf}(\mathcal{F}_{\mathcal{X}, k, x_0}) =
\text{Inf}_{x_0}(k)
$$
où les espaces situés à gauche des signes d'égalité sont donnés dans
(\ref{equation-tangent-space}) et
(\ref{equation-infinitesimal-automorphisms}), tandis que les espaces situés à
droite sont donnés par le lemme \ref{lemma-properties-lift-RS-star}.
\end{remark}

\begin{remark}[Élément canonique]
\label{remark-canonical-element}
Sous les hypothèses et avec les notations du lemme
\ref{lemma-properties-lift-RS-star}, choisissons un ouvert affine
$\Spec(\Lambda) \subset S$ tel que $\Spec(A) \to S$ corresponde à un morphisme
d'anneaux $\Lambda \to A$. Considérons le morphisme d'anneaux
$$
A \longrightarrow A[\Omega_{A/\Lambda}],
\quad
a \longmapsto (a, \text{d}_{A/\Lambda}(a))
$$
En prenant l'image réciproque de $x$ par le morphisme correspondant
$\Spec(A[\Omega_{A/\Lambda}]) \to \Spec(A)$, nous obtenons une déformation
$x_{can}$ de $x$ au-dessus de $A[\Omega_{A/\Lambda}]$. Nous l'appelons
l'{\it élément canonique}
$$
x_{can} \in T_x(\Omega_{A/\Lambda}) = \text{Lift}(x, A[\Omega_{A/\Lambda}]).
$$
Supposons ensuite que $\Lambda$ soit noethérien et que $\Lambda \to A$ soit
de type fini. Soit $k = \kappa(\mathfrak p)$ le corps résiduel en un point de
type fini $u_0$ de $U = \Spec(A)$. Soit $x_0 = x|_{u_0}$. Par (RS*) et le fait
que $A[k] = A \times_k k[k]$, l'espace $T_x(k)$ est l'espace tangent au
foncteur de déformations $\mathcal{F}_{\mathcal{X}, k, x_0}$. Par
$$
T\mathcal{F}_{U, k, u_0} =
\text{Der}_\Lambda(A, k) = \Hom_A(\Omega_{A/\Lambda}, k)
$$
(voir Théorie formelle des déformations, exemple
\ref{formal-defos-example-tangent-space-prorepresentable-functor}) et par la
fonctorialité de $T_x$, l'élément canonique produit l'application sur les
espaces tangents induite par l'objet $x$ au-dessus de $U$. Plus précisément,
$\theta \in T\mathcal{F}_{U, k, u_0}$ est envoyé sur
$T_x(\theta)(x_{can})$ dans
$T_x(k) = T\mathcal{F}_{\mathcal{X}, k, x_0}$.
\end{remark}

\begin{remark}[Automorphisme canonique]
\label{remark-canonical-isomorphism}
Soit $S$ un schéma localement noethérien. Soit $\mathcal{X}$ une catégorie
fibrée en groupoïdes sur $(\Sch/S)_{fppf}$. Supposons que $\mathcal{X}$
satisfasse à la condition (RS*). Soit $A$ une $S$-algèbre telle que
$\Spec(A) \to S$ s'envoie dans un ouvert affine, et soient $x, y$ des objets
de $\mathcal{X}$ au-dessus de $\Spec(A)$. Soit en outre $A \to B$ un morphisme
d'anneaux et soit $\alpha : x|_{\Spec(B)} \to y|_{\Spec(B)}$ un morphisme de
$\mathcal{X}$ au-dessus de $\Spec(B)$. Considérons le morphisme d'anneaux
$$
B \longrightarrow B[\Omega_{B/A}],
\quad
b \longmapsto (b, \text{d}_{B/A}(b))
$$
En prenant l'image réciproque de $\alpha$ par le morphisme correspondant
$\Spec(B[\Omega_{B/A}]) \to \Spec(B)$, nous obtenons un morphisme
$\alpha_{can}$ entre les images réciproques de $x$ et de $y$ au-dessus de
$B[\Omega_{B/A}]$. D'autre part, nous pouvons prendre l'image réciproque de
$\alpha$ par le morphisme $\Spec(B[\Omega_{B/A}]) \to \Spec(B)$ correspondant
à l'injection de $B$ dans le premier facteur de $B[\Omega_{B/A}]$. D'après la
discussion de la remarque \ref{remark-automorphisms}, nous pouvons prendre la
différence
$$
\varphi(x, y, \alpha) = \alpha_{can} - \alpha|_{\Spec(B[\Omega_{B/A}])} \in
\text{Inf}_{x|_{\Spec(B)}}(\Omega_{B/A}).
$$
Nous l'appellerons l'{\it automorphisme canonique}. Il dépend de toutes les
données $A$, $x$, $y$, $A \to B$ et $\alpha$.
\end{remark}





\section{Théories des obstructions}
\label{section-obstruction-theory}

\noindent
Dans cette section, nous décrivons ce qu'est une théorie des obstructions.
Contrairement aux espaces de déformations infinitésimales et d'automorphismes
infinitésimaux, une théorie des obstructions constitue une donnée
supplémentaire. La formulation est motivée par les résultats du lemme
\ref{lemma-properties-lift-RS-star} et de la remarque
\ref{remark-functoriality}.

\begin{definition}
\label{definition-obstruction-theory}
Soit $S$ une base localement noethérienne. Soit $\mathcal{X}$ une catégorie
fibrée en groupoïdes sur $(\Sch/S)_{fppf}$. Une {\it théorie des obstructions}
est donnée par les données suivantes :
\begin{enumerate}
\item pour toute $S$-algèbre $A$ telle que $\Spec(A) \to S$ s'envoie dans
un ouvert affine et pour tout objet $x$ de $\mathcal{X}$ au-dessus de
$\Spec(A)$, un foncteur $A$-linéaire
$$
\mathcal{O}_x : \text{Mod}_A \to \text{Mod}_A
$$
de {\it modules d'obstruction} ;
\item pour $(x, A)$ comme en (1), un morphisme d'anneaux $A \to B$, des
éléments $M \in \text{Mod}_A$, $N \in \text{Mod}_B$ et une application
$A$-linéaire $M \to N$, une application $A$-linéaire induite
$\mathcal{O}_x(M) \to \mathcal{O}_y(N)$, où $y = x|_{\Spec(B)}$ ;
\item pour toute situation de déformation $(x, A' \to A)$, un élément
{\it obstruction} $o_x(A') \in \mathcal{O}_x(I)$, où
$I = \Ker(A' \to A)$.
\end{enumerate}
Ces données sont soumises aux conditions suivantes :
\begin{enumerate}
\item[(i)] les applications de fonctorialité font des modules d'obstruction
un foncteur de la catégorie des triplets $(x, A, M)$ vers les ensembles ;
\item[(ii)] pour tout morphisme de situations de déformation
$(y, B' \to B) \to (x, A' \to A)$, l'élément $o_x(A')$ s'envoie sur
$o_y(B')$ ;
\item[(iii)] on a
$$
\text{Lift}(x, A') \not = \emptyset
\Leftrightarrow
o_x(A') = 0
$$
pour toute situation de déformation $(x, A' \to A)$.
\end{enumerate}
\end{definition}

\noindent
Cette dernière condition explique la terminologie. Le module
$\mathcal{O}_x(A')$ est appelé le {\it module d'obstruction}. L'élément
$o_x(A')$ est l'{\it obstruction}. La plupart des théories des obstructions
possèdent des propriétés supplémentaires et, pour les rendre utiles, des
conditions supplémentaires sont nécessaires. De plus, il ne s'agit que d'un
exemple de définition : nous pourrions, par exemple, n'y considérer que les
situations de déformation de type fini sur $S$.

\medskip\noindent
L'une des principales raisons d'introduire les théories des obstructions est
de vérifier l'ouverture de la versalité. Le lemme
\ref{lemma-get-openness-obstruction-theory} ci-dessous fournit un exemple de
ce type de résultat. L'idée initiale de cette méthode est due à Artin ; voir
les articles d'Artin mentionnés dans l'introduction. Elle a été reprise, par
exemple, dans les travaux de Flenner \cite{Flenner}, Hall
\cite{Hall-coherent}, Hall et Rydh \cite{rydh_axioms}, Olsson
\cite{olsson_deformation}, Olsson et Starr \cite{olsson-starr}, et Lieblich
\cite{lieblich-complexes} (ordre aléatoire des références). En outre, pour des
catégories fibrées en groupoïdes particulières, les auteurs développent
souvent un peu de théorie adaptée au problème considéré. Nous développerons
cette théorie plus loin (insérer ici une référence future).

\begin{lemma}
\label{lemma-get-openness-obstruction-theory}
\begin{reference}
C'est \cite[Theorem 4.4]{Hall-coherent}.
\end{reference}
Soit $S$ un schéma localement noethérien. Soit $\mathcal{X}$ une catégorie
fibrée en groupoïdes sur $(\Sch/S)_{fppf}$. Supposons que
\begin{enumerate}
\item $\Delta : \mathcal{X} \to \mathcal{X} \times \mathcal{X}$ soit
représentable par des espaces algébriques ;
\item $\mathcal{X}$ possède (RS*) ;
\item $\mathcal{X}$ commute aux limites ;
\item il existe une théorie des obstructions\footnote{En analysant la
démonstration, le lecteur voit qu'il suffit en fait de vérifier la
fonctorialité (ii) des classes d'obstruction de la définition
\ref{definition-obstruction-theory} pour les applications
$(y, B' \to B) \to (x, A' \to A)$ avec $B = A$ et $y = x$.} ;
\item pour un objet $x$ de $\mathcal{X}$ au-dessus de $\Spec(A)$ et des
$A$-modules $M_n$, $n \geq 1$, on ait
\begin{enumerate}
\item $T_x(\prod M_n) = \prod T_x(M_n)$,
\item $\mathcal{O}_x(\prod M_n) \to \prod \mathcal{O}_x(M_n)$
soit injective.
\end{enumerate}
\end{enumerate}
Alors $\mathcal{X}$ satisfait à l'ouverture de la versalité.
\end{lemma}

\begin{proof}
Nous le démontrons en vérifiant la condition (4) du lemme
\ref{lemma-SGE-implies-openness-versality}. Soient $(\xi_n)$ et $(R_n)$ comme
dans la remarque \ref{remark-strong-effectiveness}, tels que
$\Ker(R_m \to R_n)$ soit un idéal de carré nul pour tous $m \geq n$. Posons
$A = R_1$ et $x = \xi_1$. Notons $M_n = \Ker(R_n \to R_1)$. Alors $M_n$ est
un $A$-module. Posons $R = \lim R_n$. Soit
$$
\tilde R = \{(r_1, r_2, r_3 \ldots) \in \prod R_n
\text{ tels que tous aient la même image dans }A\}
$$
Alors $\tilde R \to A$ est surjectif, de noyau $M = \prod M_n$. Il existe une
application $R \to \tilde R$ et une application $\tilde R \to A[M]$,
$(r_1, r_2, r_3, \ldots) \mapsto
(r_1, r_2 - r_1, r_3 - r_2, \ldots)$. Ensemble, elles donnent une suite
exacte courte
$$
(x, R \to A) \to (x, \tilde R \to A) \to (x, A[M])
$$
de situations de déformation ; voir la remarque
\ref{remark-short-exact-sequence-thickenings}. La suite de noyaux associée
$0 \to \lim M_n \to M \to M \to 0$ est la suite canonique qui calcule la
limite du système de modules $(M_n)$.

\medskip\noindent
Soit $o_x(\tilde R) \in \mathcal{O}_x(M)$ l'élément d'obstruction. Puisque nous
disposons des relèvements $\xi_n$, nous voyons que $o_x(\tilde R)$ s'envoie
sur zéro dans $\mathcal{O}_x(M_n)$. Par l'hypothèse (5)(b), nous voyons que
$o_x(\tilde R) = 0$. Choisissons un relèvement $\tilde \xi$ de $x$ à
$\Spec(\tilde R)$. Soit $\tilde \xi_n$ la restriction de $\tilde \xi$ à
$\Spec(R_n)$. Il existe des éléments $t_n \in T_x(M_n)$ tels que
$t_n \cdot \tilde \xi_n = \xi_n$, d'après la partie (2)(b) du lemme
\ref{lemma-properties-lift-RS-star}. Par l'hypothèse (5)(a), nous pouvons
trouver $t \in T_x(M)$ qui s'envoie sur $t_n$ dans $T_x(M_n)$. Après avoir
remplacé $\tilde \xi$ par $t \cdot \tilde \xi$, nous constatons que
$\tilde \xi$ se restreint à $\xi_n$ au-dessus de $\Spec(R_n)$ pour tout $n$.
En particulier, puisque $\xi_{n + 1}$ se restreint à $\xi_n$ au-dessus de
$\Spec(R_n)$, la restriction $\overline{\xi}$ de $\tilde \xi$ à
$\Spec(A[M])$ a la propriété que sa restriction est la déformation triviale
au-dessus de $\Spec(A[M_n])$ pour tout $n$. Par conséquent, l'hypothèse (5)(a)
montre que $\overline{\xi}$ est la déformation triviale de $x$. En appliquant
l'axiome (RS*) à $R = \tilde R \times_{A[M]} A$, il s'ensuit que
$\tilde \xi$ est l'image réciproque d'une déformation $\xi$ de $x$ au-dessus
de $R$. Ceci achève la démonstration.
\end{proof}

\begin{example}
\label{example-global-sections}
Soit $S = \Spec(\Lambda)$ pour un anneau noethérien $\Lambda$. Soit
$W \to S$ un morphisme de schémas. Soit $\mathcal{F}$ un
$\mathcal{O}_W$-module quasi-cohérent et plat sur $S$. Considérons le foncteur
$$
F : (\Sch/S)_{fppf}^{opp} \longrightarrow \textit{Sets},
\quad
T/S \longrightarrow H^0(W_T, \mathcal{F}_T)
$$
où $W_T = T \times_S W$ est le changement de base et $\mathcal{F}_T$ l'image
réciproque de $\mathcal{F}$ sur $W_T$. Si $T = \Spec(A)$, nous écrirons
$W_T = W_A$, etc. Soit $\mathcal{X} \to (\Sch/S)_{fppf}$ la catégorie fibrée
en groupoïdes associée à $F$. Alors $\mathcal{X}$ possède une théorie des
obstructions. Plus précisément,
\begin{enumerate}
\item étant donnés $A$ sur $\Lambda$ et $x \in H^0(W_A, \mathcal{F}_A)$, nous
posons $\mathcal{O}_x(M) = H^1(W_A, \mathcal{F}_A \otimes_A M)$ ;
\item étant donnée une situation de déformation $(x, A' \to A)$, nous prenons
pour $o_x(A') \in \mathcal{O}_x(A)$ l'image de $x$ par l'application bord
$$
H^0(W_A, \mathcal{F}_A) \longrightarrow H^1(W_A, \mathcal{F}_A \otimes_A I)
$$
provenant de la suite exacte courte de modules
$$
0 \to \mathcal{F}_A \otimes_A I \to
\mathcal{F}_{A'} \to \mathcal{F}_A \to 0.
$$
\end{enumerate}
Nous avons omis certains détails, notamment la construction de la suite exacte
courte ci-dessus (elle utilise le fait que $W_A$ et $W_{A'}$ ont le même
espace topologique sous-jacent) et l'explication de la raison pour laquelle la
platitude de $\mathcal{F}$ sur $S$ implique que la suite ci-dessus est exacte
courte.
\end{example}

\begin{example}[Exemple clé]
\label{example-key}
Soit $S = \Spec(\Lambda)$ pour un anneau noethérien $\Lambda$. Prenons
$\mathcal{X} = (\Sch/X)_{fppf}$ avec $X = \Spec(R)$ et
$R = \Lambda[x_1, \ldots, x_n]/J$. Le complexe cotangent naïf
$\NL_{R/\Lambda}$ est (canoniquement) homotopiquement équivalent à
$$
J/J^2
\longrightarrow
\bigoplus\nolimits_{i = 1, \ldots, n} R\text{d}x_i,
$$
voir Algèbre, lemme \ref{algebra-lemma-NL-homotopy}. Considérons une situation
de déformation $(x, A' \to A)$. Notons $I$ le noyau de $A' \to A$. L'objet
$x$ correspond à $(a_1, \ldots, a_n)$ avec $a_i \in A$ tel que
$f(a_1, \ldots, a_n) = 0$ dans $A$ pour tout $f \in J$. Posons
\begin{align*}
\mathcal{O}_x(A')
& =
\Hom_R(J/J^2, I)/\Hom_R(R^{\oplus n}, I) \\
& =
\Ext^1_R(\NL_{R/\Lambda}, I) \\
& =
\Ext^1_A(\NL_{R/\Lambda} \otimes_R A, I).
\end{align*}
Choisissons des relèvements $a_i' \in A'$ de $a_i$ dans $A$. Alors
$o_x(A')$ est la classe de l'application $J/J^2 \to I$ définie en envoyant
$f \in J$ sur $f(a_1', \ldots, a'_n) \in I$. Nous omettons de vérifier que
$o_x(A')$ est indépendant des choix. Il est clair que, si $o_x(A') = 0$,
alors l'application se relève. Enfin, la fonctorialité est immédiate. Nous
obtenons ainsi une théorie des obstructions. Remarquons que $o_x(A')$ peut se
décrire de manière un peu plus canonique comme la composée
$$
\NL_{R/\Lambda} \to \NL_{A/\Lambda} \to \NL_{A/A'} = I[1]
$$
dans $D(A)$ ; voir Algèbre, lemme \ref{algebra-lemma-NL-surjection}, pour la
dernière identification.
\end{example}









\section{Théories naïves des obstructions}
\label{section-naive-obstruction-theory}

\noindent
Le titre de cette section renvoie au fait que nous y utiliserons le complexe
cotangent naïf. Soit $(x, A' \to A)$ une situation de déformation pour une
catégorie fibrée en groupoïdes donnée sur un schéma localement noethérien $S$.
L'exemple clé \ref{example-key} suggère que toute théorie des obstructions
devrait être étroitement liée aux morphismes de $D(A)$ ayant pour but le
complexe cotangent naïf de $A$. En explicitant cette idée, nous obtenons au
lemme \ref{lemma-characterize-versal} un critère de versalité qui conduit, au
lemme \ref{lemma-openness}, à un critère d'ouverture de la versalité. Pour
formaliser un peu davantage cette notion, nous introduisons une notion de
théorie naïve des obstructions dans la définition
\ref{definition-naive-obstruction-theory}.

\medskip\noindent
Dans la suite, nous utiliserons le complexe cotangent naïf tel qu'il est défini
dans Algèbre, section \ref{algebra-section-netherlander}. En particulier, si
$A' \to A$ est une surjection de $\Lambda$-algèbres de noyau de carré nul $I$,
alors il existe des morphismes
$$
\NL_{A'/\Lambda} \to \NL_{A/\Lambda} \to \NL_{A/A'}
$$
dont la composée est homotopiquement équivalente à zéro (voir Algèbre,
remarque \ref{algebra-remark-composition-homotopy-equivalent-to-zero}). En
général, ils ne forment pas un triangle distingué, car nous utilisons le
complexe cotangent naïf et non le complexe complet. Il existe une équivalence
d'homotopie $\NL_{A/A'} \to I[1]$ (le complexe formé de $I$ placé en degré
$-1$ ; voir Algèbre, lemme \ref{algebra-lemma-NL-surjection}). Enfin,
remarquons qu'il existe un morphisme canonique
$\NL_{A/\Lambda} \to \Omega_{A/\Lambda}$.

\begin{lemma}
\label{lemma-compute-ext-into-field}
Soit $A \to k$ un morphisme d'anneaux, où $k$ est un corps. Soit
$E \in D^-(A)$. Alors
$\Ext^i_A(E, k) = \Hom_k(H^{-i}(E \otimes^\mathbf{L} k), k)$.
\end{lemma}

\begin{proof}
Omis. Indication : remplacer $E$ par un complexe borné supérieurement de
$A$-modules libres et calculer les deux membres.
\end{proof}

\begin{lemma}
\label{lemma-construct-essential-surjection}
Soient $\Lambda \to A \to k$ des morphismes de type fini d'anneaux
noethériens, avec $k = \kappa(\mathfrak p)$ pour un idéal premier
$\mathfrak p$ de $A$. Soit $\xi : E \to \NL_{A/\Lambda}$ un morphisme de
$D^{-}(A)$ tel que $H^{-1}(\xi \otimes^{\mathbf{L}} k)$ ne soit pas surjectif.
Il existe alors une surjection $A' \to A$ de $\Lambda$-algèbres telle que
\begin{enumerate}
\item[(a)] $I = \Ker(A' \to A)$ soit de carré nul et isomorphe à $k$ comme
$A$-module ;
\item[(b)] $\Omega_{A'/\Lambda} \otimes k = \Omega_{A/\Lambda} \otimes k$ ;
\item[(c)] $E \to \NL_{A/A'}$ soit nul.
\end{enumerate}
\end{lemma}

\begin{proof}
Soient $f \in A$, $f \not \in \mathfrak p$. Supposons que $A'' \to A_f$
satisfasse à (a), (b), (c) pour le morphisme induit
$E \otimes_A A_f \to \NL_{A_f/\Lambda}$ ; voir Algèbre, lemme
\ref{algebra-lemma-localize-NL}. Nous pouvons alors poser
$A' = A'' \times_{A_f} A$ et obtenir une solution. En effet, il est clair que
$A' \to A$ satisfait à (a), puisque
$\Ker(A' \to A) = \Ker(A'' \to A) = I$. Choisissons $f'' \in A''$ relevant
$f$. Alors le localisé de $A'$ en $(f'', f)$ est isomorphe à $A''$ (par
exemple d'après Compléments sur l'algèbre, lemme
\ref{more-algebra-lemma-diagram-localize}). Ainsi, (b) et (c) sont également
clairs pour $A'$. Nous voyons de cette manière que nous pouvons remplacer $A$
par le localisé $A_f$ (un nombre fini de fois). En particulier, après un tel
remplacement, nous pouvons supposer que $\mathfrak p$ est un idéal maximal de
$A$ ; voir Morphismes, lemme \ref{morphisms-lemma-point-finite-type}.

\medskip\noindent
Choisissons une présentation $A = \Lambda[x_1, \ldots, x_n]/J$. Alors
$\NL_{A/\Lambda}$ est (canoniquement) homotopiquement équivalent à
$$
J/J^2
\longrightarrow
\bigoplus\nolimits_{i = 1, \ldots, n} A\text{d}x_i,
$$
voir Algèbre, lemme \ref{algebra-lemma-NL-homotopy}. Après localisation si
nécessaire (en utilisant le lemme de Nakayama), nous pouvons choisir des
générateurs $f_1, \ldots, f_m$ de $J$ tels que les $f_j \otimes 1$ forment une
base de $J/J^2 \otimes_A k$. De plus, après renumérotation, nous pouvons
supposer que les images de $\text{d}f_1, \ldots, \text{d}f_r$ forment une
base de l'image de $J/J^2 \otimes k \to \bigoplus k\text{d}x_i$ et que
$\text{d}f_{r + 1}, \ldots, \text{d}f_m$ s'envoient sur zéro dans
$\bigoplus k\text{d}x_i$. Avec ces choix, l'espace
$$
H^{-1}(\NL_{A/\Lambda} \otimes^{\mathbf{L}}_A k) =
H^{-1}(\NL_{A/\Lambda} \otimes_A k)
$$
a pour base $f_{r + 1} \otimes 1, \ldots, f_m \otimes 1$. En changeant encore
une fois de base, nous pouvons supposer que l'image de
$H^{-1}(\xi \otimes^{\mathbf{L}} k)$ est contenue dans le sous-espace engendré
sur $k$ par $f_{r + 1} \otimes 1, \ldots, f_{m - 1} \otimes 1$. Posons
$$
A' = \Lambda[x_1, \ldots, x_n]/(f_1, \ldots, f_{m - 1}, \mathfrak pf_m)
$$
Par construction, $A' \to A$ satisfait à (a). Puisque $\text{d}f_m$ s'envoie
sur zéro dans $\bigoplus k\text{d}x_i$, nous voyons que (b) est satisfaite.
Enfin, par construction, le morphisme induit $E \to \NL_{A/A'} = I[1]$
induit l'application nulle
$H^{-1}(E \otimes_A^\mathbf{L} k) \to I \otimes_A k$. D'après le lemme
\ref{lemma-compute-ext-into-field}, nous voyons que la composée est nulle.
\end{proof}

\noindent
Le lemme suivant est notre résultat technique clé.

\begin{lemma}
\label{lemma-characterize-versal}
Soit $S$ un schéma localement noethérien. Soit $\mathcal{X}$ une catégorie
fibrée en groupoïdes sur $(\Sch/S)_{fppf}$ satisfaisant à (RS*). Soit
$U = \Spec(A)$ un schéma affine de type fini sur $S$ qui s'envoie dans un
ouvert affine $\Spec(\Lambda)$. Soit $x$ un objet de $\mathcal{X}$ au-dessus
de $U$. Soit $\xi : E \to \NL_{A/\Lambda}$ un morphisme de $D^{-}(A)$.
Supposons que
\begin{enumerate}
\item[(i)] pour toute situation de déformation $(x, A' \to A)$, l'objet $x$
se relève à $\Spec(A')$ si et seulement si
$E \to \NL_{A/\Lambda} \to \NL_{A/A'}$ est nul ;
\item[(ii)] il existe un isomorphisme de foncteurs
$T_x(-) \to \Ext^0_A(E, -)$
tel que $E \to \NL_{A/\Lambda} \to \Omega^1_{A/\Lambda}$ corresponde à
l'élément canonique (voir la remarque \ref{remark-canonical-element}).
\end{enumerate}
Soit $u_0 \in U$ un point de type fini de corps résiduel $k = \kappa(u_0)$.
Considérons les assertions suivantes :
\begin{enumerate}
\item $x$ est versel en $u_0$ ;
\item $\xi : E \to \NL_{A/\Lambda}$ induit une surjection
$H^{-1}(E \otimes_A^{\mathbf{L}} k) \to
H^{-1}(\NL_{A/\Lambda} \otimes_A^{\mathbf{L}} k)$
et une injection
$H^0(E \otimes_A^{\mathbf{L}} k) \to
H^0(\NL_{A/\Lambda} \otimes_A^{\mathbf{L}} k)$.
\end{enumerate}
Alors on a toujours (2) $\Rightarrow$ (1), et on a (1) $\Rightarrow$ (2) si
$u_0$ est un point fermé.
\end{lemma}

\begin{proof}
Soit $\mathfrak p = \Ker(A \to k)$ l'idéal premier correspondant à $u_0$.

\medskip\noindent
Supposons que $x$ soit versel en $u_0$ et que $u_0$ soit un point fermé de
$U$. Si $H^{-1}(\xi \otimes_A^{\mathbf{L}} k)$ n'est pas surjectif, soit
$A' \to A$ une extension de noyau $I$ comme dans le lemme
\ref{lemma-construct-essential-surjection}. Puisque $u_0$ est un point fermé,
nous voyons que $I$ est un $A$-module fini, et donc que $A'$ est une
$\Lambda$-algèbre de type fini (ceci est faux si $u_0$ n'est pas fermé). En
particulier, $A'$ est noethérien. D'après la propriété (c) pour $A'$ et (i)
pour $\xi$, nous voyons que $x$ se relève en un objet $x'$ au-dessus de $A'$.
Soit $\mathfrak p' \subset A'$ le noyau de l'application surjective vers $k$.
D'après Artin--Rees (Algèbre, lemme \ref{algebra-lemma-Artin-Rees}), il existe
$n > 1$ tel que $(\mathfrak p')^n \cap I = 0$. Nous voyons alors que
$$
B' = A'/(\mathfrak p')^n \longrightarrow A/\mathfrak p^n = B
$$
est une petite extension essentielle d'anneaux locaux artiniens ; voir Théorie
formelle des déformations, lemme
\ref{formal-defos-lemma-essential-surjection}. D'autre part, puisque $x$ est
versel en $u_0$ et que $x'|_{\Spec(B')}$ est un relèvement de
$x|_{\Spec(B)}$, il existe un entier $m \geq n$ et une application
$q : A/\mathfrak p^m \to B'$ tels que la composée
$A/\mathfrak p^m \to B' \to B$ soit l'application quotient. Puisque
l'idéal maximal de $B'$ a sa puissance $n$-ième nulle, $q$ se factorise par
$B$, ce qui contredit le fait que $B' \to B$ soit une surjection essentielle.
Cette contradiction montre que $H^{-1}(\xi \otimes_A^{\mathbf{L}} k)$ est
surjectif.

\medskip\noindent
Supposons que $x$ soit versel en $u_0$. D'après le lemme
\ref{lemma-compute-ext-into-field}, l'application
$H^0(\xi \otimes_A^{\mathbf{L}} k)$ est duale de l'application
$\Ext^0_A(\NL_{A/\Lambda}, k) \to \text{Ext}^0_A(E, k)$. Remarquons que
$$
\Ext^0_A(\NL_{A/\Lambda}, k) = \text{Der}_\Lambda(A, k)
\quad\text{et}\quad
T_x(k) = \Ext^0_A(E, k)
$$
La condition (ii) assure que l'application
$\Ext^0_A(\NL_{A/\Lambda}, k) \to \text{Ext}^0_A(E, k)$ envoie un vecteur
tangent $\theta$ à $U$ en $u_0$ sur la déformation infinitésimale
correspondante de $x_0$ ; voir la remarque \ref{remark-canonical-element}.
Ainsi, si $x$ est versel, cette application est surjective ; voir Théorie
formelle des déformations, lemme \ref{formal-defos-lemma-versal-criterion}.
Par conséquent, $H^0(\xi \otimes_A^{\mathbf{L}} k)$ est injectif. Ceci achève
la démonstration de (1) $\Rightarrow$ (2) lorsque $u_0$ est un point fermé.

\medskip\noindent
Pour la suite de la démonstration, supposons que
$H^{-1}(E \otimes_A^\mathbf{L} k) \to
H^{-1}(\NL_{A/\Lambda} \otimes_A^\mathbf{L} k)$ soit surjectif et que
$H^0(E \otimes_A^\mathbf{L} k) \to
H^0(\NL_{A/\Lambda} \otimes_A^\mathbf{L} k)$
soit injectif. Posons $R = A_\mathfrak p^\wedge$ et soit $\eta$ l'objet formel
sur $R$ associé à $x|_{\Spec(R)}$. L'application $d\underline{\eta}$ sur les
espaces tangents est surjective, car elle s'identifie au dual de l'application
injective
$H^0(E \otimes_A^{\mathbf{L}} k) \to
H^0(\NL_{A/\Lambda} \otimes_A^{\mathbf{L}} k)$
(voir le paragraphe précédent). D'après Théorie formelle des déformations,
lemme \ref{formal-defos-lemma-versal-criterion}, il suffit de démontrer ce qui
suit. Soit $C' \to C$ une petite extension de $\Lambda$-algèbres locales
artiniennes de type fini, de corps résiduel $k$. Soit $R \to C$ un morphisme
de $\Lambda$-algèbres compatible aux identifications des corps résiduels. Soit
$y = x|_{\Spec(C)}$ et soit $y'$ un relèvement de $y$ à $C'$. Il s'agit de
montrer que nous pouvons relever le morphisme de $\Lambda$-algèbres
$R \to C$ en un morphisme $R \to C'$.

\medskip\noindent
Remarquons qu'il suffit de relever le morphisme de $\Lambda$-algèbres
$A \to C$. Soit $I = \Ker(C' \to C)$. Remarquons que $I$ est un espace
vectoriel de dimension $1$ sur $k$. L'obstruction $ob$ au relèvement de
$A \to C$ est un élément de $\Ext^1_A(\NL_{A/\Lambda}, I)$ ; voir l'exemple
\ref{example-key}. D'après le lemme \ref{lemma-compute-ext-into-field} et notre
hypothèse, le morphisme $\xi$ induit une injection
$$
\Ext^1_A(\NL_{A/\Lambda}, I)
\longrightarrow
\Ext^1_A(E, I)
$$
Par la construction de $ob$ et par (i), l'image de $ob$ dans
$\Ext^1_A(E, I)$ est l'obstruction au relèvement de $x$ à $A \times_C C'$.
Par (RS*), le fait que $y/C$ se relève à $y'/C'$ implique que $x$ se relève à
$A \times_C C'$. Ainsi, $ob = 0$, ce qui achève la démonstration.
\end{proof}

\noindent
Le lemme clé ci-dessus nous permet de conclure à l'ouverture de la versalité
dans certains cas.

\begin{lemma}
\label{lemma-openness}
Soit $S$ un schéma localement noethérien. Soit $\mathcal{X}$ une catégorie
fibrée en groupoïdes sur $(\Sch/S)_{fppf}$ satisfaisant à (RS*). Soit
$U = \Spec(A)$ un schéma affine de type fini sur $S$ qui s'envoie dans un
ouvert affine $\Spec(\Lambda)$. Soit $x$ un objet de $\mathcal{X}$ au-dessus
de $U$. Soit $\xi : E \to \NL_{A/\Lambda}$ un morphisme de $D^{-}(A)$.
Supposons que
\begin{enumerate}
\item[(i)] pour toute situation de déformation $(x, A' \to A)$, l'objet $x$
se relève à $\Spec(A')$ si et seulement si
$E \to \NL_{A/\Lambda} \to \NL_{A/A'}$ est nul ;
\item[(ii)] il existe un isomorphisme de foncteurs
$T_x(-) \to \Ext^0_A(E, -)$
tel que $E \to \NL_{A/\Lambda} \to \Omega^1_{A/\Lambda}$ corresponde à
l'élément canonique (voir la remarque \ref{remark-canonical-element}) ;
\item[(iii)] les groupes de cohomologie de $E$ soient des $A$-modules finis.
\end{enumerate}
Si $x$ est versel en un point fermé $u_0 \in U$, il existe alors un voisinage
ouvert $u_0 \in U' \subset U$ tel que $x$ soit versel en tout point de type
fini de $U'$.
\end{lemma}

\begin{proof}
Soit $C$ le cône de $\xi$, de sorte que nous ayons un triangle distingué
$$
E \to \NL_{A/\Lambda} \to C \to E[1]
$$
dans $D^{-}(A)$. D'après le lemme \ref{lemma-characterize-versal}, l'hypothèse
que $x$ est versel en $u_0$ implique que
$H^{-1}(C \otimes^\mathbf{L} k) = 0$. D'après Compléments sur l'algèbre, lemme
\ref{more-algebra-lemma-cut-complex-in-two}, il existe un $f \in A$ qui
n'appartient pas à l'idéal premier correspondant à $u_0$ tel que
$H^{-1}(C \otimes^\mathbf{L}_A M) = 0$ pour tout $A_f$-module $M$. En
appliquant encore le lemme \ref{lemma-characterize-versal}, nous voyons que la
versalité vaut en tous les points de type fini de l'ouvert $D(f) \subset U$.
\end{proof}

\noindent
Les lemmes techniques ci-dessus suggèrent la définition suivante.

\begin{definition}
\label{definition-naive-obstruction-theory}
Soit $S$ une base localement noethérienne. Soit $\mathcal{X}$ une catégorie
fibrée en groupoïdes sur $(\Sch/S)_{fppf}$. Supposons que $\mathcal{X}$
satisfasse à (RS*). Une {\it théorie naïve des obstructions} est donnée par
les données suivantes :
\begin{enumerate}
\item
\label{item-map}
pour toute $S$-algèbre $A$ telle que $\Spec(A) \to S$ s'envoie dans un ouvert
affine $\Spec(\Lambda) \subset S$ et pour tout objet $x$ de $\mathcal{X}$
au-dessus de $\Spec(A)$, on se donne un objet $E_x \in D^-(A)$ et un morphisme
$\xi_x : E \to \NL_{A/\Lambda}$ ;
\item
\label{item-inf}
pour $(x, A)$ comme dans (\ref{item-map}), il existe des transformations de
foncteurs
$$
\text{Inf}_x( - ) \to \Ext^{-1}_A(E_x, -)
\quad\text{et}\quad
T_x(-) \to \Ext^0_A(E_x, -)
$$
\item
\label{item-functoriality}
pour $(x, A)$ comme dans (\ref{item-map}) et un morphisme d'anneaux
$A \to B$, en posant $y = x|_{\Spec(B)}$, il existe un morphisme de
fonctorialité $E_x \to E_y$ dans $D(A)$.
\end{enumerate}
Ces données sont soumises aux conditions suivantes :
\begin{enumerate}
\item[(i)]
dans la situation de (\ref{item-functoriality}), le diagramme
$$
\xymatrix{
E_y \ar[r]_{\xi_y} & \NL_{B/\Lambda} \\
E_x \ar[u] \ar[r]^{\xi_x} & \NL_{A/\Lambda} \ar[u]
}
$$
soit commutatif dans $D(A)$ ;
\item[(ii)]
pour $(x, A)$ comme dans (\ref{item-map}) et $A \to B \to C$, en posant
$y = x|_{\Spec(B)}$ et $z = x|_{\Spec(C)}$, la composée des morphismes de
fonctorialité $E_x \to E_y$ et $E_y \to E_z$ soit le morphisme de
fonctorialité $E_x \to E_z$ ;
\item[(iii)]
les morphismes de (\ref{item-inf}) soient des isomorphismes compatibles avec
les morphismes de fonctorialité et les morphismes de la remarque
\ref{remark-functoriality} ;
\item[(iv)]
la composée $E_x \to \NL_{A/\Lambda} \to \Omega_{A/\Lambda}$ corresponde à
l'élément canonique de
$T_x(\Omega_{A/\Lambda}) = \Ext^0(E_x, \Omega_{A/\Lambda})$ ; voir la remarque
\ref{remark-canonical-element} ;
\item[(v)]
pour toute situation de déformation $(x, A' \to A)$ avec
$I = \Ker(A' \to A)$, la composée
$E_x \to \NL_{A/\Lambda} \to \NL_{A/A'}$ soit nulle dans
$$
\Hom_A(E_x, \NL_{A/\Lambda}) = \Ext^0_A(E_x, \NL_{A/A'}) =
\Ext^1_A(E_x, I)
$$
si et seulement si $x$ se relève à $A'$.
\end{enumerate}
\end{definition}

\noindent
Nous voyons ainsi en particulier que nous obtenons une théorie des
obstructions comme dans la section \ref{section-obstruction-theory} en posant
$\mathcal{O}_x( - ) = \Ext^1_A(E_x, -)$.

\begin{lemma}
\label{lemma-naive-obstruction-theory-qis}
Soient $S$ et $\mathcal{X}$ comme dans la définition
\ref{definition-naive-obstruction-theory}, et munissons $\mathcal{X}$ d'une
théorie naïve des obstructions. Soient $A \to B$ et $y \to x$ comme dans
(\ref{item-functoriality}). Soit $k$ une $B$-algèbre qui est un corps. Alors
le morphisme de fonctorialité $E_x \to E_y$ induit des bijections
$$
H^i(E_x \otimes_A^{\mathbf{L}} k) \to H^i(E_y \otimes_B^{\mathbf{L}} k)
$$
pour $i = 0, 1$.
\end{lemma}

\begin{proof}
Soit $z = x|_{\Spec(k)}$. Alors (RS*) implique que
$$
\textit{Lift}(x, A[k]) = \textit{Lift}(z, k[k])
\quad\text{et}\quad
\textit{Lift}(y, B[k]) = \textit{Lift}(z, k[k])
$$
car $A[k] = A \times_k k[k]$ et $B[k] = B \times_k k[k]$. Les propriétés
d'une théorie naïve des obstructions impliquent donc que le morphisme de
fonctorialité $E_x \to E_y$ induit des bijections
$\Ext^i_A(E_x, k) \to \text{Ext}^i_B(E_y, k)$ pour $i = -1, 0$. D'après le
lemme \ref{lemma-compute-ext-into-field}, nos applications
$H^i(E_x \otimes_A^{\mathbf{L}} k) \to H^i(E_y \otimes_B^{\mathbf{L}} k)$,
$i = 0, 1$, induisent des isomorphismes sur les espaces vectoriels duaux et
sont donc des isomorphismes.
\end{proof}

\begin{lemma}
\label{lemma-naive-obstruction-theory-gives-openness}
Soit $S$ un schéma localement noethérien. Soit
$p : \mathcal{X} \to (\Sch/S)_{fppf}^{opp}$ une catégorie fibrée en
groupoïdes. Supposons que $\mathcal{X}$ satisfasse à (RS*) et que
$\mathcal{X}$ possède une théorie naïve des obstructions. Alors l'ouverture
de la versalité vaut pour $\mathcal{X}$, pourvu que les complexes $E_x$ de la
définition \ref{definition-naive-obstruction-theory} aient des groupes de
cohomologie de type fini pour les paires $(A, x)$ où $A$ est de type fini sur
$S$.
\end{lemma}

\begin{proof}
Soit $U$ un schéma localement de type fini sur $S$, soit $x$ un objet de
$\mathcal{X}$ au-dessus de $U$, et soit $u_0$ un point de type fini de $U$ tel
que $x$ soit versel en $u_0$. Nous pouvons d'abord rétrécir $U$ en un schéma
affine tel que $u_0$ soit un point fermé et que $U \to S$ s'envoie dans un
ouvert affine $\Spec(\Lambda)$. Écrivons $U = \Spec(A)$. Soit
$\xi_x : E_x \to \NL_{A/\Lambda}$ le morphisme d'obstruction. Nous pouvons
alors appliquer le lemme \ref{lemma-openness} pour conclure.
\end{proof}









\section{Une notion duale}
\label{section-dual}

\noindent
Soit $(x, A' \to A)$ une situation de déformation pour une catégorie
$\mathcal{X}$ fibrée en groupoïdes donnée sur un schéma localement noethérien
$S$. Supposons que $\mathcal{X}$ possède une théorie des obstructions ; voir la
définition \ref{definition-obstruction-theory}. En pratique, on dispose souvent
d'un complexe $K^\bullet$ de $A$-modules et d'isomorphismes de foncteurs
$$
\text{Inf}_x(-) \to H^0(K^\bullet \otimes_A^\mathbf{L} -),\quad
T_x(-) \to H^1(K^\bullet \otimes_A^\mathbf{L} -),\quad
\mathcal{O}_x(-) \to H^2(K^\bullet \otimes_A^\mathbf{L} -)
$$
Dans cette section, nous formalisons un peu cette situation et montrons comment
elle permet, dans certains cas, de vérifier l'ouverture de la versalité.

\begin{example}
\label{example-global-sections-dual}
Soient $\Lambda, S, W, \mathcal{F}$ comme dans l'exemple
\ref{example-global-sections}. Supposons que $W \to S$ soit propre et que
$\mathcal{F}$ soit cohérent. D'après Cohomologie des schémas, remarque
\ref{coherent-remark-explain-perfect-direct-image}
il existe un complexe fini de $\Lambda$-modules projectifs de type fini
$N^\bullet$ qui calcule universellement la cohomologie de $\mathcal{F}$. En particulier,
les espaces d'obstructions de l'exemple \ref{example-global-sections} sont
$\mathcal{O}_x(M) = H^1(N^\bullet \otimes_\Lambda M)$. Ainsi, en posant
$K^\bullet = N^\bullet \otimes_\Lambda A[-1]$, nous voyons que
$\mathcal{O}_x(M) = H^2(K^\bullet \otimes_A^\mathbf{L} M)$.
\end{example}

\begin{situation}
\label{situation-dual}
Soit $S$ un schéma localement noethérien. Soit $\mathcal{X}$ une catégorie
fibrée en groupoïdes sur $(\Sch/S)_{fppf}$. Supposons que $\mathcal{X}$
satisfasse à (RS*), de sorte que nous puissions parler du foncteur $T_x(-)$ ;
voir le lemme \ref{lemma-properties-lift-RS-star}. Soit $U = \Spec(A)$ un
schéma affine de type fini sur $S$ qui s'envoie dans un ouvert affine
$\Spec(\Lambda)$. Soit $x$ un objet de $\mathcal{X}$ au-dessus de $U$.
Supposons que l'on se donne
\begin{enumerate}
\item un complexe de $A$-modules $K^\bullet$ ;
\item une transformation de foncteurs
$T_x(-) \to H^1(K^\bullet \otimes_A^\mathbf{L} -)$,
\item pour toute situation de déformation $(x, A' \to A)$ de noyau
$I = \Ker(A' \to A)$, un élément
$o_x(A') \in H^2(K^\bullet \otimes_A^\mathbf{L} I)$
\end{enumerate}
satisfaisant aux conditions (minimales) suivantes :
\begin{enumerate}
\item[(i)] la transformation
$T_x(-) \to H^1(K^\bullet \otimes_A^\mathbf{L} -)$
est un isomorphisme ;
\item[(ii)] étant donné un morphisme
$(x, A'' \to A) \to (x, A' \to A)$ de situations de déformation, l'élément
$o_x(A')$ s'envoie sur l'élément $o_x(A'')$ par l'application
$H^2(K^\bullet \otimes_A^\mathbf{L} I) \to
H^2(K^\bullet \otimes_A^\mathbf{L} I')$
où $I' = \Ker(A'' \to A)$ ;
\item[(iii)] $x$ se relève en un objet au-dessus de $\Spec(A')$ si et seulement si
$o_x(A') = 0$.
\end{enumerate}
Il est également possible d'incorporer les automorphismes infinitésimaux, mais
nous nous en abstenons afin d'obtenir le résultat le plus précis possible.
\end{situation}

\noindent
Dans la situation \ref{situation-dual}, l'objet
$K^\bullet \otimes_A^\mathbf{L} \NL_{A/\Lambda}$ jouera un rôle important.
Supposons que l'on se donne un élément
$\xi \in H^1(K^\bullet \otimes_A^\mathbf{L} \NL_{A/\Lambda})$. Alors (1) pour
toute surjection $A' \to A$ de $\Lambda$-algèbres dont le noyau $I$ est de
carré nul, l'application canonique
$\NL_{A/\Lambda} \to \NL_{A/A'} = I[1]$ envoie $\xi$ sur un élément
$\xi_{A'} \in H^2(K^\bullet \otimes_A^\mathbf{L} I)$ et (2) l'application
$\NL_{A/\Lambda} \to \Omega_{A/\Lambda}$ envoie $\xi$ sur un élément
$\xi_{can}$ de
$H^1(K^\bullet \otimes_A^\mathbf{L} \Omega_{A/\Lambda})$.

\begin{lemma}
\label{lemma-dual-obstruction}
Dans la situation \ref{situation-dual}, supposons en outre que
\begin{enumerate}
\item[(iv)] étant donnés une suite exacte courte de situations de déformation
comme dans la remarque \ref{remark-short-exact-sequence-thickenings} et un
relèvement $x'_2 \in \text{Lift}(x, A_2')$, alors
$o_x(A_3') \in H^2(K^\bullet \otimes_A^\mathbf{L} I_3)$
est égal à $\partial\theta$, où
$\theta \in H^1(K^\bullet \otimes_A^\mathbf{L} I_1)$
est l'élément correspondant à $x'_2|_{\Spec(A_1')}$ au moyen de
$A_1' = A[I_1]$ et de l'application donnée
$T_x(-) \to H^1(K^\bullet \otimes_A^\mathbf{L} -)$.
\end{enumerate}
Il existe alors un élément
$\xi \in H^1(K^\bullet \otimes_A^\mathbf{L} \NL_{A/\Lambda})$
tel que
\begin{enumerate}
\item pour toute situation de déformation $(x, A' \to A)$, on ait
$\xi_{A'} = o_x(A')$ ;
\item $\xi_{can}$ s'identifie à l'élément canonique de la remarque
\ref{remark-canonical-element} par la transformation donnée
$T_x(-) \to H^1(K^\bullet \otimes_A^\mathbf{L} -)$.
\end{enumerate}
\end{lemma}

\begin{proof}
Choisissons un morphisme $\alpha : \Lambda[x_1, \ldots, x_n] \to A$ de noyau
$J$. Posons $P = \Lambda[x_1, \ldots, x_n]$. Dans la suite de cette
démonstration, nous travaillons avec
$$
\NL(\alpha) = (J/J^2 \longrightarrow \bigoplus A \text{d}x_i)
$$
ce qui est permis par Algèbre, lemme \ref{algebra-lemma-NL-homotopy}, et
Compléments sur l'algèbre, lemme
\ref{more-algebra-lemma-derived-tor-homotopy}. Considérons l'élément
$o_x(P/J^2) \in H^2(K^\bullet \otimes_A^\mathbf{L} J/J^2)$ et le quotient
$$
C = (P/J^2 \times \bigoplus A \text{d}x_i)/(J/J^2)
$$
où $J/J^2$ est plongé diagonalement. Remarquons que $C \to A$ est une
surjection de noyau $\bigoplus A\text{d}x_i$. Il existe en outre une section
$A \to C$ de $C \to A$, donnée en envoyant la classe de $f \in P$ sur la
classe de $(f, \text{d}f)$ dans la somme amalgamée. Pour la suite, notons
$x_C$ l'image réciproque de $x$ par le morphisme correspondant
$\Spec(C) \to \Spec(A)$. Nous voyons ainsi que $o_x(C) = 0$. Nous en concluons
que $o_x(P/J^2)$ s'envoie sur zéro dans
$H^2(K^\bullet \otimes_A^\mathbf{L} \bigoplus A\text{d}x_i)$.
Il s'ensuit qu'il existe un élément
$\xi \in H^1(K^\bullet \otimes_A^\mathbf{L} \NL(\alpha))$
qui s'envoie sur $o_x(P/J^2)$.

\medskip\noindent
Remarquons que, pour toute situation de déformation $(x, A' \to A)$, il existe
un morphisme de $\Lambda$-algèbres $P/J^2 \to A'$ compatible avec les
augmentations vers $A$. Par conséquent, l'élément $\xi$ satisfait à la première
propriété du lemme par construction et à la propriété (ii) de la situation
\ref{situation-dual}.

\medskip\noindent
Remarquons que notre choix de $\xi$ n'est déterminé qu'à l'addition d'un
élément de $H^1(K^\bullet \otimes_A^\mathbf{L} \bigoplus A\text{d}x_i)$.
Nous allons montrer qu'après avoir modifié $\xi$ par un élément du groupe
précédent, nous pouvons faire en sorte que la seconde assertion du lemme soit
vraie. Soit $C' \subset C$ l'image de $P/J^2$ par le morphisme de
$\Lambda$-algèbres $P/J^2 \to C$ (inclusion du premier facteur). Observons que
$\Ker(C' \to A) = \Im(J/J^2 \to \bigoplus A\text{d}x_i)$.
Posons $\overline{C} = A[\Omega_{A/\Lambda}]$. L'application
$P/J^2 \times \bigoplus A \text{d}x_i \to \overline{C}$,
$(f, \sum f_i \text{d}x_i) \mapsto (f \bmod J, \sum f_i \text{d}x_i)$
se factorise par une application surjective $C \to \overline{C}$. Alors
$$
(x, \overline{C} \to A) \to (x, C \to A) \to (x, C' \to A)
$$
est une suite exacte courte de situations de déformation. Le scindage associé
$\overline{C} = A[\Omega_{A/\Lambda}]$ (provenant de la remarque
\ref{remark-short-exact-sequence-thickenings}) est égal au scindage donné
ci-dessus. De plus, la section $A \to C$, composée avec l'application
$C \to \overline{C}$, est l'application
$(1, \text{d}) : A \to A[\Omega_{A/\Lambda}]$ de la remarque
\ref{remark-canonical-element}. Ainsi, $x_C$ se restreint en l'élément
canonique $x_{can}$ de
$T_x(\Omega_{A/\Lambda}) = \text{Lift}(x, A[\Omega_{A/\Lambda}])$.
D'après la condition (iv), nous en concluons que $o_x(P/J^2)$ s'envoie sur
$\partial x_{can}$ dans
$$
H^1(K^\bullet \otimes_A^\mathbf{L} \Im(J/J^2 \to \bigoplus A\text{d}x_i))
$$
Par construction, $\xi$ s'envoie sur $o_x(P/J^2)$. Il s'ensuit que
$x_{can}$ et $\xi_{can}$ s'envoient sur le même élément du groupe affiché, ce
qui signifie (d'après la suite exacte longue de cohomologie) que leur différence
provient d'un élément de
$H^1(K^\bullet \otimes_A^\mathbf{L} \bigoplus A\text{d}x_i)$
comme voulu.
\end{proof}

\begin{lemma}
\label{lemma-dual-openness}
Dans la situation \ref{situation-dual}, supposons que la condition (iv) du
lemme \ref{lemma-dual-obstruction} soit satisfaite et que $K^\bullet$ soit un
objet parfait de $D(A)$. Alors, si $x$ est versel en un point fermé
$u_0 \in U$, il existe un voisinage ouvert $u_0 \in U' \subset U$ tel que $x$
soit versel en tout point de type fini de $U'$.
\end{lemma}

\begin{proof}
Nous pouvons supposer que $K^\bullet$ est un complexe fini de $A$-modules
projectifs de type fini. Ainsi, le produit tensoriel dérivé par $K^\bullet$
revient à tensoriser simplement par $K^\bullet$. Soit $E^\bullet$ le complexe
parfait dual de $K^\bullet$ ; voir Compléments sur l'algèbre, lemme
\ref{more-algebra-lemma-dual-perfect-complex}. (Ainsi,
$E^n = \Hom_A(K^{-n}, A)$ et ses différentielles sont les transposées des
différentielles de $K^\bullet$.) Soit $E \in D^{-}(A)$ l'objet représenté par
le complexe $E^\bullet[-1]$. Soit
$\xi \in H^1(\text{Tot}(K^\bullet \otimes_A \NL_{A/\Lambda}))$ l'élément
construit au lemme \ref{lemma-dual-obstruction}, et notons
$\xi : E = E^\bullet[-1] \to \NL_{A/\Lambda}$ le morphisme correspondant
(loc. cit.). Nous affirmons que la paire $(E, \xi)$ satisfait à toutes les
hypothèses du lemme \ref{lemma-openness}, ce qui achève la démonstration.

\medskip\noindent
En effet, l'hypothèse (i) du lemme \ref{lemma-openness} découle de la conclusion
(1) du lemme \ref{lemma-dual-obstruction} et du fait que
$H^2(K^\bullet \otimes_A^\mathbf{L} -) = \Ext^1(E, -)$ d'après loc. cit.
L'hypothèse (ii) du lemme \ref{lemma-openness} découle de la conclusion (2) du
lemme \ref{lemma-dual-obstruction} et du fait que
$H^1(K^\bullet \otimes_A^\mathbf{L} -) = \Ext^0(E, -)$ d'après loc. cit.
L'hypothèse (iii) du lemme \ref{lemma-openness} est claire.
\end{proof}







\section{Foncteurs commutant aux limites sur les schémas noethériens}
\label{section-noetherian}

\noindent
Il est parfois commode de considérer des foncteurs ou des champs définis
seulement sur la sous-catégorie pleine des schémas (localement) noethériens.
Dans cette section, nous étudions cette question dans le cas des espaces
algébriques.

\medskip\noindent
Soit $S$ un schéma localement noethérien. Soyons quelque peu pédants afin
d'aligner correctement nos catégories ; ceux qui font abstraction des
questions ensemblistes peuvent simplement remplacer, dans ce qui suit, les
ensembles de schémas que nous choisissons par la collection de tous les
schémas. Comme dans Topologies, remarque
\ref{topologies-remark-choice-sites}, choisissons une catégorie
$\Sch_\alpha$ de schémas contenant $S$ telle que nous obtenions les gros sites
$(\Sch/S)_{Zar}$,
$(\Sch/S)_\etale$, $(\Sch/S)_{smooth}$, $(\Sch/S)_{syntomic}$,
et $(\Sch/S)_{fppf}$, ayant tous pour même catégorie sous-jacente
$\Sch_\alpha/S$. Notons
$$
\textit{Noetherian}_\alpha \subset \Sch_\alpha
$$
la sous-catégorie pleine formée des schémas localement noethériens. Elle
détermine une sous-catégorie pleine
$$
\textit{Noetherian}_\alpha/S \subset \Sch_\alpha/S
$$
Pour $\tau \in \{Zariski, \etale, smooth, syntomic, fppf\}$, nous avons :
\begin{enumerate}
\item si $f : X \to Y$ est un morphisme de $\Sch_\alpha/S$, si $Y$ appartient
à $\textit{Noetherian}_\alpha/S$ et si $f$ est localement de type fini, alors
$X$ appartient à $\textit{Noetherian}_\alpha/S$ ;
\item pour des morphismes $f : X \to Y$ et $g : Z \to Y$ de
$\textit{Noetherian}_\alpha/S$, avec $f$ localement de type fini, le produit
fibré $X \times_Y Z$ dans $\textit{Noetherian}_\alpha/S$ existe et coïncide
avec le produit fibré dans $\Sch_\alpha/S$ ;
\item si $\{X_i \to X\}_{i \in I}$ est un recouvrement de
$(\Sch/S)_\tau$ et si $X$ appartient à $\textit{Noetherian}_\alpha/S$, alors
les objets $X_i$ appartiennent à $\textit{Noetherian}_\alpha/S$ ;
\item la catégorie $\textit{Noetherian}_\alpha/S$, munie de l'ensemble des
recouvrements de $(\Sch/S)_\tau$ dont les objets appartiennent à
$\textit{Noetherian}_\alpha/S$, est un site que nous noterons
$(\textit{Noetherian}/S)_\tau$ ;
\item le foncteur d'inclusion
$(\textit{Noetherian}/S)_\tau \to (\Sch/S)_\tau$
est pleinement fidèle, continu et cocontinu.
\end{enumerate}
D'après Sites, lemmes \ref{sites-lemma-cocontinuous-morphism-topoi} et
\ref{sites-lemma-when-shriek}, nous obtenons un morphisme de topos
$$
g_\tau : \Sh((\textit{Noetherian}/S)_\tau) \longrightarrow \Sh((\Sch/S)_\tau)
$$
dont le foncteur image inverse est la restriction des faisceaux le long du
foncteur d'inclusion
$(\textit{Noetherian}/S)_\tau \to (\Sch/S)_\tau$.

\begin{remark}[Avertissement]
\label{remark-no-fibre-products}
Le site $(\textit{Noetherian}/S)_\tau$ ne possède pas de produits fibrés.
Nous devons donc être prudents lorsque nous travaillons avec des faisceaux.
Par exemple, le foncteur d'inclusion continu
$(\textit{Noetherian}/S)_\tau \to (\Sch/S)_\tau$
ne définit pas un morphisme de sites. Voir Exemples, section
\ref{examples-section-sheaves-locally-Noetherian}, pour un exemple dans le cas
$\tau = fppf$.
\end{remark}

\noindent
Soit $F : (\textit{Noetherian}/S)_\tau^{opp} \to \textit{Sets}$ un foncteur.
Nous disons que $F$ {\it commute aux limites} si, pour toute limite dirigée de
schémas affines $X = \lim X_i$ de $(\textit{Noetherian}/S)_\tau$, on a
$F(X) = \colim F(X_i)$.

\begin{lemma}
\label{lemma-canonical-extension}
Soit $\tau \in \{Zariski, \etale, smooth, syntomic, fppf\}$. La restriction le
long du foncteur d'inclusion
$(\textit{Noetherian}/S)_\tau \to (\Sch/S)_\tau$
définit une équivalence de catégories entre
\begin{enumerate}
\item la catégorie des faisceaux commutant aux limites sur
$(\Sch/S)_\tau$ ;
\item la catégorie des faisceaux commutant aux limites sur
$(\textit{Noetherian}/S)_\tau$
\end{enumerate}
\end{lemma}

\begin{proof}
Soit $F : (\textit{Noetherian}/S)_\tau^{opp} \to \textit{Sets}$ un foncteur
qui est à la fois un faisceau et commute aux limites. D'après Topologies, lemmes
\ref{topologies-lemma-extend} et \ref{topologies-lemma-extend-sheaf-general}
il existe un unique foncteur
$F' : (\Sch/S)_\tau^{opp} \to \textit{Sets}$
qui commute aux limites, est un faisceau et se restreint en $F$. En fait, la
construction de $F'$ dans Topologies, lemme \ref{topologies-lemma-extend}, est
fonctorielle en $F$, et cette construction est un quasi-inverse de la
restriction. Certains détails sont omis.
\end{proof}

\begin{lemma}
\label{lemma-representable-limit-preserving}
Soit $X$ un objet de $(\textit{Noetherian}/S)_\tau$. Si le foncteur des points
$h_X : (\textit{Noetherian}/S)_\tau^{opp} \to \textit{Sets}$ commute aux
limites, alors $X$ est localement de présentation finie sur $S$.
\end{lemma}

\begin{proof}
Soit $V \subset X$ un sous-schéma ouvert affine qui s'envoie dans un ouvert
affine $U \subset S$. Nous pouvons écrire $V = \lim V_i$ comme limite dirigée
de schémas affines $V_i$ de présentation finie sur $U$ ; voir Algèbre, lemme
\ref{algebra-lemma-ring-colimit-fp}. Par hypothèse, la flèche $V \to X$ se
factorise sous la forme $V \to V_i \to X$ pour un certain $i$. En augmentant
$i$, nous pouvons supposer que $V_i \to X$ se factorise par $V$, car l'image
réciproque de $V \subset X$ dans $V_i$ finit par être égale à $V_i$ d'après
Limites, lemme \ref{limits-lemma-descend-opens}. Alors le morphisme identité
$V \to V$ se factorise par $V_i$ pour un certain $i$ dans la catégorie des
schémas au-dessus de $U$. Ainsi, $V \to U$ est de présentation finie ; le fait
algébrique correspondant est que, si $B$ est une $A$-algèbre telle que
$\text{id} : B \to B$ se factorise par une $A$-algèbre de présentation finie,
alors $B$ est de présentation finie sur $A$ (joli exercice). Par conséquent,
$X$ est localement de présentation finie sur $S$.
\end{proof}

\noindent
Le lemme suivant possède une variante pour les transformations représentables
par des espaces algébriques.

\begin{lemma}
\label{lemma-representable}
Soit $\tau \in \{Zariski, \etale, smooth, syntomic, fppf\}$. Soient
$F', G' : (\Sch/S)_\tau^{opp} \to \textit{Sets}$ des faisceaux commutant aux
limites. Soit $a' : F' \to G'$ une transformation de foncteurs. Notons
$a : F \to G$ la restriction de $a' : F' \to G'$ à
$(\textit{Noetherian}/S)_\tau$. Les assertions suivantes sont équivalentes :
\begin{enumerate}
\item $a'$ est représentable (comme transformation de foncteurs ; voir
Catégories, définition \ref{categories-definition-representable-morphism}) ;
\item pour tout objet $V$ de $(\textit{Noetherian}/S)_\tau$ et toute
application $V \to G$, le produit fibré
$F \times_G V : (\textit{Noetherian}/S)_\tau^{opp} \to \textit{Sets}$
est un foncteur représentable ;
\item comme en (2), mais seulement pour $V$ affine de type fini sur $S$ et
s'envoyant dans un ouvert affine de $S$.
\end{enumerate}
\end{lemma}

\begin{proof}
Supposons (1). D'après Limites d'espaces, lemme
\ref{spaces-limits-lemma-locally-finite-presentation-permanence}
la transformation $a'$ commute aux limites\footnote{Cela a un sens même si
$\tau \not = fppf$, puisque la catégorie sous-jacente à $(\Sch/S)_\tau$ est
égale à la catégorie sous-jacente à $(\Sch/S)_{fppf}$ et que l'énoncé ne fait
pas intervenir la topologie.}. Prenons $\xi : V \to G$ comme en (2). Notons
$V' = V$, mais considéré comme un objet de $(\Sch/S)_\tau$. Puisque $G$ est la
restriction de $G'$ à $(\textit{Noetherian}/S)_\tau$, nous voyons que
$\xi \in G(V)$ correspond à $\xi' \in G'(V')$. Par hypothèse,
$V' \times_{\xi', G'} F'$ est représentable par un schéma $U'$. Le morphisme
de schémas $U' \to V'$ correspondant à la projection
$V' \times_{\xi', G'} F' \to V'$ est localement de présentation finie d'après
Limites d'espaces, lemme
\ref{spaces-limits-lemma-base-change-locally-finite-presentation} et
Limites, proposition
\ref{limits-proposition-characterize-locally-finite-presentation}.
Ainsi, $U'$ est un schéma localement noethérien et $U'$ est donc isomorphe à un
objet $U$ de $(\textit{Noetherian}/S)_\tau$. Alors $U$ représente
$F \times_G V$, comme voulu.

\medskip\noindent
L'implication (2) $\Rightarrow$ (3) est immédiate. Supposons (3). Nous allons
démontrer (1). Soit $T$ un objet de $(\Sch/S)_\tau$ et soit $T \to G'$ un
morphisme. Nous devons montrer que le foncteur $F' \times_{G'} T$ est
représentable par un schéma $X$ au-dessus de $T$. Soit $\mathcal{B}$ l'ensemble
des ouverts affines de $T$ qui s'envoient dans un ouvert affine de $S$. C'est
une base de la topologie de $T$. Nous montrerons ci-dessous que, pour
$W \in \mathcal{B}$, le produit fibré $F' \times_{G'} W$ est représentable par
un schéma $X_W$ au-dessus de $W$. Si $W_1 \subset W_2$ dans $\mathcal{B}$,
nous obtenons alors un isomorphisme
$X_{W_1} \to X_{W_2} \times_{W_2} W_1$, car $X_{W_1}$ et
$X_{W_2} \times_{W_2} W_1$ représentent tous deux le foncteur
$F' \times_{G'} W_1$. Ces isomorphismes sont canoniques et satisfont à la
condition de cocycle mentionnée dans Constructions, lemme
\ref{constructions-lemma-relative-glueing}. Nous pouvons donc recoller les
schémas $X_W$ en un schéma $X$ au-dessus de $T$. La compatibilité des
applications de recollement avec les applications $X_W \to F'$ nous fournit
une application $X \to F'$. L'application qui en résulte
$X \to F' \times_{G'} T$ est un isomorphisme, comme on peut le vérifier
localement sur $T$ (car la source et le but de cette flèche sont des faisceaux
pour la topologie de Zariski).

\medskip\noindent
Soit $W$ un schéma affine qui s'envoie dans un ouvert affine $U \subset S$.
Soit $W \to G'$ une application. Toujours sous l'hypothèse (3), nous devons
montrer que $F' \times_{G'} W$ est représentable par un schéma. Nous pouvons
écrire $W = \lim V'_i$ comme limite dirigée de schémas affines $V'_i$ de
présentation finie sur $U$ ; voir Algèbre, lemme
\ref{algebra-lemma-ring-colimit-fp}. Comme $V'_i$ est de type fini sur un
schéma noethérien, nous voyons que $V'_i$ est un schéma noethérien. Notons
$V_i = V'_i$, mais considéré comme un objet de
$(\textit{Noetherian}/S)_\tau$. Comme $G'$ commute aux limites, nous pouvons
choisir un $i$ et une application $V'_i \to G'$ tels que $W \to G'$ soit la
composée $W \to V'_i \to G'$. Puisque $G$ est la restriction de $G'$ à
$(\textit{Noetherian}/S)_\tau$, le morphisme $V'_i \to G'$ équivaut à un
morphisme $V_i \to G$ (voir ci-dessus). Par l'hypothèse (3), le foncteur
$F \times_G V_i$ est représentable par un objet $X_i$ de
$(\textit{Noetherian}/S)_\tau$. Le foncteur $F \times_G V_i$ commute aux
limites, puisqu'il est la restriction de $F' \times_{G'} V'_i$ et que ce
foncteur commute aux limites d'après Limites d'espaces, lemme
\ref{spaces-limits-lemma-fibre-product-locally-finite-presentation},
l'hypothèse que $F'$ et $G'$ commutent aux limites, et Limites, remarque
\ref{limits-remark-limit-preserving}, qui nous dit que le foncteur des points
de $V'_i$ commute aux limites. D'après le lemme
\ref{lemma-representable-limit-preserving}, nous en concluons que $X_i$ est
localement de présentation finie sur $S$. Notons $X'_i = X_i$, mais considéré
comme un objet de $(\Sch/S)_\tau$. Nous voyons alors que
$F' \times_{G'} V'_i$ et le foncteur des points $h_{X'_i}$ sont tous deux des
extensions de
$h_{X_i} : (\textit{Noetherian}/S)_\tau^{opp} \to \textit{Sets}$
en des faisceaux commutant aux limites sur $(\Sch/S)_\tau$. D'après
l'équivalence de catégories du lemme \ref{lemma-canonical-extension}, nous
en déduisons que $X'_i$ représente $F' \times_{G'} V'_i$. Enfin,
$$
F' \times_{G'} W = F' \times_{G'} V'_i \times_{V'_i} W =
X'_i \times_{V'_i} W
$$
est représentable, comme voulu.
\end{proof}






\section{Espaces algébriques dans le cadre noethérien}
\label{section-algebraic-spaces-noetherian}

\noindent
Soit $S$ un schéma localement noethérien. Notons
$(\textit{Noetherian}/S)_\etale \subset (\Sch/S)_\etale$
le site étudié à la section \ref{section-noetherian}. Soit
$F : (\textit{Noetherian}/S)_\etale^{opp} \to \textit{Sets}$ un foncteur,
c'est-à-dire que $F$ est un préfaisceau sur
$(\textit{Noetherian}/S)_\etale$. Dans ce cadre, tous les axiomes [-1], [0],
[1], [2], [3], [4], [5] de la section \ref{section-axioms-functors} ont un
sens. Nous allons les passer en revue un à un et préciser exactement ce que
nous entendons.

\medskip\noindent
Axiome [-1]. Il s'agit d'une condition ensembliste que peuvent ignorer les
lecteurs qui ne s'intéressent pas aux questions ensemblistes. Elle a un
sens pour $F$, puisqu'elle concerne l'évaluation de $F$ sur des spectres de
corps de type fini sur $S$, lesquels sont des objets de
$(\textit{Noetherian}/S)_\etale$.

\medskip\noindent
Axiome [0]. C'est l'axiome selon lequel $F$ est un faisceau sur
$(\textit{Noetherian}/S)_\etale^{opp}$, c'est-à-dire satisfait à la condition
de faisceau pour les recouvrements étales.

\medskip\noindent
Axiome [1]. C'est l'axiome selon lequel $F$ commute aux limites au sens de la
section \ref{section-noetherian} : pour toute limite dirigée de schémas affines
$X = \lim X_i$ de $(\textit{Noetherian}/S)_\etale$, on a
$F(X) = \colim F(X_i)$.

\medskip\noindent
Axiome [2]. C'est l'axiome selon lequel $F$ satisfait à la condition (RS) de
Rim--Schlessinger. En examinant la définition de la condition (RS) dans la
définition \ref{definition-RS} et la discussion de la section
\ref{section-axioms-functors}, nous voyons que cela signifie ceci : pour toute
somme amalgamée $Y' = Y \amalg_X X'$ de schémas de type fini sur $S$, où
$Y, X, X'$ sont des spectres d'anneaux locaux artiniens, l'application
$$
F(Y \amalg_X X') \to F(Y) \times_{F(X)} F(X')
$$
est une bijection. Cette condition a un sens, car les schémas $X$, $X'$, $Y$
et $Y'$ appartiennent à $(\text{Noetherian}/S)_\etale$, puisqu'ils sont de type
fini sur $S$.

\medskip\noindent
Axiome [3]. C'est l'axiome selon lequel tout espace tangent $TF_{k, x_0}$ est
de dimension finie. Cela a un sens, car les espaces tangents $TF_{k, x_0}$ sont
construits à partir des évaluations de $F$ en $\Spec(k)$ et
$\Spec(k[\epsilon])$, où $k$ est un corps de type fini sur $S$ ; ils
s'obtiennent donc en évaluant sur des objets de la catégorie
$(\textit{Noetherian}/S)_\etale$.

\medskip\noindent
Axiome [4]. C'est l'axiome selon lequel tout objet formel est effectif. En
examinant la discussion des sections \ref{section-formal-objects} et
\ref{section-axioms-functors}, nous voyons que cela ne fait intervenir que
l'évaluation de notre foncteur sur des schémas noethériens ; cette condition a
donc un sens pour $F$.

\medskip\noindent
Axiome [5]. C'est l'axiome selon lequel $F$ satisfait à l'ouverture de la
versalité. Rappelons que cela signifie ce qui suit : étant donnés un schéma $U$
localement de type fini sur $S$, un élément $x \in F(U)$ et un point de type
fini $u_0 \in U$ tel que $x$ soit versel en $u_0$, il existe un voisinage
ouvert $u_0 \in U' \subset U$ tel que $x$ soit versel en tout point de type
fini de $U'$. Comme précédemment, cette vérification ne fait intervenir que
l'évaluation de notre foncteur sur des schémas noethériens.

\begin{proposition}
\label{proposition-spaces-diagonal-representable-noetherian}
Soit $S$ un schéma localement noethérien. Soit
$F : (\textit{Noetherian}/S)_\etale^{opp} \to \textit{Sets}$
un foncteur. Supposons que
\begin{enumerate}
\item $\Delta : F \to F \times F$ soit représentable (comme transformation de
foncteurs ; voir Catégories, définition
\ref{categories-definition-representable-morphism}) ;
\item $F$ satisfasse aux axiomes [-1], [0], [1], [2], [3], [4], [5]
(voir ci-dessus) ;
\item $\mathcal{O}_{S, s}$ soit un anneau G pour tout point de type fini $s$ de
$S$.
\end{enumerate}
Alors il existe un unique espace algébrique
$F' : (\Sch/S)_{fppf}^{opp} \to \textit{Sets}$
dont la restriction à $(\textit{Noetherian}/S)_\etale$ est $F$ (voir la
démonstration pour une précision).
\end{proposition}

\begin{proof}
Rappelons que les sites $(\Sch/S)_{fppf}$ et $(\Sch/S)_\etale$ ont la même
catégorie sous-jacente ; voir la discussion de la section
\ref{section-noetherian}. De même, les sites
$(\textit{Noetherian}/S)_\etale$ et
$(\textit{Noetherian}/S)_{fppf}$ ont la même catégorie sous-jacente. D'après
les axiomes [0] et [1], le foncteur $F$ est un faisceau et commute aux limites.
Soit $F' : (\Sch/S)_\etale^{opp} \to \textit{Sets}$ l'unique extension de $F$
qui soit un faisceau (pour la topologie étale) et qui commute aux limites ;
voir le lemme \ref{lemma-canonical-extension}. Alors $F'$ satisfait aux axiomes
[0] et [1] tels qu'ils sont énoncés à la section
\ref{section-axioms-functors}. D'après le lemme \ref{lemma-representable}, nous
voyons que $\Delta' : F' \to F' \times F'$ est représentable (par des schémas).
D'autre part, il est immédiatement clair que $F'$ satisfait aux axiomes [-1],
[2], [3], [4], [5] de la section \ref{section-axioms-functors}, puisque chacun
d'eux ne fait intervenir que l'évaluation de $F'$ sur des objets de
$(\textit{Noetherian}/S)_\etale$ et que nous avons supposé les conditions
correspondantes pour $F$. Par conséquent, $F'$ est un espace algébrique
d'après la proposition \ref{proposition-spaces-diagonal-representable}.
\end{proof}












\section{Théorème d'Artin sur les contractions}
\label{section-contractions}

\noindent
Dans cette section, nous utiliserons librement le langage des espaces
algébriques formels ; voir Espaces formels, section
\ref{formal-spaces-section-introduction}. Le théorème d'Artin sur les
contractions est l'un des deux théorèmes principaux de l'article d'Artin
\cite{ArtinII} ; le premier est son théorème sur les dilatations, que nous avons
énoncé et démontré dans Algébrisation des espaces formels, section
\ref{restricted-section-dilatations}.

\begin{situation}
\label{situation-contractions}
Soit $S$ un schéma localement noethérien. Soit $X'$ un espace algébrique
localement de type fini sur $S$. Soit $T' \subset |X'|$ un sous-ensemble fermé.
Soit $U' \subset X'$ le sous-espace ouvert tel que
$|U'| = |X'| \setminus T'$. Soit $W$ un espace algébrique formel localement
noethérien sur $S$, tel que $W_{red}$ soit localement de type fini sur $S$.
Enfin, soit
$$
g : X'_{/T'} \longrightarrow W
$$
une modification formelle ; voir Algébrisation des espaces formels, définition
\ref{restricted-definition-formal-modification}. Rappelons que $X'_{/T'}$
désigne le complété formel de $X'$ le long de $T'$ ; voir Espaces formels,
section \ref{formal-spaces-section-completion}.
\end{situation}

\noindent
Dans la situation ci-dessus, notre but est de démontrer qu'il existe un
morphisme propre $f : X' \to X$ d'espaces algébriques sur $S$, un sous-ensemble
fermé $T \subset |X|$ et un isomorphisme $a : X_{/T} \to W$ d'espaces
algébriques formels tels que
\begin{enumerate}
\item $T'$ soit l'image réciproque de $T$ par $|f| : |X'| \to |X|$ ;
\item $f : X' \to X$ envoie $U'$ isomorphiquement sur un sous-espace ouvert
$U$ de $X$ ;
\item $g = a \circ f_{/T}$, où $f_{/T} : X'_{/T'} \to X_{/T}$ est le morphisme
induit.
\end{enumerate}
Nous dirons que $(f : X' \to X, T, a)$ est une {\it solution}.

\medskip\noindent
Nous suivrons la stratégie d'Artin : nous construirons un foncteur $F$ sur la
catégorie des schémas localement noethériens sur $S$, nous montrerons que $F$
est un espace algébrique au moyen de la proposition
\ref{proposition-spaces-diagonal-representable-noetherian}, puis nous
démontrerons que le choix $X = F$ convient.

\begin{remark}
\label{remark-G-rings}
En particulier, nous ne pouvons pas démontrer que le résultat souhaité vaut
pour toute situation \ref{situation-contractions}, car nous devrons supposer
que les anneaux locaux de $S$ sont des anneaux G. Si vous pouvez démontrer le
résultat en général ou si vous disposez d'un contre-exemple, merci de nous le
faire savoir à l'adresse
\href{mailto:stacks.project@gmail.com}{stacks.project@gmail.com}.
\end{remark}

\noindent
Dans la situation \ref{situation-contractions}, soit $V$ un schéma localement
noethérien sur $S$. La valeur de notre foncteur $F$ en $V$ sera l'ensemble des
triplets
$$
(Z, u' : V \setminus Z \to U', \hat x : V_{/Z} \to W)
$$
satisfaisant aux conditions suivantes :
\begin{enumerate}
\item $Z \subset V$ est un sous-ensemble fermé ;
\item $u' : V \setminus Z \to U'$ est un morphisme sur $S$ ;
\item $\hat x : V_{/Z} \to W$ est un morphisme adique d'espaces algébriques
formels sur $S$ ;
\item $u'$ et $\hat x$ sont compatibles (voir ci-dessous).
\end{enumerate}
La condition de compatibilité est la suivante : par changement de base de la
modification formelle $g$, nous obtenons une modification formelle
$$
X'_{/T'} \times_{g, W, \hat x} V_{/Z} \longrightarrow V_{/Z}
$$
Voir Algébrisation des espaces formels, lemme
\ref{restricted-lemma-base-change-formal-modification}. D'après le théorème
principal sur les dilatations (Algébrisation des espaces formels, théorème
\ref{restricted-theorem-dilatations}), il existe un unique morphisme propre
$V' \to V$ d'espaces algébriques qui soit un isomorphisme au-dessus de
$V \setminus Z$ et tel que $V'_{/Z} \to V_{/Z}$ soit isomorphe à la flèche
affichée. Autrement dit, pour un certain morphisme
$\hat x' : V'_{/Z} \to X'_{/T'}$, nous avons un diagramme cartésien
$$
\xymatrix{
V'_{/Z} \ar[r] \ar[d]_{\hat x'} & V_{/Z} \ar[d]^{\hat x} \\
X'_{/T'} \ar[r]^g & W
}
$$
d'espaces algébriques formels. Nous considérerons désormais
$V \setminus Z$ comme un sous-espace ouvert de $V'$, sans autre mention. La
condition de compatibilité exige l'existence d'un morphisme $x' : V' \to X'$
qui se restreigne en $u'$ et $\hat x$ sur $V \setminus Z \subset V'$ et
$V'_{/Z}$. Autrement dit, le diagramme
$$
\xymatrix{
V \setminus Z \ar[r] \ar[d]_{u'} &
V' \ar[d]^{x'} &
V'_{/Z} \ar[l] \ar[d]^{\hat x'} \ar[r] &
V_{/Z} \ar[d]^{\hat x} \\
U' \ar[r] &
X' &
X'_{/T'} \ar[r]^g \ar[l] &
W
}
$$
est commutatif. Remarquons que, d'après Algébrisation des espaces formels,
lemme \ref{restricted-lemma-faithful}, le morphisme $x'$ est unique s'il
existe. Nous indiquerons cette situation en disant que
« {\it $V' \to V$, $\hat x'$ et $x'$ témoignent de la compatibilité entre
$u'$ et $\hat x$} ».

\begin{remark}
\label{remark-how-to-think-compatibility}
Dans la situation \ref{situation-contractions}, soit $V$ un schéma localement
noethérien sur $S$. Soit $(Z, u', \hat x)$ un triplet satisfaisant aux
conditions (1), (2) et (3) ci-dessus. Nous voulons expliquer une manière de
concevoir la condition de compatibilité (4). Elle ne sera pas mathématiquement
précise, car nous allons utiliser une catégorie fictive $\textit{An}_S$
d'espaces analytiques sur $S$ et un foncteur d'analytification fictif
$$
\left\{
\begin{matrix}
\text{espaces algébriques formels localement} \\
\text{noethériens sur }S
\end{matrix}
\right\}
\longrightarrow
\textit{An}_S,
\quad\quad
Y \longmapsto Y^{an}
$$
Par exemple, si $Y = \text{Spf}(k[[t]])$ sur $S = \Spec(k)$, alors $Y^{an}$
doit être considéré comme un disque unité ouvert. Si $Y = \Spec(k)$, alors
$Y^{an}$ est un point unique. La catégorie $\textit{An}_S$ devrait posséder
des immersions ouvertes et fermées, et nous devrions pouvoir prendre le
complémentaire ouvert d'un fermé. Étant donné $Y$, le morphisme
$Y_{red} \to Y$ devrait induire une immersion fermée
$Y_{red}^{an} \to Y^{an}$. Posons
$Y^{rig} = Y^{an} \setminus Y_{red}^{an}$, égal à ce complémentaire ouvert.
Si $Y$ est un espace algébrique et si $Z \subset Y$ est fermé, alors le
morphisme $Y_{/Z} \to Y$ devrait induire une immersion ouverte
$Y_{/Z}^{an} \to Y^{an}$, laquelle devrait à son tour induire une immersion
ouverte
$$
can : (Y_{/Z})^{rig} \longrightarrow (Y \setminus Z)^{an}
$$
De plus, étant donnée une modification formelle $g : Y' \to Y$ d'espaces
algébriques formels localement noethériens, le morphisme induit
$g^{rig} : (Y')^{rig} \to Y^{rig}$ devrait être un isomorphisme. Étant donnés
$\text{An}_S$ et le foncteur d'analytification, nous pouvons considérer la
condition selon laquelle le diagramme
$$
\xymatrix{
(V_{/Z})^{rig} \ar[rr]_{can} \ar[d]_{(g^{rig})^{-1} \circ \hat x^{an}} & &
(V \setminus Z)^{an} \ar[d]^{(u')^{an}} \\
(X'_{/T'})^{rig} \ar[rr]^{can} & & (X' \setminus T')^{an}
}
$$
soit commutatif. Cela a un sens, car
$g^{rig} : (X'_{T'})^{rig} \to W^{rig}$ est un isomorphisme et
$U' = X' \setminus T'$. Enfin, sous certaines hypothèses de fidélité du
foncteur d'analytification, cette condition sera équivalente à la condition de
compatibilité formulée ci-dessus. Nous espérons ainsi inciter le lecteur à
concevoir la compatibilité de $u'$ et $\hat x$ comme l'exigence que certaines
applications soient égales, plutôt que comme l'existence d'un certain
diagramme commutatif.
\end{remark}

\begin{lemma}
\label{lemma-functor}
Dans la situation \ref{situation-contractions}, la règle $F$ qui, à un schéma
localement noethérien $V$ sur $S$, associe l'ensemble des triplets
$(Z, u', \hat x)$ satisfaisant à la condition de compatibilité et qui, à un
morphisme $\varphi : V_2 \to V_1$ de schémas localement noethériens sur $S$,
associe l'application
$$
F(\varphi) : F(V_1) \longrightarrow F(V_2)
$$
envoyant un élément $(Z_1, u'_1, \hat x_1)$ de $F(V_1)$ sur
$(Z_2, u'_2, \hat x_2)$ dans $F(V_2)$, défini par
\begin{enumerate}
\item $Z_2 \subset V_2$ est l'image réciproque de $Z_1$ par $\varphi$ ;
\item $u'_2$ est la composée de $u'_1$ et de
$\varphi|_{V_2 \setminus Z_2} : V_2 \setminus Z_2 \to V_1 \setminus Z_1$,
\item $\hat x_2$ est la composée de $\hat x_1$ et de
$\varphi_{/Z_2} : V_{2, /Z_2} \to V_{1, /Z_1}$
\end{enumerate}
est un foncteur contravariant.
\end{lemma}

\begin{proof}
Pour vérifier la condition de compatibilité entre $u'_2$ et $\hat x_2$, soient
$V'_1 \to V_1$, $\hat x'_1$ et $x'_1$ des témoins de la compatibilité entre
$u'_1$ et $\hat x_1$. Posons $V'_2 = V_2 \times_{V_1} V'_1$, posons
$\hat x'_2$ égale à la composée de $\hat x'_1$ et de
$V'_{2, /Z_2} \to V'_{1, /Z_1}$, et posons $x'_2$
égale à la composée de $x'_1$ et de $V'_2 \to V'_1$.
Alors $V'_2 \to V_2$, $\hat x'_2$ et $x'_2$ témoignent de la compatibilité
entre $u'_2$ et $\hat x_2$. Nous omettons la vérification détaillée.
\end{proof}

\begin{lemma}
\label{lemma-solution}
Dans la situation \ref{situation-contractions}, s'il existe une solution
$(f : X' \to X, T, a)$, alors il existe une bijection fonctorielle
$F(V) = \Mor_S(V, X)$ sur la catégorie des schémas localement noethériens
$V$ sur $S$.
\end{lemma}

\begin{proof}
Soit $V$ un schéma localement noethérien sur $S$.
Soit $x : V \to X$ un morphisme sur $S$.
Nous obtenons alors un élément $(Z, u', \hat x)$ de $F(V)$ comme suit :
\begin{enumerate}
\item $Z \subset V$ est l'image réciproque de $T$ par $x$ ;
\item $u' : V \setminus Z \to U' = U$ est la restriction de
$x$ à $V \setminus Z$ ;
\item $\hat x : V_{/Z} \to W$ est la composée de
$x_{/Z} : V_{/Z} \to X_{/T}$ et de l'isomorphisme $a : X_{/T} \to W$.
\end{enumerate}
Ce triplet satisfait à la condition de compatibilité, car nous pouvons prendre
$V' = V \times_{x, X} X'$ et prendre pour $\hat x'$ le complété de la
projection $x' : V' \to X'$.

\medskip\noindent
Réciproquement, supposons donné un élément $(Z, u', \hat x)$ de $F(V)$. Nous
affirmons qu'il existe un unique morphisme $x : V \to X$ compatible avec $u'$
et $\hat x$. En effet, soient $V' \to V$, $\hat x'$ et $x'$ des témoins de la
compatibilité entre $u'$ et $\hat x$. La proposition
\ref{restricted-proposition-glue-modification}
d'Algébrisation des espaces formels est précisément le résultat dont nous avons
besoin pour trouver un unique morphisme $x : V \to X$ qui coïncide avec
$\hat x$ au-dessus de $V_{/Z}$ et avec $x'$ au-dessus de $V'$ (et, a fortiori,
avec $u'$ au-dessus de $V \setminus Z$).

\medskip\noindent
Nous omettons de vérifier que les deux constructions ci-dessus définissent des
bijections réciproques entre leurs domaines respectifs.
\end{proof}

\begin{lemma}
\label{lemma-functor-is-solution}
Dans la situation \ref{situation-contractions}, s'il existe un espace
algébrique $X$ localement de type fini sur $S$ et une bijection fonctorielle
$F(V) = \Mor_S(V, X)$ sur la catégorie des schémas localement noethériens
$V$ sur $S$, alors $X$ est une solution.
\end{lemma}

\begin{proof}
Nous devons construire un morphisme propre $f : X' \to X$, un sous-ensemble
fermé $T \subset |X|$ et un isomorphisme $a : X_{/T} \to W$ possédant les
propriétés (1), (2), (3) énumérées juste après la situation
\ref{situation-contractions}.

\medskip\noindent
La discussion de cette preuve est un peu pédante, car nous voulons faire
correspondre soigneusement les catégories sous-jacentes. Dans ce paragraphe,
nous expliquons comment le lecteur aventureux peut procéder avec moins de
timidité. Il peut en effet étendre notre définition du foncteur $F$ à tous les
espaces algébriques localement noethériens sur $S$. Il peut alors conclure que
$F$ et $X$ coïncident comme foncteurs sur la catégorie de ces espaces
algébriques, c'est-à-dire que $X$ représente $F$. On considère ensuite l'objet
universel $(T, u', \hat x)$ dans $F(X)$. Le lecteur constatera alors que, pour
le triplet $X'' \to X$, $\hat x'$, $x'$ témoignant de la compatibilité entre
$u'$ et $\hat x$, le morphisme $x' : X'' \to X'$ est un isomorphisme ; cela
fournira $f : X' \to X$ en inversant $x'$. Enfin, nous disposons déjà de
$T \subset |X|$, et le lecteur peut montrer que $\hat x$ est un isomorphisme
qui peut servir de dernier ingrédient, à savoir $a$.

\medskip\noindent
Notons $h_X(-) = \Mor_S(-, X)$ le foncteur des points de $X$ restreint à la
catégorie $(\textit{Noetherian}/S)_\etale$ de la section
\ref{section-noetherian}. D'après Limites d'espaces, remarque
\ref{spaces-limits-remark-limit-preserving}, les espaces algébriques $X$ et
$X'$ commutent aux limites. Il en va donc de même de leurs restrictions $h_X$
et $h_{X'}$. Pour construire $f$, il suffit donc de construire une
transformation $h_{X'} \to h_X = F$ ; voir le lemme
\ref{lemma-canonical-extension}. Soit donc $V \to S$ un objet de
$(\textit{Noetherian}/S)_\etale$, et soit $\tilde x : V \to X'$ un élément de
$h_{X'}(V)$. Nous obtenons alors un élément $(Z, u', \hat x)$ de $F(V)$ comme
suit :
\begin{enumerate}
\item $Z \subset V$ est l'image réciproque de $T'$ par $\tilde x$ ;
\item $u' : V \setminus Z \to U'$ est la restriction de
$\tilde x$ à $V \setminus Z$ ;
\item $\hat x : V_{/Z} \to W$ est la composée de
$x_{/Z} : V_{/Z} \to X'_{/T'}$ et de $g : X'_{/T'} \to W$.
\end{enumerate}
Ce triplet satisfait à la condition de compatibilité : tout d'abord, nous
obtenons toujours gratuitement $V' \to V$ et
$\hat x' : V'_{/Z'} \to X'_{/T'}$ (voir la discussion précédant le lemme
\ref{lemma-functor}). Il suffit ensuite de définir $x' : V' \to X'$ comme la
composée de $V' \to V$ et du morphisme $\tilde x : V \to X'$. Nous omettons
de vérifier que cela convient.

\medskip\noindent
Si $\xi : V \to X$ est un morphisme \'etale, où $V$ est un schéma, nous
obtenons $\xi = (Z, u', \hat x) \in F(V) = h_X(V) = X(V)$. Bien entendu, si
$\varphi : V' \to V$ est un autre morphisme \'etale de schémas, alors
$(Z, u', \hat x)$, ramené à $F(V')$, correspond à $\xi \circ \varphi$.
Le sous-ensemble fermé $T \subset |X|$ est simplement défini comme le
sous-ensemble fermé tel que $\xi : V \to X$, pour
$\xi = (Z, u', \hat x)$, ramène $T$ sur $Z$.

\medskip\noindent
Considérons des schémas noethériens $V$ sur $S$ et un morphisme
$\xi : V \to X$ correspondant à $(Z, u', \hat x)$ comme ci-dessus. Nous
voyons alors que $\xi(V)$ est ensemblistement contenu dans $T$ si et seulement
si $V = Z$ (comme espaces topologiques). Ainsi, $X_{/T}$ et $W$ coïncident
comme foncteurs. Cela fournit l'isomorphisme $a : X_{/T} \to W$. (Nous avons
omis ici un petit détail : pour les espaces algébriques formels localement
noethériens $X_{/T}$ et $W$, il suffit de vérifier que l'évaluation sur les
schémas localement noethériens sur $S$ donne un isomorphisme.)

\medskip\noindent
Nous omettons la preuve des conditions (1), (2) et (3).
\end{proof}

\begin{remark}
\label{remark-diagonal}
Dans la situation \ref{situation-contractions}, soit $V$ un schéma localement
noethérien sur $S$. Soient $(Z_i, u'_i, \hat x_i) \in F(V)$ pour $i = 1, 2$.
Soient $V'_i \to V$, $\hat x'_i$ et $x'_i$ des témoins de la compatibilité
entre $u'_i$ et $\hat x_i$ pour $i = 1, 2$.

\medskip\noindent
Posons $V' = V'_1 \times_V V'_2$. Notons $E' \to V'$ l'égalisateur des
morphismes
$$
V' \to V'_1 \xrightarrow{x'_1} X'
\quad\text{et}\quad
V' \to V'_2 \xrightarrow{x'_2} X'
$$
Posons $Z = Z_1 \cap Z_2$. Soit $E_W \to V_{/Z}$ l'égalisateur des morphismes
$$
V_{/Z} \to V_{/Z_1} \xrightarrow{\hat x_1} W
\quad\text{et}\quad
V_{/Z} \to V_{/Z_2} \xrightarrow{\hat x_2} W
$$
Remarquons que $E' \to V$ est séparé et localement de type fini, et que $E_W$
est un espace algébrique formel localement noethérien séparé sur $V$. Les
compatibilités entre les différents morphismes en jeu montrent que
\begin{enumerate}
\item $\Im(E' \to V) \cap (Z_1 \cup Z_2)$
est contenu dans $Z = Z_1 \cap Z_2$ ;
\item le morphisme $E' \times_V (V \setminus Z) \to V \setminus Z$
est un monomorphisme et est égal à l'égalisateur des restrictions de $u'_1$ et
$u'_2$ à $V \setminus (Z_1 \cup Z_2)$ ;
\item le morphisme $E'_{/Z} \to V_{/Z}$ se factorise par $E_W$, et le diagramme
$$
\xymatrix{
E'_{/Z} \ar[r] \ar[d] & X'_{/T'} \ar[d]^g \\
E_W \ar[r] & W
}
$$
est cartésien. En particulier, le morphisme $E'_{/Z} \to E_W$ est une
modification formelle en tant que changement de base de $g$ ;
\item $E'$, $(E' \to V)^{-1}Z$ et $E'_{/Z} \to E_W$
constituent un triplet comme dans la situation \ref{situation-contractions},
ayant pour schéma de base le schéma localement noethérien $V$ ;
\item étant donné un morphisme $\varphi : A \to V$ de schémas localement
noethériens, les assertions suivantes sont équivalentes :
\begin{enumerate}
\item $(Z_1, u'_1, \hat x_1)$ et $(Z_2, u'_2, \hat x_2)$
se restreignent au même élément de $F(A)$ ;
\item $A \setminus \varphi^{-1}(Z) \to V \setminus Z$
se factorise par $E' \times_V (V \setminus Z)$, et
$A_{/\varphi^{-1}(Z)} \to V_{/Z}$ se factorise par $E_W$.
\end{enumerate}
\end{enumerate}
Nous concluons, à l'aide des lemmes \ref{lemma-solution} et
\ref{lemma-functor-is-solution}, que s'il existe une solution $E \to V$ pour le
triplet de (4), alors $E$ représente $F \times_{\Delta, F \times F} V$ sur la
catégorie des schémas localement noethériens sur $V$.
\end{remark}

\begin{lemma}
\label{lemma-closed-immersion}
Dans la situation \ref{situation-contractions}, supposons donné un
sous-ensemble fermé $Z \subset S$ tel que
\begin{enumerate}
\item l'image réciproque de $Z$ dans $X'$ soit $T'$ ;
\item $U' \to S \setminus Z$ soit une immersion fermée ;
\item $W \to S_{/Z}$ soit une immersion fermée.
\end{enumerate}
Alors il existe une solution $(f : X' \to X, T, a)$ et, de plus,
$X \to S$ est une immersion fermée.
\end{lemma}

\begin{proof}
Supposons disposer d'un sous-schéma fermé $X \subset S$ tel que
$X \cap (S \setminus Z) = U'$ et $X_{/Z} = W$. Alors $X$ représente le
foncteur $F$ (nous omettons certains détails) et constitue donc une solution.
La recherche de $X$ est manifestement une question locale sur $S$. Nous nous
ramenons ainsi au cas traité dans le paragraphe suivant.

\medskip\noindent
Supposons $S = \Spec(A)$ affine. Soit $I \subset A$ l'idéal radical définissant
$Z$. Écrivons $I = (f_1, \ldots, f_r)$. Par hypothèse, nous disposons de
\begin{enumerate}
\item l'immersion fermée $U' \to S \setminus Z$, qui détermine des idéaux
$J_i \subset A[1/f_i]$ tels que $J_i$ et $J_j$ engendrent le même idéal dans
$A[1/f_if_j]$ ;
\item l'immersion fermée $W \to S_{/Z}$, qui est l'application
$\text{Spf}(A^\wedge/J') \to \text{Spf}(A^\wedge)$
pour un certain idéal $J' \subset A^\wedge$ dans le complété $I$-adique
$A^\wedge$ de $A$.
\end{enumerate}
Pour achever la preuve, nous devons trouver un idéal $J \subset A$ tel que
$J_i = J[1/f_i]$ et $J' = JA^\wedge$. D'après Compléments d'algèbre,
proposition \ref{more-algebra-proposition-equivalence}, il suffit de montrer
que $J_i$ et $J'$ engendrent le même idéal dans $A^\wedge[1/f_i]$ pour tout
$i$.

\medskip\noindent
Rappelons que $A' = H^0(X', \mathcal{O})$ est une $A$-algèbre finie dont la
formation commute aux changements de base plats (Cohomologie des espaces,
lemmes
\ref{spaces-cohomology-lemma-proper-over-affine-cohomology-finite} et
\ref{spaces-cohomology-lemma-flat-base-change-cohomology}). Notons
$J'' = \Ker(A \to A')$\footnote{Contrairement à ce que pourrait attendre le
lecteur, les idéaux $J$ et $J''$ ne coïncident pas en général.}.
Nous avons $J_i = J''A[1/f_i]$, comme il résulte du changement de base vers le
spectre de $A[1/f_i]$. Remarquons que nous avons un diagramme commutatif
$$
\xymatrix{
X' \ar[d] &
X'_{/T'} \times_{S_{/Z}} \text{Spf}(A^\wedge) \ar[l] \ar[d] &
X'_{/T'} \times_W \text{Spf}(A^\wedge/J') \ar@{=}[l] \ar[d] \\
\Spec(A) &
\text{Spf}(A^\wedge) \ar[l] &
\text{Spf}(A^\wedge/J') \ar[l]
}
$$
La flèche verticale médiane est le complété de la flèche verticale gauche le
long des sous-ensembles fermés évidents. D'après le théorème des fonctions
formelles, nous avons
$$
(A')^\wedge = \Gamma(X' \times_S \Spec(A^\wedge), \mathcal{O}) =
\lim H^0(X' \times_S \Spec(A/I^n), \mathcal{O})
$$
Voir Cohomologie des espaces, théorème
\ref{spaces-cohomology-theorem-formal-functions}.
Le diagramme montre que $J'$ s'envoie sur zéro dans $(A')^\wedge$. Ainsi,
$J' \subset J'' A^\wedge$. Considérons les flèches
$$
X'_{/T'} \to
\text{Spf}(A^\wedge/J''A^\wedge) \to
\text{Spf}(A^\wedge/J') = W
$$
Nous savons que la composée $g$ est une modification formelle (en particulier
rig-\'etale et rig-surjective) et que la seconde flèche est une immersion
fermée (en particulier un monomorphisme adique). Par conséquent,
$X'_{/T'} \to \text{Spf}(A^\wedge/J''A^\wedge)$ est rig-surjectif et
rig-\'etale ; voir Algébrisation des espaces formels, lemmes
\ref{restricted-lemma-rig-surjective-alternative-permanence} et
\ref{restricted-lemma-rig-etale-alternative-permanence}.
En appliquant Algébrisation des espaces formels, lemmes
\ref{restricted-lemma-rig-etale-descent} et
\ref{restricted-lemma-permanence-rig-surjective}
nous concluons que $\text{Spf}(A^\wedge/J''A^\wedge) \to W$ est rig-\'etale et
rig-surjectif. D'après Algébrisation des espaces formels, lemme
\ref{restricted-lemma-closed-immersion-rig-smooth-rig-surjective}
nous concluons que $I^n J'' A^\wedge \subset J'$ pour un certain $n > 0$.
Il s'ensuit que $J'' A^\wedge[1/f_i] = J' A^\wedge[1/f_i]$, et nous en
déduisons $J_i A^\wedge[1/f_i] = J' A^\wedge[1/f_i]$ pour tout $i$, comme
souhaité.
\end{proof}

\begin{lemma}
\label{lemma-diagonal-contractions}
Dans la situation \ref{situation-contractions}, supposons que $X' \to S$ et
$W \to S$ soient séparés. Alors la diagonale $\Delta : F \to F \times F$ est
représentable par des immersions fermées.
\end{lemma}

\begin{proof}
Combiner le lemme \ref{lemma-closed-immersion} avec la discussion de la
remarque \ref{remark-diagonal}.
\end{proof}

\begin{lemma}
\label{lemma-sheaf}
Dans la situation \ref{situation-contractions}, le foncteur $F$ possède la
propriété de faisceau pour tous les recouvrements \'etales de schémas
localement noethériens sur $S$.
\end{lemma}

\begin{proof}
Omis. Indication : les morphismes peuvent être définis localement pour la
topologie \'etale.
\end{proof}

\begin{lemma}
\label{lemma-limit-preserving}
Dans la situation \ref{situation-contractions}, le foncteur $F$ commute aux
limites : pour toute limite dirigée $V = \lim V_\lambda$ de schémas affines
noethériens sur $S$, nous avons $F(V) = \colim F(V_\lambda)$.
\end{lemma}

\begin{proof}
Cette preuve est d'une longueur absurde. Elle consiste en grande partie en des
arguments classiques sur les limites et la localisation \'etale. Nous
conseillons vivement au lecteur de passer directement à la dernière partie de
la preuve, où il se produit quelque chose d'intéressant.

\medskip\noindent
Soit $V = \lim_{\lambda \in \Lambda} V_i$ une limite dirigée de schémas sur
$S$, où $V$ et les $V_\lambda$ sont noethériens et où les morphismes de
transition sont affines. Voir Limites, section \ref{limits-section-limits},
pour des résultats sur les limites de schémas. Nous allons démontrer que
$\colim F(V_\lambda) \to F(V)$ est bijective.

\medskip\noindent
Preuve de l'injectivité : notations.
Soit $\lambda \in \Lambda$, et soient
$\xi_{\lambda, 1}, \xi_{\lambda, 2} \in F(V_\lambda)$ des éléments qui se
restreignent au même élément de $F(V)$. Écrivons
$\xi_{\lambda, 1} = (Z_{\lambda, 1}, u'_{\lambda, 1},
\hat x_{\lambda, 1})$ et
$\xi_{\lambda, 2} = (Z_{\lambda, 2}, u'_{\lambda, 2}, \hat x_{\lambda, 2})$.

\medskip\noindent
Preuve de l'injectivité : coïncidence des $Z_{\lambda, i}$.
Puisque $Z_{\lambda, 1}$ et $Z_{\lambda, 2}$ se restreignent au même
sous-ensemble fermé de $V$, nous pouvons, après avoir augmenté $i$, supposer
que $Z_{\lambda, 1} = Z_{\lambda, 2}$ ; voir Limites, lemme
\ref{limits-lemma-inverse-limit-top}, et Topologie, lemme
\ref{topology-lemma-describe-limits}. Notons $Z_\lambda \subset V_\lambda$
leur valeur commune ; pour $\mu \geq \lambda$, notons
$Z_\mu \subset V_\mu$ son image réciproque dans $V_\mu$, et notons $Z$ son
image réciproque dans $V$. Nous utiliserons plus bas le fait que
$Z = \lim_{\mu \geq \lambda} Z_\mu$ comme schémas, si nous considérons $Z$ et
les $Z_\mu$ comme des sous-schémas fermés réduits.

\medskip\noindent
Preuve de l'injectivité : coïncidence des $u'_{\lambda, i}$.
Puisque $U'$ est localement de type fini sur $S$ et que les restrictions de
$u'_{\lambda, 1}$ et $u'_{\lambda, 2}$ à $V \setminus Z$ sont égales, nous
pouvons, après avoir augmenté $\lambda$, supposer que
$u'_{\lambda, 1} = u'_{\lambda, 2}$ ; voir Limites, proposition
\ref{limits-proposition-characterize-locally-finite-presentation}. Notons
$u'_\lambda$ leur valeur commune et $u'$ la restriction à $V \setminus Z$.

\medskip\noindent
Preuve de l'injectivité : reformulation.
À ce stade, nous avons
$\xi_{\lambda, 1} = (Z_\lambda, u'_\lambda, \hat x_{\lambda, 1})$ et
$\xi_{\lambda, 2} = (Z_\lambda, u'_\lambda, \hat x_{\lambda, 2})$.
Le principal problème de cette partie de la preuve est de montrer que les
morphismes $\hat x_{\lambda, 1}$ et $\hat x_{\lambda, 2}$ deviennent égaux
après avoir augmenté $\lambda$.

\medskip\noindent
Preuve de l'injectivité : coïncidence des
$\hat x_{\lambda, i}|_{Z_\lambda}$.
Considérons les morphismes
$\hat x_{\lambda, 1}|_{Z_\lambda}, \hat x_{\lambda, 2}|_{Z_\lambda} :
Z_\lambda \to W_{red}$.
Ces morphismes se restreignent au même morphisme $Z \to W_{red}$. Puisque
$W_{red}$ est un schéma localement de type fini sur $S$, la proposition
\ref{limits-proposition-characterize-locally-finite-presentation} de Limites
montre qu'après avoir remplacé $\lambda$ par un indice plus grand, nous pouvons
supposer
$\hat x_{\lambda, 1}|_{Z_\lambda} = \hat x_{\lambda, 2}|_{Z_\lambda} :
Z_\lambda \to W_{red}$.

\medskip\noindent
Preuve de l'injectivité : fin.
Nous allons maintenant appliquer la discussion de la remarque
\ref{remark-diagonal} à $V_\lambda$ et aux deux éléments
$\xi_{\lambda, 1}, \xi_{\lambda, 2} \in F(V_\lambda)$. Nous obtenons
\begin{enumerate}
\item $e_\lambda : E_\lambda' \to V_\lambda$
séparé et localement de type fini ;
\item $e_\lambda^{-1}(V_\lambda \setminus Z_\lambda) \to
V_\lambda \setminus Z_\lambda$ qui est un isomorphisme ;
\item un monomorphisme $E_{W, \lambda} \to V_{\lambda, /Z_\lambda}$ qui est
l'égalisateur de $\hat x_{\lambda, 1}$ et $\hat x_{\lambda, 2}$ ;
\item une modification formelle
$E'_{\lambda, /Z_\lambda} \to E_{W, \lambda}$.
\end{enumerate}
L'assertion (2) résulte de l'assertion (2) de la remarque
\ref{remark-diagonal} et du travail préparatoire ci-dessus, qui a donné
$u'_{\lambda, 1} = u'_{\lambda, 2} = u'_\lambda$.
Puisque $Z_\lambda = (V_{\lambda, /Z_\lambda})_{red}$ se factorise par
$E_{W, \lambda}$ en raison de l'égalité
$\hat x_{\lambda, 1}|_{Z_\lambda} = \hat x_{\lambda, 2}|_{Z_\lambda}$
le lemme \ref{formal-spaces-lemma-monomorphism-iso-over-red} d'Espaces formels
montre que $E_{W, \lambda} \to V_{\lambda, /Z_\lambda}$ est une immersion
fermée. L'assertion (4) de la remarque \ref{remark-diagonal} et le lemme
\ref{lemma-closed-immersion}, appliqué au triplet
$E_\lambda'$, $e_\lambda^{-1}(Z_\lambda)$,
$E'_{\lambda, /Z_\lambda} \to E_{W, \lambda}$ sur $V_\lambda$, montrent qu'il
existe une immersion fermée $E_\lambda \to V_\lambda$ qui est une solution
pour ce triplet. Utilisons maintenant l'assertion (5) de la remarque
\ref{remark-diagonal}, qui, combinée au lemme \ref{lemma-solution}, affirme que
$E_\lambda$ est l'« égalisateur » de $\xi_{\lambda, 1}$ et de
$\xi_{\lambda, 2}$. En particulier, $V \to V_\lambda$ se factorise par
$E_\lambda$. En utilisant une fois encore Limites, proposition
\ref{limits-proposition-characterize-locally-finite-presentation}, nous
trouvons alors $\mu \geq \lambda$ tel que $V_\mu \to V_\lambda$ se factorise
par $E_\lambda$ ; les images réciproques de $\xi_{\lambda, 1}$ et de
$\xi_{\lambda, 2}$ dans $V_\mu$ sont donc égales, comme souhaité.

\medskip\noindent
Preuve de la surjectivité : énoncé.
Soit $\xi = (Z, u', \hat x)$ un élément de $F(V)$. Nous devons trouver un
$\lambda \in \Lambda$ et un élément $\xi_\lambda \in F(V_\lambda)$ se
restreignant à $\xi$.

\medskip\noindent
Preuve de la surjectivité : la question est locale pour la topologie \'etale.
Par l'unicité démontrée dans la partie précédente de la preuve et par la
propriété de faisceau de $F$ dans le lemme \ref{lemma-sheaf}, le problème est
local sur $V$ pour la topologie \'etale. Plus précisément, soit $v \in V$.
Nous affirmons qu'il suffit de trouver un morphisme \'etale
$(\tilde V, \tilde v) \to (V, v)$, un
$\lambda$, un morphisme \'etale $\tilde V_\lambda \to V_\lambda$ et un élément
$\tilde \xi_\lambda \in F(\tilde V_\lambda)$ tels que
$\tilde V = \tilde V_\lambda \times_{V_\lambda} V$
et $\xi|_U = \tilde \xi_\lambda|_U$. Nous omettons une preuve détaillée de
cette affirmation\footnote{Pour la démontrer, on rassemble une famille de
morphismes $\tilde V \to V$ en un recouvrement \'etale fini et on montre que
les morphismes correspondants $\tilde V_\lambda \to V_\lambda$ forment aussi
un recouvrement \'etale (quitte à augmenter $\lambda$). On utilise ensuite
l'injectivité pour voir que les éléments $\tilde \xi_\lambda$ se recollent
(quitte à augmenter $\lambda$), puis la propriété de faisceau de $F$ pour
descendre ces éléments en un élément de $F(V_\lambda)$.}.

\medskip\noindent
Preuve de la surjectivité : reformulation du problème.
Rappelons que tout morphisme \'etale $(\tilde V, \tilde v) \to (V, v)$, avec
$\tilde V$ affine, est le changement de base d'un morphisme \'etale
$\tilde V_\lambda \to V_\lambda$, où $\tilde V_\lambda$ est affine, pour un
certain $\lambda$ ; voir par exemple Topologies, lemme
\ref{topologies-lemma-limit-fppf-topology}. Étant donné $\tilde V_\lambda$,
nous avons
$\tilde V = \lim_{\mu \geq \lambda} \tilde V_\lambda \times_{V_\lambda} V_\mu$.
Ainsi, étant donné $(\tilde V, \tilde v) \to (V, v)$ \'etale avec $\tilde V$
affine, nous pouvons remplacer $(V, v)$ par $(\tilde V, \tilde v)$ et $\xi$
par la restriction de $\xi$ à $\tilde V$.

\medskip\noindent
Preuve de la surjectivité : réduction à une base affine. En particulier,
supposons que $\tilde S \subset S$ soit un sous-schéma ouvert affine tel que
$v \in V$ s'envoie sur un point de $\tilde S$. D'après le paragraphe
précédent, nous pouvons alors remplacer $V$ par
$\tilde V = \tilde S \times_S V$. Dans ce cas, il suffit
de résoudre le problème pour le foncteur $F$ restreint à la catégorie des
schémas localement noethériens sur $\tilde S$. Ce foncteur est précisément
celui associé à la situation tout entière après changement de base à
$\tilde S$. Pour le reste de la preuve, nous pouvons et allons donc supposer
que $S = \Spec(R)$ est un schéma affine noethérien.

\medskip\noindent
Preuve de la surjectivité : cas facile.
Si $v \in V \setminus Z$, nous pouvons prendre $\tilde V = V \setminus Z$.
Cet ouvert descend en un sous-schéma ouvert
$\tilde V_\lambda \subset V_\lambda$ pour un certain $\lambda$, d'après
Limites, lemme \ref{limits-lemma-descend-opens}. Ensuite, quitte à augmenter
$\lambda$, nous pouvons supposer qu'il existe un morphisme
$u'_\lambda : \tilde V_\lambda \to U'$ se restreignant à $u'$. Le choix
$\tilde \xi_\lambda = (\emptyset, u'_\lambda, \emptyset)$
donne l'élément souhaité de $F(\tilde V_\lambda)$.

\medskip\noindent
Preuve de la surjectivité : cas difficile et réduction au cas où $W$ est affine.
Le cas le plus difficile est celui où $v \in Z \subset V$.
Nous affirmons que nous pouvons nous ramener au cas où $W$ est un
schéma formel affine ; nous invitons le lecteur à sauter cet argument\footnote{La
méthode d'Artin pour démontrer ce lemme consiste à contourner ce point, ce qui
lui permet d'éviter de démontrer d'abord l'injectivité. Plus précisément, Artin
travaille systématiquement avec des recouvrements \'etales affines finis de tous
les espaces en jeu, en conservant pendant la preuve les applications entre eux.
Rétrospectivement, cette méthode serait peut-être préférable à la nôtre.}.
En effet, nous pouvons choisir un morphisme \'etale
$\tilde W \to W$, où $\tilde W$ est un espace algébrique formel affine,
tel que l'image de $v$ par $\hat x : V_{/Z} \to W$ appartienne
à l'image de $\tilde W \to W$ (sur les réductions).
Les morphismes
$$
p : \tilde W \times_{W, g} X'_{/T'} \longrightarrow X'_{/T'}
$$
et
$$
q : \tilde W \times_{W, \hat x} V_{/Z} \to V_{/Z}
$$
sont des morphismes \'etales d'espaces algébriques formels localement
noethériens. Par un cas facile du théorème d'Algébrisation des espaces formels
\ref{restricted-theorem-dilatations-general}
il existe un morphisme d'espaces algébriques $\tilde X' \to X'$
qui est localement de type fini, est un isomorphisme au-dessus de $U'$ et
tel que $\tilde X'_{/T'} \to X'_{/T'}$ soit isomorphe à $p$.
D'après Algébrisation des espaces formels, lemme
\ref{restricted-lemma-output-etale}, le morphisme $\tilde X' \to X'$ est
\'etale. Notons $\tilde T' \subset |\tilde X'|$ l'image inverse de $T'$.
Notons $\tilde U' \subset \tilde X'$ le sous-espace ouvert complémentaire.
Notons $\tilde g' : \tilde X'_{/\tilde T'} \to \tilde W$
la modification formelle obtenue de $g$ par le changement de base
$\tilde W \to W$. Nous voyons alors que
$$
\tilde X',\ \tilde T',\ \tilde U',\ \tilde W,
\ \tilde g : \tilde X'_{/\tilde T'} \to \tilde W
$$
est un autre exemple de la situation \ref{situation-contractions}.
Notons $\tilde F$ le foncteur construit à partir de ce triplet.
Il existe une transformation de foncteurs
$$
\tilde F \longrightarrow F
$$
construite de manière évidente à l'aide des morphismes $\tilde X' \to X'$ et
$\tilde W \to W$ ; les détails sont omis.

\medskip\noindent
Preuve de la surjectivité : cas difficile et réduction au cas où $W$ est
affine, deuxième partie. D'après le même théorème que ci-dessus, il existe un
morphisme d'espaces algébriques $\tilde V \to V$ qui est localement de type
fini, est un isomorphisme au-dessus de $V \setminus Z$ et tel que
$\tilde V_{/Z} \to V_{/Z}$ soit isomorphe à $q$.
Notons $\tilde Z \subset \tilde V$ l'image inverse de $Z$.
D'après Algébrisation des espaces formels, lemmes
\ref{restricted-lemma-output-etale} et \ref{restricted-lemma-output-separated},
le morphisme $\tilde V \to V$ est \'etale et séparé.
En particulier, $\tilde V$ est un schéma (localement noethérien) ; voir par
exemple Morphismes d'espaces, proposition
\ref{spaces-morphisms-proposition-locally-quasi-finite-separated-over-scheme}.
Nous disposons du morphisme $u'$, que nous pouvons considérer comme un morphisme
$$
\tilde u' : \tilde V \setminus \tilde Z \longrightarrow \tilde U'
$$
où $\tilde U' \subset \tilde X'$ est l'ouvert qui s'envoie isomorphiquement
sur $U'$. Nous disposons d'un morphisme
$$
\tilde {\hat x} :
\tilde V_{/\tilde Z} = \tilde W \times_{W, \hat x} V_{/Z}
\longrightarrow
\tilde W
$$
Il s'agit simplement ici de la projection. Nous obtenons donc le triplet
$$
\tilde \xi = (\tilde Z, \tilde u', \tilde {\hat x}) \in \tilde F(\tilde V)
$$
Nous omettons de démontrer la condition de compatibilité ; voici quelques
indications : si $V' \to V$, $\hat x'$ et $x'$ attestent la compatibilité entre
$u'$ et $\hat x$, on pose $\tilde V' = V' \times_V \tilde V$, qui est muni des
morphismes $\tilde{\hat x}'$ et $\tilde x'$, et on vérifie que cela convient.
L'image de $\tilde \xi$ par la transformation $\tilde F \to F$
est la restriction de $\xi$ à $\tilde V$.

\medskip\noindent
Preuve de la surjectivité : cas difficile et réduction au cas où $W$ est
affine, troisième partie. Par notre choix de $\tilde W \to W$, il existe un
ouvert affine $\tilde V_{open} \subset \tilde V$ (les notations commencent à
manquer) dont l'image dans $V$ contient le point choisi $v \in V$.
D'après le cas étudié au paragraphe suivant et les remarques faites plus haut,
nous pouvons alors faire descendre $\tilde \xi|_{\tilde V_{open}}$
en un élément $\tilde \xi_\lambda$ de $\tilde F$ au-dessus de
$\tilde V_{\lambda, open}$, pour un certain morphisme \'etale
$\tilde V_{\lambda, open} \to V_\lambda$
dont le changement de base à $V$ est $\tilde V_{open}$.
En appliquant la transformation de foncteurs $\tilde F \to F$, nous obtenons
l'élément de $F(\tilde V_{\lambda, open})$ que nous cherchions.
Nous sommes ainsi ramenés au cas traité au paragraphe suivant.

\medskip\noindent
Preuve de la surjectivité : cas où $W$ est affine. Nous avons $v \in Z \subset V$
et $W$ est un espace algébrique formel affine. Rappelons que
$$
\xi = (Z, u', \hat x) \in F(V)
$$
Nous pouvons encore remplacer $V$ par un voisinage \'etale de $v$.
En particulier, nous pouvons et allons supposer $V$ et $V_\lambda$ affines.

\medskip\noindent
Preuve de la surjectivité : descente de $Z$. Nous pouvons trouver un $\lambda$
et un sous-schéma fermé $Z_\lambda \subset V_\lambda$ tels que
$Z$ soit le changement de base de $Z_\lambda$ à $V$. Voir
Limites, lemme \ref{limits-lemma-descend-finite-presentation}.
Attention : nous ne savons pas (et, en général, ce ne sera pas vrai)
que $Z_\lambda$ soit un sous-schéma fermé réduit de $V_\lambda$.
Pour $\mu \geq \lambda$, notons $Z_\mu \subset V_\mu$ le sous-schéma
image inverse dans $V_\mu$. Nous utiliserons plus loin que
$Z = \lim_{\mu \geq \lambda} Z_\mu$ en tant que schémas.

\medskip\noindent
Preuve de la surjectivité : descente de $u'$.
Puisque $U'$ est localement de type fini sur $S$,
nous pouvons, quitte à augmenter $\lambda$, supposer
qu'il existe un morphisme
$u'_\lambda : V_\lambda \setminus Z_\lambda \to U'$
dont la restriction à $V \setminus Z$ est $u'$.
Voir Limites, proposition
\ref{limits-proposition-characterize-locally-finite-presentation}.
Pour $\mu \geq \lambda$, nous noterons $u'_\mu$ la restriction
de $u'_\lambda$ à $V_\mu \setminus Z_\mu$.

\medskip\noindent
Preuve de la surjectivité : descente d'un témoin.
Soient $V' \to V$, $\hat x'$ et $x'$ des données attestant la compatibilité entre
$u'$ et $\hat x$. En utilisant les mêmes références que ci-dessus, nous pouvons
supposer (quitte à augmenter $\lambda$) qu'il existe un morphisme
$V'_\lambda \to V_\lambda$ de type fini dont le changement de base à $V$
est $V' \to V$. Quitte à augmenter $\lambda$, nous pouvons supposer que
$V'_\lambda \to V_\lambda$ est propre
(Limites, lemme \ref{limits-lemma-eventually-proper}).
Ensuite, nous pouvons supposer que $V'_\lambda \to V_\lambda$ est un isomorphisme
au-dessus de $V_\lambda \setminus Z_\lambda$
(Limites, lemme \ref{limits-lemma-descend-isomorphism}).
Ensuite, nous pouvons supposer qu'il existe un morphisme
$x'_\lambda : V'_\lambda \to X'$ dont la restriction à $V'$ est $x'$.
En augmentant encore $\lambda$, nous pouvons supposer que $x'_\lambda$
coïncide avec $u'_\lambda$ au-dessus de $V_\lambda \setminus Z_\lambda$.
Pour $\mu \geq \lambda$, nous notons $V'_\mu$ et $x'_\mu$ le changement de
base de $V'_\lambda$ et la restriction de $x'_\lambda$, respectivement.

\medskip\noindent
Preuve de la surjectivité : algèbre.
Écrivons $W = \text{Spf}(B)$, $V = \Spec(A)$ et,
pour $\mu \geq \lambda$, écrivons $V_\mu = \Spec(A_\mu)$.
Notons $I_\mu \subset A_\mu$ et $I \subset A$
les idéaux définissant $Z_\mu$ et $Z$.
Nous avons alors $I_\lambda A_\mu = I_\mu$ et $I_\lambda A = I$.
Le morphisme $\hat x$ détermine, et est déterminé par, un
homomorphisme continu d'anneaux
$$
(\hat x)^\sharp : B \longrightarrow A^\wedge
$$
où $A^\wedge$ est le complété $I$-adique de $A$.
Pour achever la preuve, nous devons montrer que cet homomorphisme descend en un
homomorphisme à valeurs dans $A_\mu^\wedge$ pour un $\mu$ suffisamment grand,
où $A_\mu^\wedge$ est le complété $I_\mu$-adique de $A_\mu$.
Ce fait n'est pas trivial ; Artin écrit dans son article
\cite{ArtinII} : « Puisque les données (3.5) font intervenir des complétés
$I$-adiques, qui ne commutent pas aux limites inductives, la vérification est
assez délicate. C'est l'analogue algébrique d'une preuve de convergence en analyse. »

\medskip\noindent
Preuve de la surjectivité : algèbre, d'autres anneaux.
Posons
$$
C_\mu = \Gamma(V'_\mu, \mathcal{O})
\quad\text{et}\quad
C = \Gamma(V', \mathcal{O})
$$
Observons que les homomorphismes d'anneaux $A \to C$ et
$A_\mu \to C_\mu$ sont finis, puisque $V' \to V$ et
$V'_\mu \to V_\mu$ sont des morphismes propres ; voir Cohomologie des espaces,
lemme
\ref{spaces-cohomology-lemma-proper-over-affine-cohomology-finite}.
Puisque $V = \lim V_\mu$ et $V' = \lim V'_\mu$,
nous avons
$$
A = \colim A_\mu
\quad\text{et}\quad
C = \colim C_\mu
$$
d'après Limites, lemme \ref{limits-lemma-descend-section}\footnote{Nous ne
savons pas que $C_\mu = C_\lambda \otimes_{A_\lambda} A_\mu$, car les
divers morphismes ne sont pas plats.}. Pour un élément
$a \in I$, resp.\ $a \in I_\mu$, les homomorphismes $A_a \to C_a$,
resp.\ $(A_\mu)_a \to (C_\mu)_a$, sont des isomorphismes par changement de
base plat (Cohomologie des espaces, lemme
\ref{spaces-cohomology-lemma-flat-base-change-cohomology}).
Par conséquent, le noyau et le conoyau de $A \to C$ ont leur support contenu
dans $V(I)$, et il en va de même pour $A_\mu \to C_\mu$.
Nous en concluons que le noyau et le conoyau de $A \to C$
sont annulés par une puissance de $I$, et que le noyau et le conoyau de
$A_\mu \to C_\mu$ sont annulés par une puissance de $I_\mu$ ; voir
Algèbre, lemme \ref{algebra-lemma-Noetherian-power-ideal-kills-module}.

\medskip\noindent
Preuve de la surjectivité : algèbre, d'autres homomorphismes d'anneaux.
Notons $Z_n \subset V$ le $n$-ième voisinage infinitésimal
de $Z$, et notons $Z_{\mu, n} \subset V_\mu$
le $n$-ième voisinage infinitésimal de $Z_\mu$.
Par le théorème des fonctions formelles
(Cohomologie des espaces, théorème
\ref{spaces-cohomology-theorem-formal-functions})
nous avons
$$
C^\wedge = \lim_n H^0(V' \times_V Z_n, \mathcal{O})
\quad\text{et}\quad
C_\mu^\wedge =
\lim_n H^0(V'_\mu \times_{V_\mu} Z_{\mu, n}, \mathcal{O})
$$
où $C^\wedge$ et $C_\mu^\wedge$ sont les complétés respectifs pour
$I$ et $I_\mu$.
En composant le complété du morphisme
$x'_\mu : V'_\mu \to X'$ avec le morphisme $g : X'_{/T'} \to W$, nous obtenons
$$
g \circ x'_{\mu, /Z_\mu} :
V'_{\mu, /Z_\mu} = \colim V_\mu' \times_{V_\mu} Z_{\mu, n}
\longrightarrow
W
$$
et donc, d'après la description ci-dessus du complété
$C_\mu^\wedge$, nous obtenons un homomorphisme continu d'anneaux
$$
(g \circ x'_{\mu, /Z_\mu})^\sharp : B \longrightarrow C_\mu^\wedge
$$
Le fait que $V' \to V$, $\hat x'$ et $x'$ attestent la compatibilité
entre $u'$ et $\hat x$ entraîne la
commutativité du diagramme suivant
$$
\xymatrix{
C_\mu^\wedge \ar[r] &
C^\wedge \\
B \ar[u]^{(g \circ x'_{\mu, /Z_\mu})^\sharp} \ar[r]^{(\hat x)^\sharp} &
A^\wedge \ar[u]
}
$$

\medskip\noindent
Preuve de la surjectivité : autres arguments algébriques. Rappelons que les
$A$-modules de type fini $\Ker(A \to C)$ et $\Coker(A \to C)$ sont annulés
par une puissance de $I$ et que, de même, les $A_\mu$-modules de type fini
$\Ker(A_\mu \to C_\mu)$ et $\Coker(A_\mu \to C_\mu)$ sont annulés
par une puissance de $I_\mu$. Ces modules sont donc égaux
à leurs complétés. Puisque le passage au complété $I$-adique est exact dans
la catégorie des $A$-modules de type fini (voir
Algèbre, section \ref{algebra-section-completion-noetherian}),
il s'ensuit que
$$
\Coker(A^\wedge \to C^\wedge) = \Coker(A \to C)
$$
et il en va de même pour les noyaux et pour les homomorphismes
$A_\mu \to C_\mu$. Bien entendu, nous avons aussi
$$
\Ker(A \to C) = \colim \Ker(A_\mu \to C_\mu)
\quad\text{et}\quad
\Coker(A \to C) = \colim \Coker(A_\mu \to C_\mu)
$$
Rappelons que $S = \Spec(R)$ est affine. Tous les homomorphismes d'anneaux
ci-dessus sont des homomorphismes de $R$-algèbres, puisque tous les morphismes
sont des morphismes au-dessus de $S$. D'après
Algébrisation des espaces formels, lemme
\ref{restricted-lemma-Noetherian-finite-type-red}
nous voyons que $B$ est topologiquement de type fini sur $R$.
L'algèbre $B$ est, disons, topologiquement engendrée par $b_1, \ldots, b_n$.
Choisissons un $\mu$ (par exemple $\lambda$) et considérons
les éléments
$$
\text{images de }
(g \circ x'_{\mu, /Z_\mu})^\sharp(b_1)
, \ldots,
(g \circ x'_{\mu, /Z_\mu})^\sharp(b_n)
\text{ dans }\Coker(A_\mu \to C_\mu)
$$
L'image de ces éléments dans $\Coker(\alpha)$ est nulle
par commutativité du carré ci-dessus. Puisque
$\Coker(A \to C) = \colim \Coker(A_\mu \to C_\mu)$ et que ces
conoyaux sont égaux à leurs complétés,
nous voyons que, quitte à augmenter $\mu$, nous pouvons supposer que toutes ces
images sont nulles. Le homomorphisme continu
$(g \circ x'_{\mu, /Z_\mu})^\sharp$ a donc son image contenue
dans $\Im(A_\mu \to C_\mu)$.
Choisissons des éléments $a_{\mu, j} \in (A_\mu)^\wedge$ qui s'envoient sur
$(g \circ x'_{\mu, /Z_\mu})^\sharp(b_1)$ dans $(C_\mu)^\wedge$.
Alors $a_{\mu, j} \in A_\mu^\wedge$ et $(\hat x)^\sharp(b_j) \in A^\wedge$
ont la même image dans $C^\wedge$ par commutativité du
carré ci-dessus. Comme
$\Ker(A \to C) = \colim \Ker(A_\mu \to C_\mu)$ et que ces noyaux
sont égaux à leurs complétés, nous pouvons, quitte à augmenter
$\mu$, modifier nos choix de $a_{\mu, j}$ de sorte que
l'image de $a_{\mu, j}$ dans $A^\wedge$ soit égale à $(\hat x)^\sharp(b_j)$.

\medskip\noindent
Preuve de la surjectivité : derniers arguments algébriques.
Soit $\mathfrak b \subset B$ l'idéal des éléments topologiquement
nilpotents. Soit $J \subset R[x_1, \ldots, x_n]$ l'idéal
formé des $h(x_1, \ldots, x_n)$ tels que
$h(b_1, \ldots, b_n) \in \mathfrak b$. Nous obtenons alors une surjection
continue de $R$-algèbres topologiques
$$
\Phi : R[x_1, \ldots, x_n]^\wedge \longrightarrow B,\quad
x_j \longmapsto b_j
$$
où le membre de gauche est complété relativement à $J$.
Puisque $R[x_1, \ldots, x_n]$ est noethérien, nous pouvons choisir
des générateurs $h_1, \ldots, h_m$ de $J$. Par commutativité
du carré ci-dessus, nous voyons que $h_j(a_{\mu, 1}, \ldots, a_{\mu, n})$ est
un élément de $A_\mu^\wedge$ dont l'image dans $A^\wedge$ est contenue
dans $IA^\wedge$. En effet, l'homomorphisme d'anneaux $(\hat x)^\sharp$ est continu
et $IA^\wedge$ est l'idéal des éléments topologiquement nilpotents
de $A^\wedge$, car $A^\wedge/IA^\wedge = A/I$ est réduit.
(Voir Algèbre, section \ref{algebra-section-completion-noetherian},
pour les résultats sur les complétés des anneaux noethériens.)
Puisque $A/I = \colim A_\mu/I_\mu$, nous concluons que, quitte à augmenter
$\mu$, nous pouvons supposer que $h_j(a_{\mu, 1}, \ldots, a_{\mu, n})$ appartient à
$I_\mu A_\mu^\wedge$. En particulier, les éléments
$h_j(a_{\mu, 1}, \ldots, a_{\mu, n})$ de $A_\mu^\wedge$
sont topologiquement nilpotents dans $A_\mu^\wedge$.
Nous obtenons ainsi un homomorphisme continu de $R$-algèbres
$$
\Psi : R[x_1, \ldots, x_n]^\wedge \longrightarrow A_\mu^\wedge,\quad
x_j \longmapsto a_{\mu, j}
$$
Pour obtenir la conclusion voulue, nous devons voir si $\Ker(\Phi)$ est
annulé par $\Psi$. Cela peut être faux, mais nous pouvons faire en sorte
que ce soit vrai quitte à augmenter $\mu$. En effet, puisque
$R[x_1, \ldots, x_n]^\wedge$ est noethérien,
nous pouvons choisir des générateurs $g_1, \ldots, g_l$ de l'idéal
$\Ker(\Phi)$. Nous voyons alors que
$$
\Psi(g_1), \ldots, \Psi(g_l) \in
\Ker(A_\mu^\wedge \to C_\mu^\wedge) = \Ker(A_\mu \to C_\mu)
$$
ont une image nulle dans $\Ker(A \to C) = \colim \Ker(A_\mu \to C_\mu)$.
Quitte à augmenter $\mu$ comme ci-dessus, nous obtenons donc le résultat voulu.

\medskip\noindent
Preuve de la surjectivité : conclusion. L'homomorphisme continu d'anneaux
$B \to (A_\mu)^\wedge$ construit ci-dessus détermine un morphisme
$\hat x_\mu : V_{\mu, /Z_\mu} \to W$.
La compatibilité de $\hat x_\mu$ et de $u'_\mu$ résulte
du fait que l'homomorphisme d'anneaux $B \to (A_\mu)^\wedge$
est, par construction, compatible avec l'homomorphisme $A_\mu \to C_\mu$.
En fait, la compatibilité est attestée par le morphisme propre
$V'_\mu \to V_\mu$ et les morphismes
$x'_\mu$ et $\hat x'_\mu = x'_{\mu, /Z_\mu}$ que nous avons utilisés dans
la construction. Cela achève la démonstration.
\end{proof}

\begin{lemma}
\label{lemma-rs}
In Situation \ref{situation-contractions} the functor $F$ satisfies
the Rim-Schlessinger condition (RS).
\end{lemma}

\begin{proof}
Recall that the condition only involves the evaluation $F(V)$ of the functor
$F$ on schemes $V$ over $S$ which are spectra of Artinian local rings
and the restriction maps $F(V_2) \to F(V_1)$ for morphisms $V_1 \to V_2$
of schemes over $S$ which are spectra of Artinian local rings.
Thus let $V/S$ be the spetruim of an Artinian local ring.
If $\xi = (Z, u', \hat x) \in F(V)$ then either $Z = \emptyset$
or $Z = V$ (set theoretically). In the first case we see that
$\hat x$ is a morphism from the empty formal algebraic space
into $W$. In the second case we see that $u'$ is a morphism from
the empty scheme into $X'$ and we see that $\hat x : V \to W$
is a morphism into $W$. We conclude that
$$
F(V) = U'(V) \amalg W(V)
$$
and moreover for $V_1 \to V_2$ as above the induced map
$F(V_2) \to F(V_1)$ is compatible with this decomposition.
Hence it suffices to prove that both $U'$ and $W$ satisfy the
Rim-Schlessinger condition. For $U'$ this follows from
Lemma \ref{lemma-algebraic-stack-RS}.
To see that it is true for $W$, we write $W = \colim W_n$ as in
Formal Spaces, Lemma
\ref{formal-spaces-lemma-structure-locally-noetherian}.
Say $V = \Spec(A)$ with $(A, \mathfrak m)$ an Artinian local ring.
Pick $n \geq 1$ such that $\mathfrak m^n = 0$. Then we have
$W(V) = W_n(V)$. Hence we see that the Rim-Schlessinger condition
for $W$ follows from the Rim-Schlessinger condition for $W_n$ for
all $n$ (which in turn follows from
Lemma \ref{lemma-algebraic-stack-RS}).
\end{proof}

\begin{lemma}
\label{lemma-finite-dim}
In Situation \ref{situation-contractions} the tangent spaces of
the functor $F$ are finite dimensional.
\end{lemma}

\begin{proof}
In the proof of Lemma \ref{lemma-rs} we have seen that
$F(V) = U'(V) \amalg W(V)$ if $V$ is the spectrum of an Artinian
local ring. The tangent spaces are computed entirely from evaluations
of $F$ on such schemes over $S$.
Hence it suffices to prove that the tangent spaces
of the functors $U'$ and $W$ are finite dimensional.
For $U'$ this follows from
Lemma \ref{lemma-finite-dimension}.
Write $W = \colim W_n$ as in the proof of Lemma \ref{lemma-rs}.
Then we see that the tangent spaces of $W$ are equal to the
tangent spaces of $W_2$, as to get at the tangent space
we only need to evaluate $W$ on spectra of Artinian local
rings $(A, \mathfrak m)$ with $\mathfrak m^2 = 0$.
Then again we see that the tangent spaces of $W_2$ have
finite dimension by
Lemma \ref{lemma-finite-dimension}.
\end{proof}

\begin{lemma}
\label{lemma-formal-object-effective}
In Situation \ref{situation-contractions} assume $X' \to S$ is separated.
Then every formal object for $F$ is effective.
\end{lemma}

\begin{proof}
A formal object $\xi = (R, \xi_n)$ of $F$ consists of a Noetherian
complete local $S$-algebra $R$ whose residue field is of finite type
over $S$, together with elements $\xi_n \in F(\Spec(R/\mathfrak m^n))$
for all $n$ such that $\xi_{n + 1}|_{\Spec(R/\mathfrak m^n)} = \xi_n$.
By the discussion in the proof of Lemma \ref{lemma-rs}
we see that either $\xi$ is a formal object of $U'$ or a formal
object of $W$. In the first case we see that $\xi$ is effective
by Lemma \ref{lemma-effective}. The second case is the interesting case.
Set $V = \Spec(R)$. We will construct an element
$(Z, u', \hat x) \in F(V)$ whose image in $F(\Spec(R/\mathfrak m^n))$
is $\xi_n$ for all $n \geq 1$.

\medskip\noindent
We may view the collection of elements $\xi_n$ as a morphism
$$
\xi : \text{Spf}(R) \longrightarrow W
$$
of locally Noetherian formal algebraic spaces over $S$. Observe that $\xi$
is {\it not} an adic morphism in general. To fix this, let $I \subset R$
be the ideal corresponding to the formal closed subspace
$$
\text{Spf}(R) \times_{\xi, W} W_{red} \subset \text{Spf}(R)
$$
Note that $I \subset \mathfrak m_R$. Set $Z = V(I) \subset V = \Spec(R)$.
Since $R$ is $\mathfrak m_R$-adically complete it is a fortiori
$I$-adically complete (Algebra, Lemma \ref{algebra-lemma-complete-by-sub}).
Moreover, we claim that for each $n \geq 1$ the morphism
$$
\xi|_{\text{Spf}(R/I^n)} :
\text{Spf}(R/I^n)
\longrightarrow
W
$$
actually comes from a morphism
$$
\xi'_n :  \Spec(R/I^n) \longrightarrow W
$$
Namely, this follows from writing $W = \colim W_n$ as in the
proof of Lemma \ref{lemma-rs}, noticing that $\xi|_{\text{Spf}(R/I^n)}$
maps into $W_n$, and applying Formal Spaces, Lemma
\ref{formal-spaces-lemma-map-into-algebraic-space}
to algebraize this to a morphism $\Spec(R/I^n) \to W_n$
as desired. Let us denote $\text{Spf}'(R) = V_{/Z}$ the formal spectrum
of $R$ endowed with the $I$-adic topology -- equivalently the formal
completion of $V$ along $Z$. Using the morphisms
$\xi'_n$ we obtain an adic morphism
$$
\hat x = (\xi'_n) : \text{Spf}'(R) \longrightarrow W
$$
of locally Noetherian formal algebraic spaces over $S$.
Consider the base change
$$
\text{Spf}'(R) \times_{\hat x, W, g} X'_{/T'} \longrightarrow \text{Spf}'(R)
$$
This is a formal modification by
Algebraization of Formal Spaces, Lemma
\ref{restricted-lemma-base-change-formal-modification}.
Hence by the main theorem on dilatations
(Algebraization of Formal Spaces, Theorem \ref{restricted-theorem-dilatations})
we obtain a proper morphism
$$
V' \longrightarrow V = \Spec(R)
$$
which is an isomorphism over $\Spec(R) \setminus V(I)$ and
whose completion recovers the formal modification above, in other words
$$
V' \times_{\Spec(R)} \Spec(R/I^n) =
\Spec(R/I^n) \times_{\xi'_n, W, g} X'_{/T'}
$$
This in particular tells us we have a compatible system of morphisms
$$
V' \times_{\Spec(R)} \Spec(R/I^n) \longrightarrow X' \times_S \Spec(R/I^n)
$$
Hence by Grothendieck's algebraization theorem (in the form of
More on Morphisms of Spaces, Lemma
\ref{spaces-more-morphisms-lemma-algebraize-morphism})
we obtain a morphism
$$
x' : V' \to X'
$$
over $S$ recovering the morphisms displayed above. Then finally
setting $u' : V \setminus Z \to X'$ the restriction of $x'$ to
$V \setminus Z \subset V'$ gives the third component of our
desired element $(Z, u', \hat x) \in F(V)$.
\end{proof}

\begin{lemma}
\label{lemma-openness-smoothness}
Let $S$ be a locally Noetherian scheme. Let $V$ be a scheme locally
of finite type over $S$. Let $Z \subset V$ be closed. Let $W$ be
a locally Noetherian formal algebraic space over $S$ such that
$W_{red}$ is locally of finite type over $S$. Let $g : V_{/Z} \to W$
be an adic morphism of formal algebraic spaces over $S$. Let $v \in V$
be a closed point such that $g$ is versal at $v$ (as in
Section \ref{section-axioms-functors}).
Then after replacing $V$ by an open neighbourhood of $v$ the
morphism $g$ is smooth (see proof).
\end{lemma}

\begin{proof}
Since $g$ is adic it is representable by algebraic spaces (Formal Spaces,
Section \ref{formal-spaces-section-adic}).
Thus by saying $g$ is smooth we mean that $g$ should be smooth
in the sense of
Bootstrap, Definition \ref{bootstrap-definition-property-transformation}.

\medskip\noindent
Write $W = \colim W_n$ as in Formal Spaces, Lemma
\ref{formal-spaces-lemma-structure-locally-noetherian}.
Set $V_n = V_{/Z} \times_{\hat x, W} W_n$.
Then $V_n$ is a closed subscheme with underlying set $Z$.
Smoothness of $V \to W$ is equivalent to the smoothness
of all the morphisms $V_n \to W_n$ (this holds because any morphism
$T \to W$ with $T$ a quasi-compact scheme factors through
$W_n$ for some $n$). We know that the morphism $V_n \to W_n$
is smooth at $v$ by
Lemma \ref{lemma-base-change-versal}\footnote{The lemma applies since
the diagonal of $W$ is representable by algebraic spaces and
locally of finite type, see Formal Spaces, Lemma
\ref{formal-spaces-lemma-diagonal-morphism-formal-algebraic-spaces}
and we have seen that $W$ has (RS) in the proof of Lemma \ref{lemma-rs}.}.
Of course this means that given any $n$ we can shrink $V$
such that $V_n \to W_n$ is smooth. The problem is to find
an open which works for all $n$ at the same time.

\medskip\noindent
The question is local on $V$, hence we may assume $S = \Spec(R)$ and
$V = \Spec(A)$ are affine.

\medskip\noindent
In this paragraph we reduce to the case where $W$ is an affine formal
algebraic space. Choose an affine formal scheme $W'$ and an \'etale morphism
$W' \to W$ such that the image of $v$ in $W_{red}$ is in the
image of $W'_{red} \to W_{red}$. Then $V_{/Z} \times_{g, W} W' \to V_{/Z}$
is an adic \'etale morphism of formal algebraic spaces over $S$
and $V_{/Z} \times_{g, W} W'$ is an affine formal algebraic space.
By Algebraization of Formal Spaces,
Lemma \ref{restricted-lemma-algebraize-rig-etale-affine}
there exists an \'etale morphism $\varphi : V' \to V$ of affine schemes
such that the completion of $V'$ along $Z' = \varphi^{-1}(Z)$
is isomorphic to $V_{/Z} \times_{g, W} W'$ over $V_{/Z}$.
Observe that $v$ is the image of some $v' \in V'$.
Since smoothness is preserved under base change we see that
$V'_n \to W'_n$ is smooth for all $n$. In the next paragraph
we show that after replacing $V'$ by an open neighbourhood of $v'$
the morphisms $V'_n \to W'_n$ are smooth for all $n$.
Then, after we replace $V$ by the open image of $V' \to V$,
we obtain that $V_n \to W_n$ is smooth by \'etale descent of smoothness.
Some details omitted.

\medskip\noindent
Assume $S = \Spec(R)$, $V = \Spec(A)$, $Z = V(I)$, and $W = \text{Spf}(B)$.
Let $v$ correspond to the maximal ideal $I \subset \mathfrak m \subset A$.
We are given an adic continuous $R$-algebra homomorphism
$$
B \longrightarrow A^\wedge
$$
Let $\mathfrak b \subset B$ be the ideal of topologically nilpotent
elements (this is the maximal ideal of definition of the Noetherian adic
topological ring $B$). Observe that $\mathfrak b A^\wedge$ and
$IA^\wedge$ are both ideals of definition of the Noetherian adic
ring $A^\wedge$. Also, $\mathfrak m A^\wedge$ is a maximal ideal
of $A^\wedge$ containing both $\mathfrak b A^\wedge$ and $IA^\wedge$.
We are given that
$$
B_n = B/\mathfrak b^n \to A^\wedge/\mathfrak b^n A^\wedge = A_n
$$
is smooth at $\mathfrak m$ for all $n$. By the discussion above
we may and do assume that $B_1 \to A_1$ is a smooth ring map.
Denote $\mathfrak m_1 \subset A_1$ the maximal ideal corresponding
to $\mathfrak m$. Since smoothness implies flatness, we see that:
for all $n \geq 1$ the map
$$
\mathfrak b^n/\mathfrak b^{n + 1} \otimes_{B_1} (A_1)_{\mathfrak m_1}
\longrightarrow
\left(\mathfrak b^nA^\wedge/\mathfrak b^{n + 1}A^\wedge\right)_{\mathfrak m_1}
$$
is an isomorphism (see
Algebra, Lemma \ref{algebra-lemma-what-does-it-mean-again}).
Consider the Rees algebra
$$
B' = \bigoplus\nolimits_{n \geq 0} \mathfrak b^n/\mathfrak b^{n + 1}
$$
which is a finite type graded algebra over the Noetherian ring $B_1$ and
the Rees algebra
$$
A' = \bigoplus\nolimits_{n \geq 0}
\mathfrak b^nA^\wedge/\mathfrak b^{n + 1}A^\wedge
$$
which is a a finite type graded algebra over the Noetherian ring $A_1$.
Consider the homomorphism of graded $A_1$-algebras
$$
\Psi : B' \otimes_{B_1} A_1 \longrightarrow A'
$$
By the above this map is an isomorphism after localizing at
the maximal ideal $\mathfrak m_1$ of $A_1$.
Hence $\Ker(\Psi)$, resp.\ $\Coker(\Psi)$ is a finite module
over $B' \otimes_{B_1} A_1$, resp.\ $A'$ whose localization
at $\mathfrak m_1$ is zero. It follows that after replacing
$A_1$ (and correspondingly $A$) by a principal localization
we may assume $\Psi$ is an isomorphism. (This is the key step of the proof.)
Then working backwards we see that $B_n \to A_n$ is flat, see
Algebra, Lemma \ref{algebra-lemma-what-does-it-mean-again}.
Hence $A_n \to B_n$ is smooth (as a flat ring map with smooth
fibres, see Algebra, Lemma \ref{algebra-lemma-flat-fibre-smooth})
and the proof is complete.
\end{proof}

\begin{lemma}
\label{lemma-openness-versality}
In Situation \ref{situation-contractions} the functor
$F$ satisfies openness of versality.
\end{lemma}

\begin{proof}
We have to show the following. Given a scheme $V$ locally of finite type over
$S$, given $\xi \in F(V)$, and given a finite type point $v_0 \in V$ such that
$\xi$ is versal at $v_0$, after replacing $V$ by an open neighbourhood
of $v_0$ we have that $\xi$ is versal at every finite type point of $V$.
Write $\xi = (Z, u', \hat x)$.

\medskip\noindent
First case: $v_0 \not \in Z$. Then we can first replace $V$ by
$V \setminus Z$. Hence we see that $\xi = (\emptyset, u', \emptyset)$
and the morphism $u' : V \to X'$ is versal at $v_0$.
By More on Morphisms of Spaces, Lemma
\ref{spaces-more-morphisms-lemma-lifting-along-artinian-at-point}
this means that $u' : V \to X'$ is smooth at $v_0$.
Since the set of a points where a morphism is smooth is open,
we can after shrinking $V$ assume $u'$ is smooth.
Then the same lemma tells us that $\xi$ is versal at every
point as desired.

\medskip\noindent
Second case: $v_0 \in Z$. Write $W = \colim W_n$ as in
Formal Spaces, Lemma \ref{formal-spaces-lemma-structure-locally-noetherian}.
By Lemma \ref{lemma-openness-smoothness} we may assume $\hat x : V_{/Z} \to W$
is a smooth morphism of formal algebraic spaces. It follows immediately
that $\xi = (Z, u', \hat x)$ is versal at all finite type points of $Z$.
Let $V' \to V$, $\hat x'$, and $x'$ witness the compatibility between $u'$
and $\hat x$. We see that $\hat x' : V'_{/Z} \to X'_{/T'}$ is smooth as a
base change of $\hat x$. Since $\hat x'$ is the completion of
$x' : V' \to X'$ this implies that $x' : V' \to X'$ is smooth at all
points of $(V' \to V)^{-1}(Z) = |x'|^{-1}(T') \subset |V'|$
by the already used More on Morphisms of Spaces, Lemma
\ref{spaces-more-morphisms-lemma-lifting-along-artinian-at-point}.
Since the set of smooth points of a morphism is open, we see that
the closed set of points $B \subset |V'|$ where $x'$ is not smooth
does not meet $(V' \to V)^{-1}(Z)$. Since $V' \to V$ is proper and
hence closed, we see that $(V' \to V)(B) \subset V$ is a closed
subset not meeting $Z$. Hence after shrinking $V$ we may assume
$B = \emptyset$, i.e., $x'$ is smooth. By the discussion in the previous
paragraph this exactly means that $\xi$ is versal at all finite type
points of $V$ not contained in $Z$ and the proof is complete.
\end{proof}

\noindent
Here is the final result.

\begin{theorem}
\label{theorem-contractions}
\begin{reference}
\cite[Theorem 3.1]{ArtinII}
\end{reference}
Let $S$ be a locally Noetherian scheme such that $\mathcal{O}_{S, s}$
is a G-ring for all finite type points $s \in S$. Let $X'$ be an algebraic
space locally of finite type over $S$. Let $T' \subset |X'|$ be a closed
subset. Let $W$ be a locally Noetherian formal algebraic space over $S$
with $W_{red}$ locally of finite type over $S$. Finally, we let
$$
g : X'_{/T'} \longrightarrow W
$$
be a formal modification, see Algebraization of Formal Spaces, Definition
\ref{restricted-definition-formal-modification}. If $X'$ and $W$ are
separated\footnote{See Remark \ref{remark-separated-needed}.} over $S$, then
there exists a proper morphism $f : X' \to X$ of algebraic spaces over $S$,
a closed subset $T \subset |X|$, and an isomorphism $a : X_{/T} \to W$
of formal algebraic spaces such that
\begin{enumerate}
\item $T'$ is the inverse image of $T$ by $|f| : |X'| \to |X|$,
\item $f : X' \to X$ maps $X' \setminus T'$ isomorphically to
$X \setminus T$, and
\item $g = a \circ f_{/T}$ where $f_{/T} : X'_{/T'} \to X_{/T}$
is the induced morphism.
\end{enumerate}
In other words, $(f : X' \to X, T, a)$ is a solution as defined earlier in
this section.
\end{theorem}

\begin{proof}
Let $F$ be the functor constructed using $X'$, $T'$, $W$, $g$ in this section.
By Lemma \ref{lemma-functor-is-solution} it suffices to show that
$F$ corresponds to an algebraic space $X$ locally of finite type over $S$.
In order to do this, we will apply
Proposition \ref{proposition-spaces-diagonal-representable-noetherian}.
Namely, by Lemma \ref{lemma-diagonal-contractions}
the diagonal of $F$ is representable by closed immersions
and by
Lemmas \ref{lemma-sheaf}, \ref{lemma-limit-preserving},
\ref{lemma-rs}, \ref{lemma-finite-dim},
\ref{lemma-formal-object-effective}, and \ref{lemma-openness-versality}
we have axioms [0], [1], [2], [3], [4], and [5].
\end{proof}

\begin{remark}
\label{remark-separated-needed}
The proof of Theorem \ref{theorem-contractions} uses that $X'$ and $W$
are separated over $S$ in two places. First, the proof uses this in showing
$\Delta : F \to F \times F$ is representable by algebraic spaces.
This use of the assumption can be entirely avoided by proving
that $\Delta$ is representable by applying the theorem in the
separated case to the triples
$E'$, $(E' \to V)^{-1}Z$, and $E'_{/Z} \to E_W$
found in Remark \ref{remark-diagonal} (this is the usual bootstrap
procedure for the diagonal). Thus the proof of
Lemma \ref{lemma-formal-object-effective} is the only
place in our proof of Theorem \ref{theorem-contractions}
where we really need to use that $X' \to S$ is separated.
The reader checks that we use the assumption only to obtain
the morphism $x' : V' \to X'$. The existence of $x'$ can be shown,
using results in the literature, if $X' \to S$ is quasi-separated, see
More on Morphisms of Spaces, Remark
\ref{spaces-more-morphisms-remark-weaken-separation-axioms-question}.
We conclude the theorem holds as stated with
``separated'' replaced by ``quasi-separated''. If we ever need this
we will precisely state and carefully prove this here.
\end{remark}















\begin{multicols}{2}[\section{Autres chapitres}]
\noindent
Préliminaires
\begin{enumerate}
\item \hyperref[introduction-section-phantom]{Introduction}
\item \hyperref[conventions-section-phantom]{Conventions}
\item \hyperref[sets-section-phantom]{Théorie des ensembles}
\item \hyperref[categories-section-phantom]{Catégories}
\item \hyperref[topology-section-phantom]{Topologie}
\item \hyperref[sheaves-section-phantom]{Faisceaux sur les espaces}
\item \hyperref[sites-section-phantom]{Sites et faisceaux}
\item \hyperref[stacks-section-phantom]{Champs}
\item \hyperref[fields-section-phantom]{Corps}
\item \hyperref[algebra-section-phantom]{Algèbre commutative}
\item \hyperref[brauer-section-phantom]{Groupes de Brauer}
\item \hyperref[homology-section-phantom]{Algèbre homologique}
\item \hyperref[derived-section-phantom]{Catégories dérivées}
\item \hyperref[simplicial-section-phantom]{Méthodes simpliciales}
\item \hyperref[more-algebra-section-phantom]{Compléments d'algèbre}
\item \hyperref[smoothing-section-phantom]{Lissage des morphismes d'anneaux}
\item \hyperref[modules-section-phantom]{Faisceaux de modules}
\item \hyperref[sites-modules-section-phantom]{Modules sur les sites}
\item \hyperref[injectives-section-phantom]{Objets injectifs}
\item \hyperref[cohomology-section-phantom]{Cohomologie des faisceaux}
\item \hyperref[sites-cohomology-section-phantom]{Cohomologie sur les sites}
\item \hyperref[dga-section-phantom]{Algèbre différentielle graduée}
\item \hyperref[dpa-section-phantom]{Algèbre à puissances divisées}
\item \hyperref[sdga-section-phantom]{Faisceaux différentiels gradués}
\item \hyperref[hypercovering-section-phantom]{Hyperrecouvrements}
\end{enumerate}
Schémas
\begin{enumerate}
\setcounter{enumi}{25}
\item \hyperref[schemes-section-phantom]{Schémas}
\item \hyperref[constructions-section-phantom]{Constructions de schémas}
\item \hyperref[properties-section-phantom]{Propriétés des schémas}
\item \hyperref[morphisms-section-phantom]{Morphismes de schémas}
\item \hyperref[coherent-section-phantom]{Cohomologie des schémas}
\item \hyperref[divisors-section-phantom]{Diviseurs}
\item \hyperref[limits-section-phantom]{Limites de schémas}
\item \hyperref[varieties-section-phantom]{Variétés}
\item \hyperref[topologies-section-phantom]{Topologies sur les schémas}
\item \hyperref[descent-section-phantom]{Descente}
\item \hyperref[perfect-section-phantom]{Catégories dérivées de schémas}
\item \hyperref[more-morphisms-section-phantom]{Compléments sur les morphismes}
\item \hyperref[flat-section-phantom]{Compléments sur la platitude}
\item \hyperref[groupoids-section-phantom]{Schémas en groupoïdes}
\item \hyperref[more-groupoids-section-phantom]{Compléments sur les schémas en groupoïdes}
\item \hyperref[etale-section-phantom]{Morphismes étales de schémas}
\end{enumerate}
Thèmes de théorie des schémas
\begin{enumerate}
\setcounter{enumi}{41}
\item \hyperref[chow-section-phantom]{Homologie de Chow}
\item \hyperref[intersection-section-phantom]{Théorie de l'intersection}
\item \hyperref[pic-section-phantom]{Schémas de Picard des courbes}
\item \hyperref[weil-section-phantom]{Théories cohomologiques de Weil}
\item \hyperref[adequate-section-phantom]{Modules adéquats}
\item \hyperref[dualizing-section-phantom]{Complexes dualisants}
\item \hyperref[duality-section-phantom]{Dualité pour les schémas}
\item \hyperref[discriminant-section-phantom]{Discriminants et différentes}
\item \hyperref[derham-section-phantom]{Cohomologie de de Rham}
\item \hyperref[local-cohomology-section-phantom]{Cohomologie locale}
\item \hyperref[algebraization-section-phantom]{Géométrie algébrique et formelle}
\item \hyperref[curves-section-phantom]{Courbes algébriques}
\item \hyperref[resolve-section-phantom]{Résolution des surfaces}
\item \hyperref[models-section-phantom]{Réduction semi-stable}
\item \hyperref[functors-section-phantom]{Foncteurs et morphismes}
\item \hyperref[equiv-section-phantom]{Catégories dérivées de variétés}
\item \hyperref[pione-section-phantom]{Groupes fondamentaux de schémas}
\item \hyperref[etale-cohomology-section-phantom]{Cohomologie étale}
\item \hyperref[crystalline-section-phantom]{Cohomologie cristalline}
\item \hyperref[proetale-section-phantom]{Cohomologie pro-étale}
\item \hyperref[relative-cycles-section-phantom]{Cycles relatifs}
\item \hyperref[more-etale-section-phantom]{Compléments de cohomologie étale}
\item \hyperref[trace-section-phantom]{La formule des traces}
\end{enumerate}
Espaces algébriques
\begin{enumerate}
\setcounter{enumi}{64}
\item \hyperref[spaces-section-phantom]{Espaces algébriques}
\item \hyperref[spaces-properties-section-phantom]{Propriétés des espaces algébriques}
\item \hyperref[spaces-morphisms-section-phantom]{Morphismes d'espaces algébriques}
\item \hyperref[decent-spaces-section-phantom]{Espaces algébriques décents}
\item \hyperref[spaces-cohomology-section-phantom]{Cohomologie des espaces algébriques}
\item \hyperref[spaces-limits-section-phantom]{Limites d'espaces algébriques}
\item \hyperref[spaces-divisors-section-phantom]{Diviseurs sur les espaces algébriques}
\item \hyperref[spaces-over-fields-section-phantom]{Espaces algébriques sur des corps}
\item \hyperref[spaces-topologies-section-phantom]{Topologies sur les espaces algébriques}
\item \hyperref[spaces-descent-section-phantom]{Descente et espaces algébriques}
\item \hyperref[spaces-perfect-section-phantom]{Catégories dérivées d'espaces}
\item \hyperref[spaces-more-morphisms-section-phantom]{Compléments sur les morphismes d'espaces}
\item \hyperref[spaces-flat-section-phantom]{Platitude sur les espaces algébriques}
\item \hyperref[spaces-groupoids-section-phantom]{Groupoïdes dans les espaces algébriques}
\item \hyperref[spaces-more-groupoids-section-phantom]{Compléments sur les groupoïdes dans les espaces}
\item \hyperref[bootstrap-section-phantom]{Amorçage}
\item \hyperref[spaces-pushouts-section-phantom]{Sommes amalgamées d'espaces algébriques}
\end{enumerate}
Thèmes de géométrie
\begin{enumerate}
\setcounter{enumi}{81}
\item \hyperref[spaces-chow-section-phantom]{Groupes de Chow des espaces}
\item \hyperref[groupoids-quotients-section-phantom]{Quotients de groupoïdes}
\item \hyperref[spaces-more-cohomology-section-phantom]{Compléments sur la cohomologie des espaces}
\item \hyperref[spaces-simplicial-section-phantom]{Espaces simpliciaux}
\item \hyperref[spaces-duality-section-phantom]{Dualité pour les espaces}
\item \hyperref[formal-spaces-section-phantom]{Espaces algébriques formels}
\item \hyperref[restricted-section-phantom]{Algébrisation des espaces formels}
\item \hyperref[spaces-resolve-section-phantom]{Résolution des surfaces revisitée}
\end{enumerate}
Théorie des déformations
\begin{enumerate}
\setcounter{enumi}{89}
\item \hyperref[formal-defos-section-phantom]{Théorie formelle des déformations}
\item \hyperref[defos-section-phantom]{Théorie des déformations}
\item \hyperref[cotangent-section-phantom]{Le complexe cotangent}
\item \hyperref[examples-defos-section-phantom]{Problèmes de déformation}
\end{enumerate}
Champs algébriques
\begin{enumerate}
\setcounter{enumi}{93}
\item \hyperref[algebraic-section-phantom]{Champs algébriques}
\item \hyperref[examples-stacks-section-phantom]{Exemples de champs}
\item \hyperref[stacks-sheaves-section-phantom]{Faisceaux sur les champs algébriques}
\item \hyperref[criteria-section-phantom]{Critères de représentabilité}
\item \hyperref[artin-section-phantom]{Axiomes d'Artin}
\item \hyperref[quot-section-phantom]{Espaces de Quot et de Hilbert}
\item \hyperref[stacks-properties-section-phantom]{Propriétés des champs algébriques}
\item \hyperref[stacks-morphisms-section-phantom]{Morphismes de champs algébriques}
\item \hyperref[stacks-limits-section-phantom]{Limites de champs algébriques}
\item \hyperref[stacks-cohomology-section-phantom]{Cohomologie des champs algébriques}
\item \hyperref[stacks-perfect-section-phantom]{Catégories dérivées de champs}
\item \hyperref[stacks-introduction-section-phantom]{Introduction aux champs algébriques}
\item \hyperref[stacks-more-morphisms-section-phantom]{Compléments sur les morphismes de champs}
\item \hyperref[stacks-geometry-section-phantom]{La géométrie des champs}
\end{enumerate}
Thèmes de théorie des espaces de modules
\begin{enumerate}
\setcounter{enumi}{107}
\item \hyperref[moduli-section-phantom]{Champs de modules}
\item \hyperref[moduli-curves-section-phantom]{Espaces de modules de courbes}
\end{enumerate}
Divers
\begin{enumerate}
\setcounter{enumi}{109}
\item \hyperref[examples-section-phantom]{Exemples}
\item \hyperref[exercises-section-phantom]{Exercices}
\item \hyperref[guide-section-phantom]{Guide bibliographique}
\item \hyperref[desirables-section-phantom]{Desiderata}
\item \hyperref[coding-section-phantom]{Style de programmation}
\item \hyperref[obsolete-section-phantom]{Obsolète}
\item \hyperref[fdl-section-phantom]{Licence de documentation libre GNU}
\item \hyperref[index-section-phantom]{Index généré automatiquement}
\end{enumerate}
\end{multicols}


\bibliography{my}
\bibliographystyle{amsalpha}

\end{document}
